% Leçon 4 — Expression écrite : écriture des caractères de l'amour et de la famille ; types de phrases étudiées — Chinois 1ère A — Cours signé OnBuch+
% Programme officiel MINESEC. Compiler avec : tectonic ZHA04-ecrit-ecriture-des-caracteres-de-l-amour-et-de-la.tex
\newcommand{\DOCMATIERE}{Chinois}
\newcommand{\DOCNIVEAU}{1ère A}
\newcommand{\DOCMODULE}{Module 1 : Vie familiale et intégration sociale}
\newcommand{\DOCLECON}{ZHA04}
\newcommand{\DOCTITRE}{Leçon 4 — Expression écrite : écriture des caractères de l'amour et de la famille ; types de phrases étudiées}
% OnBuch+ — préambule « cours signés » ENRICHI (compilé avec tectonic / XeLaTeX)
% Système visuel « Pop » d'OnBuch+ : fond crème (jamais blanc), bordures encre
% épaisses, ombres « dures » décalées, titres Archivo Black en capitales.
% Version MATHÉMATIQUES : graphes (pgfplots/tikz), boîtes pédagogiques
% (définition, propriété, méthode, attention, le savais-tu, exercice résolu,
% exercice, TP, bilan), page de garde et signature OnBuch+ ancrée en bas.
%
% Macros attendues dans main.tex AVANT \input{preamble} :
%   \DOCMATIERE  \DOCNIVEAU  \DOCMODULE  \DOCLECON (numéro)  \DOCTITRE
\documentclass[11pt]{article}

\usepackage{fontspec}
\usepackage[french]{babel} % français = langue principale ; le chinois via \zh{..} / textebox (xeCJK)
\usepackage{xeCJK}
\usepackage{pifont} % \ding{51} \ding{55} utilisés par les modèles
\usepackage[a4paper,margin=2cm,top=3.4cm,bottom=2.8cm,headheight=1.4cm,footskip=1.3cm]{geometry}
\usepackage{amsmath,amssymb,amsfonts}
\usepackage{mathtools} % \xmapsto, \coloneqq... souvent utilisés par les modèles
\usepackage{cancel} % simplifications barrées (bilans de forces)
% \abs{z} (module, valeur absolue) et \norm{u} (norme), utilisés par les modèles
\providecommand{\abs}{}\let\abs\relax\DeclarePairedDelimiter{\abs}{\lvert}{\rvert}
\providecommand{\norm}{}\let\norm\relax\DeclarePairedDelimiter{\norm}{\lVert}{\rVert}
\usepackage[table]{xcolor} % [table] : fournit aussi \rowcolors (alternance de lignes)
\usepackage{pagecolor}
\usepackage{tikz}
\usetikzlibrary{arrows.meta,positioning,calc,shapes.geometric,shapes.misc,shapes.arrows,decorations.pathmorphing,decorations.pathreplacing,decorations.markings,patterns,shadows,mindmap,trees,fit,backgrounds,matrix,intersections,angles,quotes}
\usepackage{pgfplots}
\usepackage[version=4]{mhchem} % équations nucléaires \ce{^{235}_{92}U}
\usepackage[european,siunitx]{circuitikz} % schémas électriques
\pgfplotsset{compat=1.18}
\usepgfplotslibrary{fillbetween,groupplots} % aires entre courbes (calcul intégral)
\usepackage{enumitem}
\usepackage{fancyhdr}
% [most] charge toutes les bibliothèques tcolorbox (listings, minted, raster,
% documentation...) et gonfle inutilement la mémoire TeX sur un document long
% (~300 boîtes) : on ne charge que ce qui est réellement utilisé.
\usepackage[skins,breakable]{tcolorbox}
\usepackage{graphicx}
\usepackage{colortbl}
\usepackage{booktabs}
\usepackage{tabularx}
\usepackage{diagbox} % case d'en-tête diagonale des tableaux à double entrée
% Colonnes extensibles centrée / gauche / droite (C, L, R), souvent utilisées par les modèles
\newcolumntype{C}{>{\centering\arraybackslash}X}
\newcolumntype{L}{>{\raggedright\arraybackslash}X}
\newcolumntype{R}{>{\raggedleft\arraybackslash}X}
\usepackage{array}
\usepackage{multicol}
\usepackage{multirow}
\usepackage{siunitx}
% parse-numbers=false : accepte aussi \SI{2,5 \times 10^{-3}}{...}
\sisetup{locale=FR,output-decimal-marker={,},group-minimum-digits=4,parse-numbers=false}
\DeclareSIUnit\ohm{\ensuremath{\symup{\Omega}}} % U+2126 absent de la police
\DeclareSIUnit\division{div} % réglages d'oscilloscope (ms/div, V/div)
\DeclareSIUnit\tr{tr} % tours (vitesse de rotation, ex. \SI{1500}{\tr\per\minute})
\DeclareSIUnit\molar{M} % concentration molaire (ex. \SI{3}{\molar}, \SI{1}{\micro\molar})
\DeclareSIUnit\an{an} % année (le modèle confond parfois avec \year, un compteur TeX) — ex. \SI{5}{\kilo\meter\per\an}
\DeclareSIUnit\FCFA{FCFA} % franc CFA (ex. \SI{500}{\FCFA\per\kilo\gram})
\DeclareSIUnit\degreeBrix{\textdegree Brix} % teneur en sucre d'un sirop/jus (réfractométrie)
\DeclareSIUnit\month{mois} % le modèle utilise parfois \month, un compteur TeX, comme unité de durée
\usepackage{qrcode}
% \degree : commande fréquemment utilisée par les modèles pour un angle ou une
% température en dehors de \SI{}{} (non fournie par les paquets chargés).
\providecommand{\degree}{^{\circ}}
% \qed (fin de démonstration), utilisé par les modèles sans amsthm
\providecommand{\qed}{\hfill$\square$}
% \barre : double barre des valeurs interdites dans un tableau de variations (tkz-tab)
\providecommand{\barre}{\ensuremath{\|}}
% environnement proof (amsthm non chargé) utilisé par les modèles
\ifdefined\proof\else\newenvironment{proof}[1][Démonstration]{\par\noindent\textit{#1.}\ }{\qed\par}\fi
\usepackage{lastpage}
\usepackage{hyperref}
\hypersetup{hidelinks}

% ============================================================
% Palette « Pop » OnBuch+ — mêmes hex que la classe P (plus_ui.dart)
% ============================================================
\definecolor{poporange}{HTML}{F59321}
\definecolor{popdark}{HTML}{E07A0C}
\definecolor{popcream}{HTML}{FFF6E8}
\definecolor{popink}{HTML}{1C1714}
\definecolor{poppurple}{HTML}{7A5AE0}
\definecolor{popgreen}{HTML}{2E9E6A}
\definecolor{poppink}{HTML}{E0457B}
\definecolor{popblue}{HTML}{2F7DE1}
\definecolor{popgold}{HTML}{C98A12}
\definecolor{popmuted}{HTML}{8A7F76}
\definecolor{popline}{HTML}{EADFD2}
\definecolor{popgray}{HTML}{8A7F76}
\colorlet{popgrey}{popgray}
% Alias tolérants (le modèle nomme parfois les couleurs différemment)
\colorlet{popwhite}{white}
\colorlet{popblack}{popink}
\colorlet{popred}{poppink}
\colorlet{popyellow}{popgold}
% Teintes claires (fonds de tableaux/encadrés) — le modèle nomme parfois
% les couleurs avec un suffixe L (ex. popblueL) sans les avoir définies.
\colorlet{poporangeL}{poporange!20}
\colorlet{popdarkL}{popdark!20}
\colorlet{poppurpleL}{poppurple!20}
\colorlet{popgreenL}{popgreen!20}
\colorlet{poppinkL}{poppink!20}
\colorlet{popblueL}{popblue!20}
\colorlet{popgoldL}{popgold!20}
\colorlet{popredL}{popred!20}
\colorlet{popgrayL}{popgray!20}
% Teintes claires pour fonds de tableaux / zones
\colorlet{popblueL}{popblue!10!white}
\colorlet{poporangeL}{poporange!14!white}
\colorlet{popgreenL}{popgreen!12!white}
\colorlet{poppurpleL}{poppurple!12!white}
\colorlet{poppinkL}{poppink!12!white}
\colorlet{popgoldL}{popgold!14!white}

\pagecolor{popcream}

% ============================================================
% Typographie
% ============================================================
\defaultfontfeatures{Ligatures=TeX}
\setmainfont{PlusJakartaSans-Medium.ttf}[Path=../../../fonts/,BoldFont=PlusJakartaSans-Bold.ttf]
% Symboles, API et emojis absents de la police principale : repli sur DejaVu Sans (caractères actifs)
\newfontface\symfont{DejaVuSans.ttf}[Path=../../../fonts/]
\makeatletter
\newcommand\symactive[1]{\catcode#1=13 \begingroup\lccode`\~=#1 \lowercase{\endgroup\def~}{{\symfont\char#1}}}
\newcommand\symblank[1]{\catcode#1=13 \begingroup\lccode`\~=#1 \lowercase{\endgroup\def~}{}}
\newcommand\symhyph[1]{\catcode#1=13 \begingroup\lccode`\~=#1 \lowercase{\endgroup\def~}{\mbox{-}}}
\newcount\symc
\newcommand\symrange[2]{\symc=#1 \@whilenum\symc<#2 \do{\expandafter\symactive\expandafter{\the\symc}\advance\symc by 1 }}
\newcommand\symblankrange[2]{\symc=#1 \@whilenum\symc<#2 \do{\expandafter\symblank\expandafter{\the\symc}\advance\symc by 1 }}
\makeatother
\symrange{"0250}{"02FF}\symactive{"2713}\symactive{"2714}\symactive{"2717}\symactive{"2718}\symactive{"2610}\symactive{"2611}\symactive{"2612}
\symrange{"2600}{"27BF}
\catcode"2705=13 \begingroup\lccode`\~="2705 \lowercase{\endgroup\def~}{{\symfont\char"2713}}\symblank{"26BD}\symblank{"26A1}
\symrange{"2460}{"257F}\symactive{"223C}\symactive{"2248}\symactive{"2260}
\symhyph{"2011}\symhyph{"2E1A}\symblankrange{"1F300}{"1FAFF}\symblank{"FE0F}
% Caractères chinois : WenQuanYi Zen Hei (sous-ensemble GB2312 + pinyin), gras/italique simulés
\setCJKmainfont{WenQuanYiZenHei-subset.ttf}[Path=../../../fonts/,AutoFakeBold=3,AutoFakeSlant=0.2]
\xeCJKsetup{CJKspace=false}
% Pinyin : voyelles à caron (ǎ ǐ ǒ ǔ ǖ ǘ ǚ ǜ) absentes de la police principale -> caractères actifs
\newfontface\pyfont{WenQuanYiZenHei-subset.ttf}[Path=../../../fonts/]
\newcommand\pyactive[1]{\catcode#1=13 \begingroup\lccode`\~=#1 \lowercase{\endgroup\def~}{{\pyfont\char#1}}}
\pyactive{"01CD}\pyactive{"01CE}\pyactive{"01CF}\pyactive{"01D0}\pyactive{"01D1}\pyactive{"01D2}\pyactive{"01D3}\pyactive{"01D4}\pyactive{"01D5}\pyactive{"01D6}\pyactive{"01D7}\pyactive{"01D8}\pyactive{"01D9}\pyactive{"01DA}\pyactive{"01DB}\pyactive{"01DC}
% Maths en sans-serif (Fira Math) avec chiffres et lettres droites de Plus
% Jakarta, pour que les formules se fondent dans le texte.
\usepackage[math-style=ISO,bold-style=ISO]{unicode-math}
\setmathfont{FiraMath-Regular.otf}[Path=../../../fonts/]
\setmathfont{PlusJakartaSans-Medium.ttf}[Path=../../../fonts/,range={up/num,up/{Latin,latin},\mathup}]
\usepackage{icomma} % virgule décimale sans espace en mode math
% Caractères mathématiques Unicode absents des polices de texte (Plus Jakarta,
% Archivo Black) : rendus par leur équivalent mathématique, en texte comme en
% formule — sinon ils s'affichent en carré vide (ex. « Arithmétique dans ℤ »).
\usepackage{newunicodechar}
\newunicodechar{ℝ}{\ensuremath{\mathbb{R}}}
\newunicodechar{ℤ}{\ensuremath{\mathbb{Z}}}
\newunicodechar{ℂ}{\ensuremath{\mathbb{C}}}
\newunicodechar{ℕ}{\ensuremath{\mathbb{N}}}
\newunicodechar{ℚ}{\ensuremath{\mathbb{Q}}}
\newunicodechar{𝔻}{\ensuremath{\mathbb{D}}}
\newunicodechar{α}{\ensuremath{\alpha}}
\newunicodechar{β}{\ensuremath{\beta}}
\newunicodechar{γ}{\ensuremath{\gamma}}
\newunicodechar{δ}{\ensuremath{\delta}}
\newunicodechar{ε}{\ensuremath{\varepsilon}}
\newunicodechar{θ}{\ensuremath{\theta}}
\newunicodechar{λ}{\ensuremath{\lambda}}
\newunicodechar{μ}{\ensuremath{\mu}}
\newunicodechar{σ}{\ensuremath{\sigma}}
\newunicodechar{τ}{\ensuremath{\tau}}
\newunicodechar{φ}{\ensuremath{\varphi}}
\newunicodechar{ψ}{\ensuremath{\psi}}
\newunicodechar{ω}{\ensuremath{\omega}}
\newunicodechar{Σ}{\ensuremath{\Sigma}}
% majuscules produites par \MakeUppercase dans les titres (α-aminés -> Α)
\newunicodechar{Α}{\ensuremath{\alpha}}
\newunicodechar{Β}{\ensuremath{\beta}}
\newunicodechar{Γ}{\ensuremath{\Gamma}}
\newunicodechar{Δ}{\ensuremath{\Delta}}
\newunicodechar{Ε}{\ensuremath{\varepsilon}}
\newunicodechar{Θ}{\ensuremath{\Theta}}
\newunicodechar{Λ}{\ensuremath{\Lambda}}
\newunicodechar{Μ}{\ensuremath{\mu}}
\newunicodechar{Τ}{\ensuremath{\tau}}
\newunicodechar{Φ}{\ensuremath{\Phi}}
\newunicodechar{Ψ}{\ensuremath{\Psi}}
\newunicodechar{Ω}{\ensuremath{\Omega}}
\newunicodechar{↦}{\ensuremath{\mapsto}}
\newunicodechar{∈}{\ensuremath{\in}}
\newunicodechar{∉}{\ensuremath{\notin}}
\newunicodechar{∧}{\ensuremath{\wedge}}
\newunicodechar{⇔}{\ensuremath{\Leftrightarrow}}
\newunicodechar{⟺}{\ensuremath{\Longleftrightarrow}}
\newunicodechar{⇒}{\ensuremath{\Rightarrow}}
\newunicodechar{⟹}{\ensuremath{\Longrightarrow}}
\newunicodechar{∘}{\ensuremath{\circ}}
\newunicodechar{∪}{\ensuremath{\cup}}
\newunicodechar{∩}{\ensuremath{\cap}}
\newunicodechar{≡}{\ensuremath{\equiv}}
\newunicodechar{⊂}{\ensuremath{\subset}}
\newunicodechar{⊥}{\ensuremath{\perp}}
\newunicodechar{∃}{\ensuremath{\exists}}
\newunicodechar{∀}{\ensuremath{\forall}}
\newunicodechar{∛}{\ensuremath{\sqrt[3]{\,}}}
\newunicodechar{∜}{\ensuremath{\sqrt[4]{\,}}}
\newunicodechar{⃗}{\ensuremath{{}^{\to}}}
\newunicodechar{ⁿ}{\ifmmode^{n}\else\textsuperscript{\ensuremath{n}}\fi}
\newunicodechar{ˣ}{\ifmmode^{x}\else\textsuperscript{\ensuremath{x}}\fi}
\newunicodechar{ᵇ}{\ifmmode^{b}\else\textsuperscript{\ensuremath{b}}\fi}
\newunicodechar{ʳ}{\ifmmode^{r}\else\textsuperscript{\ensuremath{r}}\fi}
\newunicodechar{ᵘ}{\ifmmode^{u}\else\textsuperscript{\ensuremath{u}}\fi}
\newunicodechar{ʸ}{\ifmmode^{y}\else\textsuperscript{\ensuremath{y}}\fi}
\newunicodechar{ᵃ}{\ifmmode^{a}\else\textsuperscript{\ensuremath{a}}\fi}
\newunicodechar{ᵉ}{\ifmmode^{e}\else\textsuperscript{\ensuremath{e}}\fi}
\newunicodechar{ᵏ}{\ifmmode^{k}\else\textsuperscript{\ensuremath{k}}\fi}
\newunicodechar{ᵖ}{\ifmmode^{p}\else\textsuperscript{\ensuremath{p}}\fi}
\newunicodechar{ᴺ}{\ifmmode^{N}\else\textsuperscript{\ensuremath{N}}\fi}
\newunicodechar{ᶜ}{\ifmmode^{c}\else\textsuperscript{\ensuremath{c}}\fi}
\newunicodechar{ʰ}{\ifmmode^{h}\else\textsuperscript{\ensuremath{h}}\fi}
\newunicodechar{ᵠ}{\ifmmode^{q}\else\textsuperscript{\ensuremath{q}}\fi}
\newunicodechar{ₙ}{\ifmmode_{n}\else\textsubscript{\ensuremath{n}}\fi}
\newunicodechar{ₐ}{\ifmmode_{a}\else\textsubscript{\ensuremath{a}}\fi}
\newunicodechar{ᵢ}{\ifmmode_{i}\else\textsubscript{\ensuremath{i}}\fi}
\newunicodechar{ᵣ}{\ifmmode_{r}\else\textsubscript{\ensuremath{r}}\fi}
\newunicodechar{ₖ}{\ifmmode_{k}\else\textsubscript{\ensuremath{k}}\fi}
\newunicodechar{ᵦ}{\ifmmode_{\beta}\else\textsubscript{\ensuremath{\beta}}\fi}
\newunicodechar{ₓ}{\ifmmode_{x}\else\textsubscript{\ensuremath{x}}\fi}
\newfontfamily\popdisplay{ArchivoBlack-Regular.ttf}[Path=../../../fonts/]
\newfontfamily\poplogo{SpaceGrotesk-Bold.ttf}[Path=../../../fonts/]
\newfontfamily\popsemi{PlusJakartaSans-SemiBold.ttf}[Path=../../../fonts/]
\newfontfamily\popxbold{PlusJakartaSans-ExtraBold.ttf}[Path=../../../fonts/]
\newfontfamily\popxboldls{PlusJakartaSans-ExtraBold.ttf}[Path=../../../fonts/,LetterSpace=3.0]

\setlist[enumerate]{leftmargin=1.7em,itemsep=5pt,topsep=4pt}
\setlist[itemize]{leftmargin=1.7em,itemsep=4pt,topsep=4pt,label={\color{poporange}\textbullet}}
\setlength{\parindent}{0pt}
\setlength{\parskip}{5pt}
\renewcommand{\arraystretch}{1.25}

% pgfplots : style Pop par défaut
\pgfplotsset{
  every axis/.append style={axis line style={popink,line width=1pt},tick style={popink},
    grid style={popline,line width=0.5pt},label style={font=\small},tick label style={font=\footnotesize}},
  popcurve/.style={poporange,line width=1.8pt,no markers,smooth},
}
% Même style disponible dans un \draw TikZ simple (hors pgfplots)
\tikzset{popcurve/.style={poporange,line width=1.8pt}}
\tikzset{popgrid/.style={popline,very thin}}
\colorlet{popcurve}{poporange} % le nom du style est parfois utilisé comme couleur

% ============================================================
% Pill / squiggle
% ============================================================
\newcommand{\poppill}[3]{%
  \tikz[baseline=-0.55ex]\node[draw=popink,line width=1.2pt,rounded corners=2.6pt,%
    fill=#2,inner xsep=6.5pt,inner ysep=3pt]{%
    \popxboldls\fontsize{7}{7}\selectfont\color{#3}\MakeUppercase{#1}};%
}
\newcommand{\popsquiggle}[1]{%
  \tikz[baseline=-0.4ex]{
    \draw[#1,line width=2.6pt,line cap=round]
      plot [smooth] coordinates {(0,0.09) (0.34,0.01) (0.68,0.21) (1.02,0.01) (1.36,0.10)};
  }%
}
% Mot-clé surligné (marqueur orange sous le texte)
\newcommand{\cle}[1]{{\popxbold\color{popdark}#1}}

% ============================================================
% En-tête / pied
% ============================================================
\pagestyle{fancy}
\fancyhf{}
\renewcommand{\headrulewidth}{0pt}
\renewcommand{\footrulewidth}{0pt}
\lhead{\raisebox{-0.15ex}{\poplogo\fontsize{13}{13}\selectfont\color{popink}OnBuch\color{poporange}+}}
\rhead{\poppill{\DOCMATIERE\ \textbullet\ \DOCNIVEAU\ \textbullet\ Leçon \DOCLECON}{white}{popink}}
\lfoot{\popxbold\fontsize{7.6}{7.6}\selectfont\color{popmuted}\MakeUppercase{Cours signé OnBuch+}\ \textbullet\ {\popsemi\DOCTITRE}}
\rfoot{\tikz[baseline=-0.6ex]\node[rounded corners=8.5pt,fill=popink,minimum height=17pt,minimum width=17pt,inner xsep=5pt,inner sep=0]{\popxbold\fontsize{8}{8}\selectfont\color{popcream}\thepage\,/\,\pageref*{LastPage}};}

% ============================================================
% Titres
% ============================================================
\newcommand{\chaptitre}[2]{%
  \begin{tcolorbox}[enhanced,colback=poporange,colframe=popink,boxrule=2.6pt,arc=11pt,
      drop shadow=popink,left=15pt,right=15pt,top=12pt,bottom=11pt,before skip=3pt,after skip=18pt]
    \poppill{#2}{popink}{popcream}\par\vspace{9pt}
    {\raggedright\hyphenpenalty=10000\exhyphenpenalty=10000\popdisplay\fontsize{22}{24}\selectfont\color{popcream}\MakeUppercase{#1}\par}
  \end{tcolorbox}
}
% Section numérotée : pastille orange avec le numéro + titre Archivo
\newcounter{popsec}
\newcommand{\coursec}[1]{%
  \stepcounter{popsec}\setcounter{popsub}{0}%
  \par\vspace{16pt}\noindent
  \begin{minipage}{\linewidth}
  \tikz[baseline=-0.7ex]\node[draw=popink,line width=1.6pt,fill=poporange,rounded corners=4pt,minimum size=22pt,inner sep=0]{\popdisplay\fontsize{12}{12}\selectfont\color{popcream}\thepopsec};%
  \hspace{8pt}{\popdisplay\fontsize{15}{17}\selectfont\color{popink}\MakeUppercase{#1}}\par
  \vspace{4pt}\noindent{\color{poporange}\rule{36pt}{3.6pt}}
  \end{minipage}\par\nopagebreak\vspace{8pt}
}
\newcounter{popsub}
\newcommand{\courssub}[1]{%
  \stepcounter{popsub}%
  \par\vspace{9pt}\noindent{\popxbold\fontsize{11.5}{13}\selectfont\color{popink}{\color{poporange}\thepopsec.\thepopsub}\hspace{6pt}#1}\par\nopagebreak\vspace{4pt}
}

% ============================================================
% Boîtes pédagogiques — même recette : fond blanc, bordure épaisse colorée,
% ombre dure encre, badge plein. Titre optionnel [#1].
% ============================================================
\tcbset{popbase/.style={breakable,enhanced,colback=white,boxrule=2.3pt,arc=9pt,
  shadow={3pt}{-3pt}{0pt}{popink},left=14pt,right=14pt,top=11pt,bottom=11pt,before skip=11pt,after skip=15pt,
  fontupper=\popsemi\fontsize{10.6}{14.5}\selectfont\color{popink},
  fontlower=\popsemi\fontsize{10.6}{14.5}\selectfont\color{popink}}}
% Coche dessinée (le glyphe ✓ n'existe pas dans les polices du document)
\newcommand{\popcheck}[1]{\tikz[baseline=-0.5ex]\draw[#1,line width=1.6pt,line cap=round,line join=round](-0.13,0.0)--(-0.04,-0.09)--(0.13,0.1);}
\providecommand{\coche}{\popcheck{popgreen}}
\providecommand{\resultat}[1]{\par\smallskip\noindent\fcolorbox{popgreen}{popgreenL}{\parbox{\dimexpr\linewidth-2\fboxsep-2\fboxrule\relax}{\textbf{Résultat.} #1}}\par\smallskip}
\newcommand{\popboxhead}[3]{\poppill{#1}{#2}{white}\ifx\relax#3\relax\else\hspace{6pt}{\popxbold\fontsize{10.6}{12}\selectfont\color{popink}#3}\fi\par\vspace{7pt}}

\newtcolorbox{aretenir}[1][]{popbase,colframe=popgreen,before upper={\popboxhead{À retenir}{popgreen}{#1}}}
% Texte en chinois (document de lecture, dialogue, extrait) : cadre
% d'étude. Un titre optionnel (source, auteur) s'affiche dans l'en-tête.
\newtcolorbox{textebox}[1][]{popbase,colframe=popblue,before upper={\popboxhead{文本 · Texte}{popblue}{#1}},after upper={}}
% Marques ✓ / ✗ (forme correcte / fautive) souvent utilisées par les modèles
\providecommand{\cross}{\textcolor{poppink}{\ensuremath{\times}}}
\providecommand{\xmark}{\cross}\providecommand{\crossmark}{\cross}\providecommand{\cmark}{\ensuremath{\checkmark}}
% Ligne de réponse / justification à compléter (inventée par les modèles)
\providecommand{\justification}[1]{\par\noindent\textit{Justification~:} \dotfill\par}
% \textipa (package tipa non chargé) : la police ne couvre pas l'API -> simple passage
\providecommand{\textipa}[1]{#1}
% Soulignement (ulem non chargé) : \uline, \ul
\providecommand{\uline}[1]{\underline{#1}}\providecommand{\ul}[1]{\underline{#1}}
% \glossaire{\item ...} inventée par les modèles : liste de glossaire
\providecommand{\glossaire}[1]{\begin{itemize}#1\end{itemize}}
% \alert (beamer) inventée par les modèles : texte mis en évidence
\providecommand{\alert}[1]{\textbf{\textcolor{poppink}{#1}}}
% variantes ASCII d'ordinaux écrites en macro
\providecommand{\iere}{\textsuperscript{re}}\providecommand{\ieme}{\textsuperscript{e}}\providecommand{\ere}{\textsuperscript{re}}\providecommand{\eme}{\textsuperscript{e}}\providecommand{\er}{\textsuperscript{er}}\providecommand{\ie}{\textsuperscript{e}}
% Ordinaux français écrits en macro par les modèles : 1\ière, 3\ième
\newcommand{\ière}{\textsuperscript{re}}\newcommand{\ième}{\textsuperscript{e}}
% \exemple (inventée par les modèles) : étiquette « Exemple : »
\providecommand{\exemple}{\textbf{Exemple~:}\ }
% \square utilisable aussi hors mode math (cases à cocher)
\AtBeginDocument{\let\squareMath\square\renewcommand{\square}{\ensuremath{\squareMath}}}
% Mot ou phrase en chinois à l'intérieur d'un paragraphe français
\newcommand{\zh}[1]{\ifmmode\text{#1}\else{#1}\fi}
\let\eng\zh \let\esp\zh \let\all\zh % alias tolérés
% flèches d'intonation inventées par les modèles
\providecommand{\textsearrow}{}\renewcommand{\textsearrow}{\ensuremath{\searrow}}\providecommand{\textnearrow}{}\renewcommand{\textnearrow}{\ensuremath{\nearrow}}
% \fr{..} : mot français (inventé par les modèles)
\providecommand{\fr}[1]{\textit{#1}}
% zhbox : encadré de texte chinois inventé par les modèles
\newenvironment{zhbox}{\begin{quote}}{\end{quote}}
% \vs : « versus » inventé par les modèles
\providecommand{\vs}{\textit{vs}\ }
% \blank : case à compléter inventée par les modèles
\providecommand{\blank}[1][]{\underline{\hspace{1.6em}}}
\newtcolorbox{exemplebox}[1][]{popbase,colframe=poporange,before upper={\popboxhead{Exemple}{poporange}{#1}}}
\newtcolorbox{definition}[1][]{popbase,colframe=popblue,before upper={\popboxhead{Définition}{popblue}{#1}}}
\newtcolorbox{propriete}[1][]{popbase,colframe=poppurple,before upper={\popboxhead{Propriété}{poppurple}{#1}}}
\newtcolorbox{methode}[1][]{popbase,colframe=popgold,before upper={\popboxhead{Méthode}{popgold}{#1}}}
\newtcolorbox{attention}[1][]{popbase,colframe=poppink,before upper={\popboxhead{Attention}{poppink}{#1}}}
\newtcolorbox{experience}[1][]{popbase,colframe=popblue,colback=popblueL,before upper={\popboxhead{Expérience}{popblue}{#1}}}
% « Le savais-tu ? » : carte violette pleine (contexte, culture, Cameroun)
\newtcolorbox{savaistu}[1][]{popbase,colback=poppurple,colframe=popink,
  fontupper=\popsemi\fontsize{10.4}{14.2}\selectfont\color{white},
  before upper={\poppill{Le savais-tu\,?}{white}{poppurple}\ifx\relax#1\relax\else\hspace{6pt}{\popxbold\color{white}#1}\fi\par\vspace{7pt}}}
% Exercice résolu : énoncé en haut, \tcblower puis la solution sur fond crème
\newcounter{popexo}\newcounter{popexr}
\newtcolorbox{exoresolu}[1][]{popbase,colframe=poporange,colbacklower=popcream,
  segmentation style={popink,line width=1.2pt,dashed},
  before upper={\stepcounter{popexr}\popboxhead{Exercice résolu \thepopexr}{poporange}{#1}},
  before lower={\poppill{Solution}{popink}{popcream}\par\vspace{6pt}}}
% Exercice (non résolu en place) — la correction est dans la partie Corrigés
\newtcolorbox{exercice}[1][]{popbase,colframe=popink,
  before upper={\stepcounter{popexo}\popboxhead{Exercice \thepopexo}{popink}{#1}}}
% Corrigé
\newtcolorbox{corrige}[1][]{popbase,colframe=popgreen,colback=popgreenL,
  before upper={\popboxhead{Corrigé}{popgreen}{#1}}}
% Objectifs / prérequis / bilan
\newtcolorbox{objectifs}{popbase,colframe=popink,colback=poporangeL,
  before upper={\popboxhead{Objectifs de la leçon}{popink}{}}}
\newtcolorbox{prerequis}{popbase,colframe=popink,colback=popblueL,
  before upper={\popboxhead{Prérequis}{popblue}{}}}
\newtcolorbox{bilan}{popbase,colframe=popink,colback=white,boxrule=2.8pt,
  before upper={\popboxhead{Fiche bilan}{poporange}{L'essentiel à connaître pour l'examen}}}
% Figure encadrée avec légende
\newtcolorbox{popfigure}[1][]{enhanced,breakable=false,colback=white,colframe=popline,boxrule=1.4pt,arc=8pt,
  left=10pt,right=10pt,top=10pt,bottom=8pt,before skip=10pt,after skip=12pt,halign=center,
  lower separated=false,halign lower=center,
  fontlower=\popsemi\fontsize{9}{11.5}\selectfont\color{popmuted},
  #1}
\newcommand{\legende}[1]{\par\vspace{6pt}{\popsemi\fontsize{9}{11.5}\selectfont\color{popmuted}#1}}

% ============================================================
% Signature OnBuch+
% ============================================================
\newcommand{\poplogoinline}[1]{{\poplogo\fontsize{#1}{#1}\selectfont\color{popink}OnBuch\color{poporange}+}}
\newcommand{\signatureonbuch}{%
  \begin{tcolorbox}[enhanced,colback=popink,colframe=popink,boxrule=2.2pt,arc=10pt,
      drop shadow=poporange,left=14pt,right=14pt,top=10pt,bottom=10pt,before skip=0pt,after skip=0pt]
    \tikz[baseline=-0.7ex]\node[circle,fill=popgreen,draw=popcream,line width=1.3pt,minimum size=22pt,inner sep=0]{\popcheck{white}};%
    \hspace{9pt}\begin{minipage}[c]{0.62\linewidth}
      {\popxbold\fontsize{12}{13}\selectfont\color{popcream}Signé\ \textbullet\ {\poplogo OnBuch\color{poporange}+}}\par\vspace{2pt}
      {\raggedright\popsemi\fontsize{8}{10.5}\selectfont\color{popline}\MakeUppercase{Équipe pédagogique OnBuch+}\\\MakeUppercase{Conforme au programme officiel MINESEC}\par}
    \end{minipage}\hfill
    {\poplogo\fontsize{22}{22}\selectfont\color{popcream}OnBuch\color{poporange}+}
  \end{tcolorbox}%
}

% ============================================================
% Page de garde
% #1 titre · #2 matière · #3 niveau · #4 module · #5 n° leçon · #6 durée ·
% #7 description / objectifs (texte ou itemize)
% ============================================================
\newcommand{\pagegarde}[7]{%
  \thispagestyle{empty}
  \newgeometry{a4paper,margin=1.7cm,top=1.6cm,bottom=1.4cm}%
  \begin{tikzpicture}[remember picture,overlay]
    \fill[poporange!18] ([xshift=-3.2cm,yshift=-3.6cm]current page.north east) circle (4.6cm);
    \fill[poppurple!14] ([xshift=2.4cm,yshift=7.6cm]current page.south west) circle (3.8cm);
  \end{tikzpicture}%
  {\poplogo\fontsize{30}{30}\selectfont\color{popink}OnBuch\color{poporange}+}\hfill
  \raisebox{0.6ex}{\poppill{Cours signé \textbullet\ Premium}{popink}{popcream}}\par
  \vspace{1pt}\popsquiggle{poppurple}\par
  \vspace{8pt}
  \begin{tcolorbox}[enhanced,colback=poporange,colframe=popink,boxrule=2.8pt,arc=13pt,
      drop shadow=popink,left=18pt,right=18pt,top=12pt,bottom=13pt,after skip=10pt]
    \poppill{#2 \textbullet\ #3}{popink}{popcream}\hspace{5pt}\poppill{#4}{white}{popink}\par\vspace{8pt}
    {\popdisplay\fontsize{14}{14}\selectfont\color{popink}LEÇON #5}\par\vspace{4pt}
    {\raggedright\hyphenpenalty=10000\exhyphenpenalty=10000\popdisplay\fontsize{25}{27}\selectfont\color{popcream}\MakeUppercase{#1}\par}
    \vspace{8pt}
    {\popxbold\fontsize{9.5}{9.5}\selectfont\color{popink}\MakeUppercase{Durée conseillée : #6}}
  \end{tcolorbox}
  \begin{tcolorbox}[enhanced,colback=white,colframe=popink,boxrule=2.2pt,arc=10pt,
      drop shadow=popink,left=16pt,right=16pt,top=10pt,bottom=10pt,after skip=8pt]
    \poppill{Dans cette leçon}{poppurple}{white}\par\vspace{5pt}
    \popsemi\fontsize{9.3}{12.6}\selectfont\color{popink}#7
  \end{tcolorbox}
  \noindent\begin{minipage}[t]{0.487\linewidth}
    \begin{tcolorbox}[enhanced,colback=popgreen,colframe=popink,boxrule=2.2pt,arc=10pt,
        drop shadow=popink,left=11pt,right=11pt,top=10pt,bottom=10pt,height=3.1cm,valign=center]
      \tcbox[colback=white,colframe=popink,boxrule=1.3pt,arc=5pt,boxsep=0pt,nobeforeafter,
        left=3pt,right=3pt,top=3pt,bottom=3pt,valign=center]{\qrcode[height=1.9cm]{https://play.google.com/store/apps/details?id=com.luvvix.onbuch&hl=fr}}%
      \hspace{8pt}\begin{minipage}[c]{0.5\linewidth}
        {\popdisplay\fontsize{11}{12.5}\selectfont\color{white}\MakeUppercase{Télécharge OnBuch}}\par\vspace{4pt}
        {\raggedright\popsemi\fontsize{8.4}{11}\selectfont\color{white}Gratuit \textbullet\ Google Play\\Tuteur IA, annales, concours, résultats\par}
      \end{minipage}
    \end{tcolorbox}
  \end{minipage}\hfill
  \begin{minipage}[t]{0.487\linewidth}
    \begin{tcolorbox}[enhanced,colback=poppurple,colframe=popink,boxrule=2.2pt,arc=10pt,
        drop shadow=popink,left=11pt,right=11pt,top=10pt,bottom=10pt,height=3.1cm,valign=center]
      \tikz[baseline=-0.7ex]\node[circle,fill=white,draw=popink,line width=1.3pt,minimum size=1.6cm,inner sep=0]{\popdisplay\fontsize{24}{24}\selectfont\color{poppurple}+};%
      \hspace{8pt}\begin{minipage}[c]{0.6\linewidth}
        {\popdisplay\fontsize{11}{12.5}\selectfont\color{white}\MakeUppercase{Passe à OnBuch+}}\par\vspace{4pt}
        {\raggedright\popsemi\fontsize{8.4}{11}\selectfont\color{white}Tous les cours signés, Tuteur IA illimité, fiches PDF — onglet OnBuch+ de l'app\par}
      \end{minipage}
    \end{tcolorbox}
  \end{minipage}
  \vspace{8pt}
  \popcontenu
  \vfill
  \signatureonbuch
  \restoregeometry
  \newpage
}

% Rangée « Ce cours contient » (page de garde)
\newcommand{\poptile}[3]{% #1 couleur #2 gros texte #3 légende
  \begin{tcolorbox}[enhanced,colback=#1,colframe=popink,boxrule=2pt,arc=9pt,shadow={2.5pt}{-2.5pt}{0pt}{popink},
      left=6pt,right=6pt,top=9pt,bottom=9pt,halign=center,valign=center,height=1.75cm,width=\linewidth,nobeforeafter]
    {\popdisplay\fontsize{12.5}{13}\selectfont\color{white}\MakeUppercase{#2}}\par\vspace{3pt}
    {\centering\popsemi\fontsize{7.8}{9.5}\selectfont\color{white}#3\par}
  \end{tcolorbox}}
\newcommand{\popcontenu}{%
  \par\noindent{\popxboldls\fontsize{7.5}{7.5}\selectfont\color{popmuted}CE COURS SIGNÉ CONTIENT}\par\vspace{7pt}
  \noindent\begin{minipage}[t]{0.235\linewidth}\poptile{popblue}{Cours}{complet, illustré, pas à pas}\end{minipage}\hfill
  \begin{minipage}[t]{0.235\linewidth}\poptile{poporange}{Méthodes}{et exercices résolus}\end{minipage}\hfill
  \begin{minipage}[t]{0.235\linewidth}\poptile{poppink}{Exercices}{type examen + corrigés}\end{minipage}\hfill
  \begin{minipage}[t]{0.235\linewidth}\poptile{popgreen}{Fiche bilan}{l'essentiel à retenir}\end{minipage}\par}

% Sommaire
\newtcolorbox{sommaire}{enhanced,colback=white,colframe=popink,boxrule=2.2pt,arc=10pt,
  drop shadow=popink,left=18pt,right=18pt,top=13pt,bottom=13pt,before skip=6pt,after skip=20pt,
  before upper={\poppill{Sommaire}{poppurple}{white}\par\vspace{9pt}\popsemi\fontsize{10.6}{14.5}\selectfont\color{popink}}}

% Signature de fin de cours — poussée tout en bas de la dernière page
\newcommand{\findecours}{%
  \par\vfill\nopagebreak
  \begin{center}\popsquiggle{poporange}\end{center}\vspace{4pt}
  \signatureonbuch
}
\AtBeginDocument{\def\llbracket{\mathopen{[\mkern-3.5mu[}}\def\rrbracket{\mathclose{]\mkern-3.5mu]}}} % intervalles d'entiers (glyphe ⟦ absent de la police)
% Glyphes absents de Fira Math : redéfinis à partir de glyphes disponibles
\AtBeginDocument{%
  \def\setminus{\mathbin{\backslash}}\let\smallsetminus\setminus
  \def\vdots{\mathord{\vbox{\baselineskip3.2pt\lineskiplimit0pt\kern4pt\hbox{.}\hbox{.}\hbox{.}}}}%
  \def\ddots{\mathinner{\mkern1mu\raise6.5pt\hbox{.}\mkern2mu\raise3.6pt\hbox{.}\mkern2mu\raise0.7pt\hbox{.}\mkern1mu}}%
  \def\triangle{\mathord{\scalebox{0.9}{$\symup{\Delta}$}}}%
  \def\nabla{\mathord{\raisebox{\depth}{\scalebox{1}[-1]{$\symup{\Delta}$}}}}%
  \def\checkmark{\popcheck{popdark}}%
  \def\star{\mathbin{\ast}}\def\bigstar{\mathord{\ast}}%
  \def\diamond{\mathbin{\scalebox{0.6}{$\Diamond$}}}%
  \def\bigcup{\mathop{\mathchoice{\raisebox{-0.3em}{\scalebox{1.7}{$\cup$}}}{\scalebox{1.25}{$\cup$}}{\cup}{\cup}}\displaylimits}%
  \def\bigcap{\mathop{\mathchoice{\raisebox{-0.3em}{\scalebox{1.7}{$\cap$}}}{\scalebox{1.25}{$\cap$}}{\cap}{\cap}}\displaylimits}%
  \def\lesseqqgtr{\mathrel{\lessgtr}}\def\bigoplus{\mathop{\oplus}\displaylimits}%
}
\providecommand{\ssi}{si et seulement si} % abréviation fréquente
\AtBeginDocument{\providecommand{\wideparen}{\overparen}} % arc de cercle

%% FIGFIX-BEGIN v3 — correctifs globaux des figures (docs/onbuch-generation/tools/figfix)
\usepackage{adjustbox}\usepackage{etoolbox}
% toute figure TikZ/pgfplots plus large que la ligne est réduite à la largeur disponible
\BeforeBeginEnvironment{tikzpicture}{\begin{adjustbox}{max width=\linewidth}}
\AfterEndEnvironment{tikzpicture}{\end{adjustbox}}
% légendes pgfplots lisibles par-dessus les courbes
\pgfplotsset{every axis/.append style={legend style={fill=white,fill opacity=0.92,text opacity=1,draw=black!15,inner sep=3pt}}}
% étiquettes d'arêtes (syntaxe edge["..."]) avec fond blanc
\tikzset{every edge quotes/.append style={fill=white,inner sep=1.2pt}}
% bulles de cartes mentales : texte à la ligne dans la bulle, bulle un peu plus grande
\tikzset{every concept/.append style={text width=2.0cm,align=center,minimum size=2.3cm,inner sep=2pt}}
%% FIGFIX-END
\begin{document}

\pagegarde{Leçon 4 — Expression écrite : écriture des caractères de l'amour et de la famille ; types de phrases étudiées}{Chinois}{1ère A}{Module 1 : Vie familiale et intégration sociale}{ZHA04}{2 h}{Cette leçon d'expression écrite apprend à tracer correctement les caractères du thème de l'amour et de la famille (traits, radicaux, caractères proches) et à rédiger un texte simple en chinois à l'aide de connecteurs. Elle entraîne aussi à la traduction dans les deux sens (thème et version), compétence clé de l'examen.
\vspace{4pt}
\begin{multicols}{2}\small
\begin{itemize}
\item À la fin de cette leçon, je sais écrire correctement les caractères du thème de l'amour et de la famille en respectant l'ordre des traits.
\item Je sais identifier les radicaux des caractères étudiés et les utiliser pour mémoriser et distinguer les caractères proches.
\item Je sais distinguer les caractères faciles à confondre (爱/受, 妈/吗, 他/她/它).
\item Je sais rédiger un texte simple de 5 à 8 phrases sur ma famille en chinois.
\item Je sais utiliser les connecteurs 和, 也, 但是, 因为…所以… pour lier mes phrases.
\item Je sais traduire des phrases simples du français vers le chinois (thème).
\end{itemize}\end{multicols}}

\begin{sommaire}
\begin{multicols}{2}
\begin{enumerate}
\item Modèle de texte commenté
\item Écriture des caractères : traits et radicaux
\item Caractères proches : ne pas confondre
\item Structure et connecteurs du texte
\item Thème : français → chinois
\item Version : chinois → français
\item Grille d'évaluation et entraînement à l'examen
\item Activité d'intégration
\item Méthodes et exercices résolus
\item Exercices
\item Corrigés des exercices
\item Fiche bilan
\end{enumerate}
\end{multicols}
\end{sommaire}

\begin{objectifs}
\begin{itemize}
\item À la fin de cette leçon, je sais écrire correctement les caractères du thème de l'amour et de la famille en respectant l'ordre des traits.
\item Je sais identifier les radicaux des caractères étudiés et les utiliser pour mémoriser et distinguer les caractères proches.
\item Je sais distinguer les caractères faciles à confondre (爱/受, 妈/吗, 他/她/它).
\item Je sais rédiger un texte simple de 5 à 8 phrases sur ma famille en chinois.
\item Je sais utiliser les connecteurs 和, 也, 但是, 因为…所以… pour lier mes phrases.
\item Je sais traduire des phrases simples du français vers le chinois (thème).
\item Je sais traduire des phrases simples du chinois vers le français (version).
\item Je sais appliquer les types de phrases étudiés (déclarative, interrogative avec 吗/呢, négative avec 不/没) dans mes productions.
\end{itemize}
\end{objectifs}

\begin{prerequis}
\begin{itemize}
\item Vocabulaire de la famille : 爸爸, 妈妈, 哥哥, 姐姐, 弟弟, 妹妹, 家, 爱 (leçons ZHA01-ZHA03)
\item Structure de base Sujet + Verbe + Complément et la phrase avec 是
\item Négations 不 et 没, particule interrogative 吗
\item Règles générales d'ordre des traits (de haut en bas, de gauche à droite)
\item Adjectifs simples : 好, 大, 小, 漂亮, 高兴
\end{itemize}
\end{prerequis}

\newpage

\coursec{Modèle de texte commenté}

\courssub{Situation d'accroche et problématique}

Pour la fête des mères, Amina, élève de Première A au lycée de Yaoundé, veut écrire une carte en chinois à sa correspondante de Chengdu. Elle a déjà préparé ses phrases en français : « J'aime ma famille. Ma mère est gentille, mon père est travailleur. » Mais devant sa feuille blanche, elle hésite : comment écrire \zh{爱} (ài, aimer) sans le confondre avec \zh{受} (shòu, recevoir) ? Comment construire correctement une phrase en chinois, dans le bon ordre, avec les bons mots-outils ? Comment donner à son petit texte un début, un milieu et une fin, comme le font les élèves chinois ?

Cette leçon lui donne les clés pour écrire \emph{juste, beau et clair}. Nous allons d'abord lire ensemble un texte modèle, le comprendre, le traduire, puis observer comment il est construit : sa structure, ses types de phrases, ses mots de liaison. Ce modèle servira ensuite de plan pour la propre production d'Amina — et pour la vôtre.

\textbf{Problématique :} comment un texte simple en chinois est-il organisé, et quels outils (structures, connecteurs, types de phrases) faut-il maîtriser pour le comprendre et l'imiter ?

\courssub{Lecture du texte modèle}

\begin{textebox}[我爱我的家 — Texte modèle]
我爱我的家。\\
Wǒ ài wǒ de jiā.\\
« J'aime ma famille. »\\[0.5em]
我家有五口人：爸爸、妈妈、哥哥、妹妹和我。\\
Wǒ jiā yǒu wǔ kǒu rén : bàba, māma, gēge, mèimei hé wǒ.\\
« Ma famille compte cinq personnes : papa, maman, mon grand frère, ma petite sœur et moi. »\\[0.5em]
我的爸爸是老师，我的妈妈是医生。\\
Wǒ de bàba shì lǎoshī, wǒ de māma shì yīshēng.\\
« Mon père est professeur, ma mère est médecin. »\\[0.5em]
哥哥很喜欢足球，妹妹很可爱。\\
Gēge hěn xǐhuan zúqiú, mèimei hěn kě'ài.\\
« Mon frère aime beaucoup le football, ma sœur est très mignonne. »\\[0.5em]
我们都爱我们的家。\\
Wǒmen dōu ài wǒmen de jiā.\\
« Nous aimons tous notre famille. »
\end{textebox}

\begin{definition}[Glossaire du texte modèle]
\begin{center}
\begin{tabularx}{\linewidth}{|X|X|X|X|}
\hline
\rowcolor{popblueL} Caractères & Pinyin & Nature & Traduction \\
\hline
爱 & ài & verbe & aimer \\
\hline
家 & jiā & nom & famille, maison \\
\hline
有 & yǒu & verbe & avoir, il y a \\
\hline
口 & kǒu & classificateur / nom & (classificateur des membres de la famille) ; bouche \\
\hline
人 & rén & nom & personne \\
\hline
爸爸 & bàba & nom & papa \\
\hline
妈妈 & māma & nom & maman \\
\hline
哥哥 & gēge & nom & grand frère \\
\hline
妹妹 & mèimei & nom & petite sœur \\
\hline
和 & hé & conjonction & et \\
\hline
是 & shì & verbe & être \\
\hline
老师 & lǎoshī & nom & professeur, enseignant \\
\hline
医生 & yīshēng & nom & médecin \\
\hline
很 & hěn & adverbe & très \\
\hline
喜欢 & xǐhuan & verbe & aimer, apprécier \\
\hline
足球 & zúqiú & nom & football \\
\hline
可爱 & kě'ài & adjectif & mignon, adorable \\
\hline
都 & dōu & adverbe & tous, tout \\
\hline
我们 & wǒmen & pronom & nous \\
\hline
的 & de & particule & (marque de possession : de) \\
\hline
\end{tabularx}
\end{center}
\end{definition}

\courssub{Repérage de la structure du texte}

Lisons le texte une deuxième fois, non plus pour le traduire, mais pour observer \emph{comment il est fabriqué}. Un bon texte chinois, même très court, suit un plan clair. Notre modèle comporte cinq étapes :

\begin{enumerate}
\item \textbf{Phrase d'ouverture} : \zh{我爱我的家。} — l'auteur annonce son sujet et son sentiment. C'est la porte d'entrée du texte.
\item \textbf{Présentation des membres} : \zh{我家有五口人：爸爸、妈妈、哥哥、妹妹和我。} — on découvre qui compose la famille.
\item \textbf{Professions} : \zh{我的爸爸是老师，我的妈妈是医生。} — on précise ce que font les parents.
\item \textbf{Goûts et qualités} : \zh{哥哥很喜欢足球，妹妹很可爱。} — on ajoute un détail vivant sur les enfants.
\item \textbf{Phrase de conclusion} : \zh{我们都爱我们的家。} — on boucle le texte en reprenant l'idée du début (\zh{爱} revient, comme un écho).
\end{enumerate}

Remarquez l'élégance de la construction : la dernière phrase répond à la première. On appelle cela une \cle{conclusion en écho}. C'est une technique très appréciée dans l'écriture chinoise, et très facile à réutiliser dans vos propres textes.

\begin{popfigure}
\begin{tikzpicture}[mindmap, grow cyclic,
  every node/.style={concept, align=center, text=white},
  root concept/.append style={concept color=popblue, font=\Large\bfseries},
  level 1/.append style={level distance=3.2cm, sibling angle=72},
  level 1 concept/.append style={font=\small},
  concept color=popblue]
  \node[root concept, align=center]{我的家\\ Wǒ de jiā}
    child[concept color=poporange] { node[align=center]{ouverture\\ 我爱我的家} }
    child[concept color=popgreen] { node[align=center]{membres\\ 五口人} }
    child[concept color=poppurple] { node[align=center]{professions\\ 老师、医生} }
    child[concept color=poppink] { node[align=center]{goûts\\ 足球、可爱} }
    child[concept color=popgold] { node[align=center]{conclusion\\ 都爱我们的家} };
\end{tikzpicture}
\legende{Carte mentale de la structure du texte modèle : cinq étapes autour du thème « ma famille ».}
\end{popfigure}

\begin{methode}[Analyser un texte modèle en quatre étapes]
\begin{enumerate}
\item \textbf{Lire et traduire.} Lis le texte à voix haute avec le pinyin, puis traduis-le phrase par phrase. Vérifie que chaque phrase est comprise avant de passer à la suivante.
\item \textbf{Repérer la structure.} Découpe le texte en parties (ouverture, développement, conclusion) et donne un titre à chaque partie. Demande-toi : « Que fait cette phrase dans le texte ? »
\item \textbf{Souligner les connecteurs et les mots-outils.} Entoure les mots de liaison (\zh{和}, \zh{都}), la particule \zh{的}, les verbes de structure (\zh{是}, \zh{有}). Ce sont les « charnières » du texte.
\item \textbf{Réutiliser le plan.} Recopie le plan à trous (ouverture → membres → professions → goûts → conclusion) et remplis-le avec \emph{ta} famille, tes mots, tes phrases. Le modèle devient ton gabarit.
\end{enumerate}
\end{methode}

\coursec{Écriture des caractères : traits et radicaux}

\courssub{Principes généraux de l'ordre des traits}

Avant d'étudier les caractères du thème, rappelons les règles fondamentales qui garantissent une écriture lisible, fluide et conforme aux standards (HSK / examen officiel) :

\begin{propriete}[Règles d'ordre des traits (ordre standard)]
\begin{enumerate}
\item \textbf{Haut avant bas} (从上到下) : \zh{三} (sān), \zh{王} (wáng).
\item \textbf{Gauche avant droite} (从左到右) : \zh{好} (hǎo), \zh{妈} (mā).
\item \textbf{Horizontal avant vertical} (先横后竖) : \zh{十} (shí), \zh{干} (gān).
\item \textbf{Extérieur avant intérieur, puis fermeture} (先外后内，再封口) : \zh{口} (kǒu), \zh{国} (guó), \zh{家} (jiā).
\item \textbf{Trait central avant les ailes} (先中间后两边) : \zh{小} (xiǎo), \zh{水} (shuǐ).
\item \textbf{Point en haut à droite à la fin} (右上点最后写) : \zh{文} (wén), \zh{武} (wǔ).
\end{enumerate}
\end{propriete}

\courssub{Analyse des caractères clés du thème}

Le tableau ci-dessous détaille l'écriture des caractères essentiels pour parler de sa famille et de l'amour. Maîtrisez leur ordre des traits et leur radical (clé de dictionnaire) : c'est la condition pour les retrouver dans un dictionnaire et pour les écrire vite et bien.

\begin{exemplebox}[Fiche d'écriture des caractères du thème]
\begin{center}
\begin{tabularx}{\linewidth}{|X|X|X|X|X|}
\hline
\rowcolor{popblueL} Caractère & Pinyin & Radical (Clé) & Nb traits & Ordre des traits (description) \\
\hline
爱 & ài & \zh{爫} (zhuǎ, griffe) & 10 & 1. \zh{爫} (3 traits : horizontale, crochet, horizontale) 2. \zh{冖} (mì, couverture, 1 trait) 3. \zh{友} (yǒu, 4 traits) 4. \zh{心} (xīn, cœur, 3 traits : point, point, crochet). \textbf{Astuce} : « L'amitié (\zh{友}) sous un toit (\zh{冖}) touche le cœur (\zh{心}). » \\
\hline
家 & jiā & \zh{宀} (mián, toit) & 10 & 1. \zh{宀} (3 traits : point, horizontale pliée, point) 2. \zh{豕} (shǐ, cochon, 7 traits : haut → bas, gauche → droite). \textbf{Astuce} : « Sous le toit (\zh{宀}), un cochon (\zh{豕}) = la famille (richesse d'autrefois). » \\
\hline
爸 & bà & \zh{父} (fù, père) & 8 & 1. \zh{父} (4 traits : horizontale, verticale, horizontale, horizontale) — \emph{radical à droite ici} 2. \zh{巴} (bā, 4 traits). \textbf{Note} : \zh{父} est à la fois radical et phonétique. \\
\hline
妈 & mā & \zh{女} (nǚ, femme) & 6 & 1. \zh{女} (3 traits : trait oblique, trait oblique, horizontale) 2. \zh{马} (mǎ, cheval, 3 traits). \textbf{Astuce} : « La femme (\zh{女}) sur le cheval (\zh{马}) = maman. » \\
\hline
哥 & gē & \zh{亻} (rén, homme) & 10 & 1. \zh{亻} (2 traits) 2. \zh{可} (kě, 5 traits : horizontale, verticale, horizontale, verticale, horizontale). \\
\hline
妹 & mèi & \zh{女} (nǚ, femme) & 8 & 1. \zh{女} (3 traits) 2. \zh{未} (wèi, 5 traits : horizontale, verticale, horizontale, verticale, horizontale). \\
\hline
口 & kǒu & \zh{口} (kǒu, bouche) & 3 & 1. Verticale gauche 2. Horizontale pliée (haut + droite) 3. Horizontale basse. \textbf{Attention} : pas un carré en 4 traits ! \\
\hline
人 & rén & \zh{人} (rén, homme) & 2 & 1. Trait oblique gauche 2. Trait oblique droite (traverse le premier). \\
\hline
\end{tabularx}
\end{center}
\end{exemplebox}

\begin{experience}[Atelier d'écriture guidée : « Mon arbre généalogique »]
\begin{enumerate}
\item Sur une feuille à carreaux (format chinois), écris chaque caractère du tableau ci-dessus \textbf{trois fois} en respectant scrupuleusement l'ordre des traits.
\item Entoure ton meilleur exemplaire de chaque caractère au stylo rouge.
\item Échange ta feuille avec ton voisin : il vérifie l'ordre des traits en te regardant écrire le caractère \zh{爱} et \zh{家} au tableau.
\item Défis bonus : écris \zh{弟} (dì, petit frère) et \zh{姐} (jiě, grande sœur) en déduisant leur structure à partir de \zh{哥} et \zh{妹} (radical \zh{亻} / \zh{女} + phonétique \zh{弟} / \zh{且}).
\end{enumerate}
\end{experience}

\coursec{Caractères proches : ne pas confondre}

\courssub{Les pièges visuels fréquents}

En chinois, une petite différence de trait ou de composant change radicalement le sens. Voici les paires les plus dangereuses pour le thème « famille / amour », spécialement sélectionnées pour les francophones.

\begin{attention}[Paires de caractères à ne pas confondre]
\begin{center}
\begin{tabularx}{\linewidth}{|X|X|X|X|X|}
\hline
\rowcolor{popblueL} Caractère 1 & Pinyin / Sens & Caractère 2 & Pinyin / Sens & Différence clé / Mnémotechnique \\
\hline
爱 & ài (aimer) & 受 & shòu (recevoir, subir) & \textbf{Bas} : \zh{心} (cœur) vs \zh{又} (main droite). \emph{Aimer, c'est donner son cœur ; recevoir, c'est tendre la main.} \\
\hline
好 & hǎo (bon) & 妹 & mèi (petite sœur) & \textbf{Gauche} : \zh{女} (femme) commun. \textbf{Droite} : \zh{子} (enfant) vs \zh{未} (pas encore). \emph{Une femme + un enfant = bon ; une femme + « pas encore » adulte = petite sœur.} \\
\hline
姐 & jiě (grande sœur) & 妹 & mèi (petite sœur) & \textbf{Droite} : \zh{且} (qiě, de plus) vs \zh{未} (wèi, pas encore). \emph{La grande sœur « en plus » (âgée) ; la petite sœur « pas encore » (jeune).} \\
\hline
爸 & bà (papa) & 父 & fù (père, formel) & \zh{爸} = \zh{父} (sens/phonétique) + \zh{巴} (phonétique). \zh{父} seul = terme formel / écrit. \\
\hline
妈 & mā (maman) & 马 & mǎ (cheval) & \zh{妈} = \zh{女} (femme) + \zh{马} (phonétique). \zh{马} seul = animal. \\
\hline
哥 & gē (grand frère) & 歌 & gē (chanter) & \zh{哥} = \zh{亻} (homme) + \zh{可} (phonétique). \zh{歌} = \zh{欠} (souffle/bouche) + \zh{可}. \emph{Le frère est un homme ; la chanson sort de la bouche.} \\
\hline
家 & jiā (famille/maison) & 豕 & shǐ (cochon, radical) & \zh{家} = \zh{宀} (toit) + \zh{豕}. \zh{豕} seul = radical « cochon » (rarement seul). \\
\hline
\end{tabularx}
\end{center}
\end{attention}

\begin{exercice}[Entraînement à la discrimination visuelle]
Parmi les paires suivantes, entoure le caractère qui correspond au sens donné.
\begin{enumerate}
\item « Aimer » : \zh{爱} \quad \zh{受}
\item « Petite sœur » : \zh{妹} \quad \zh{好}
\item « Papa (familier) » : \zh{爸} \quad \zh{父}
\item « Chanter » : \zh{哥} \quad \zh{歌}
\item « Famille » : \zh{家} \quad \zh{豕}
\end{enumerate}
\end{exercice}

\begin{corrige}[Exercice discrimination visuelle]
1. \zh{爱} (cœur en bas) 2. \zh{妹} (femme + « pas encore ») 3. \zh{爸} (radical père + 巴) 4. \zh{歌} (radical souffle/bouche) 5. \zh{家} (toit + cochon).
\end{corrige}

\coursec{Structure et connecteurs du texte}

\courssub{Observation des types de phrases}

Notre texte modèle utilise trois grands types de phrases déclaratives. Observons-les avant d'énoncer les règles.

\textbf{Type 1 : la phrase d'identité avec \zh{是}.}

\zh{我的妈妈是医生。} \emph{(Wǒ de māma shì yīshēng.)} « Ma mère est médecin. »

\begin{propriete}[La phrase Sujet + 是 + Nom]
Structure : \textbf{Sujet + 是 (shì) + nom / groupe nominal}.\\
Emploi : présenter une identité, une profession, une nationalité, une définition. \zh{是} relie deux noms (ou pronoms), \textbf{jamais} un nom et un adjectif.\\
Comparaison avec le français : le français dit « ma mère est \emph{gentille} » avec « être » ; le chinois dit \zh{我的妈妈很好} \emph{sans} \zh{是}, car devant un adjectif on utilise \zh{很} (ou un autre adverbe de degré).\\
Exemples : \zh{我是学生。} (Wǒ shì xuésheng.) « Je suis élève. » / \zh{他是喀麦隆人。} (Tā shì Kāmàilóng rén.) « Il est Camerounais. » / \zh{这是书。} (Zhè shì shū.) « Ceci est un livre. »
\end{propriete}

\textbf{Type 2 : la phrase de possession / existence avec \zh{有}.}

\zh{我家有五口人。} \emph{(Wǒ jiā yǒu wǔ kǒu rén.)} « Ma famille compte cinq personnes. »

\begin{propriete}[La phrase Sujet + 有 + Complément]
Structure : \textbf{Sujet + 有 (yǒu) + nombre + classificateur + nom}.\\
Emploi : exprimer la possession (« avoir ») ou l'existence (« il y a »). \textbf{Règle d'or} : le nombre ne se colle jamais directement au nom ; il faut un \cle{classificateur} (量词 liàngcí) entre les deux. Pour les membres d'une famille (et les personnes en général), le classificateur est \zh{口} (kǒu).\\
Négation : \zh{没有} (méiyǒu), \textbf{jamais} \zh{不有}. Exemple : \zh{我家没有哥哥。} (Wǒ jiā méiyǒu gēge.) « Ma famille n'a pas de grand frère. »\\
Autres classificateurs courants : \zh{个} (gè, général), \zh{本} (běn, livres), \zh{张} (zhāng, objets plats), \zh{只} (zhī, animaux).
\end{propriete}

\textbf{Type 3 : la phrase descriptive avec les adverbes \zh{很} et \zh{都}.}

\zh{哥哥很喜欢足球。} \emph{(Gēge hěn xǐhuan zúqiú.)} « Mon frère aime beaucoup le football. »\\
\zh{妹妹很可爱。} \emph{(Mèimei hěn kě'ài.)} « Ma sœur est très mignonne. »\\
\zh{我们都爱我们的家。} \emph{(Wǒmen dōu ài wǒmen de jiā.)} « Nous aimons tous notre famille. »

\begin{propriete}[Place des adverbes 很 et 都]
Structure : \textbf{Sujet + Adverbe + Verbe / Adjectif}.\\
L'adverbe se place \textbf{toujours avant} le verbe ou l'adjectif, et \textbf{après le sujet}.\\
\zh{很} (hěn, très) : accompagne presque toujours un adjectif seul (verbe adjectival) pour combler le vide syntaxique : \zh{妹妹很可爱} (pas *\zh{妹妹可爱}). Devant un verbe d'action, il marque l'intensité : \zh{很喜欢}.\\
\zh{都} (dōu, tous) : résume un groupe \textbf{pluriel} déjà mentionné (sujet ou objet antérieur). Il faut au moins deux entités dans le référent. \zh{都} se place avant le verbe (ou l'adjectif).\\
Comparaison avec le français : en français, « tous » peut se déplacer (« nous aimons tous » / « nous tous aimons ») ; en chinois, \zh{都} est \textbf{figé} en position préverbale.
\end{propriete}

\begin{exemplebox}[Tableau récapitulatif des types de phrases du modèle]
\begin{center}
\begin{tabularx}{\linewidth}{|X|X|X|}
\hline
\rowcolor{popblueL} Structure & Exemple en caractères + pinyin & Traduction \\
\hline
Sujet + 是 + Nom & 我的爸爸是老师。 Wǒ de bàba shì lǎoshī. & Mon père est professeur. \\
\hline
Sujet + 有 + Nombre + 口 + Nom & 我家有五口人。 Wǒ jiā yǒu wǔ kǒu rén. & Ma famille compte cinq personnes. \\
\hline
Sujet + 很 + Adjectif & 妹妹很可爱。 Mèimei hěn kě'ài. & Ma petite sœur est très mignonne. \\
\hline
Sujet + 很 + Verbe + Objet & 哥哥很喜欢足球。 Gēge hěn xǐhuan zúqiú. & Mon frère aime beaucoup le football. \\
\hline
Sujet + 都 + Verbe & 我们都爱我们的家。 Wǒmen dōu ài wǒmen de jiā. & Nous aimons tous notre famille. \\
\hline
\end{tabularx}
\end{center}
\end{exemplebox}

\courssub{Les mots de liaison du texte : 和 et 都}

Soulignons maintenant dans le texte les deux \cle{connecteurs} qui tiennent les phrases ensemble.

\textbf{1. 和 (hé), « et » (coordination nominale).} Dans \zh{爸爸、妈妈、哥哥、妹妹和我}, le mot \zh{和} relie le dernier élément de l'énumération, comme le « et » français avant le dernier mot d'une liste. Les éléments précédents sont séparés par la virgule chinoise d'énumération \zh{、} (dùnhào), qui ne s'emploie qu'entre des mots de même nature (noms, verbes, adjectifs).

\begin{attention}[和 ne relie que des noms / pronoms, jamais des phrases complètes]
Erreur classique de francophone : traduire « et » par \zh{和} partout.\\
Faux : \zh{我爱我的家和我的家很大。} (\zh{和} relie deux phrases).\\
Juste : \zh{我爱我的家，我的家也很大。} (Wǒ ài wǒ de jiā, wǒ de jiā yě hěn dà.) « J'aime ma famille, et ma famille est aussi très grande. » (Utiliser \zh{也} ou simple virgule).\\
Juste : \zh{我和你去北京。} (Wǒ hé nǐ qù Běijīng.) « Toi et moi allons à Pékin. » (\zh{和} relie deux pronoms).
\end{attention}

\textbf{2. 都 (dōu), « tous » (quantificateur distributif).} Dans \zh{我们都爱我们的家}, l'adverbe \zh{都} résume tout le groupe présenté avant (les cinq personnes). Il donne à la conclusion sa force : personne n'est oublié, tout le monde aime la famille. C'est le mot qui transforme une simple fin de texte en vraie conclusion en écho.

\begin{savaistu}[Pourquoi compte-t-on les familles en « bouches » ?]
En Chine, on compte les membres d'une famille avec le classificateur \zh{口} (kǒu), qui signifie littéralement « bouche » : \zh{五口人} (wǔ kǒu rén), c'est « cinq bouches », c'est-à-dire cinq personnes à nourrir. Cette image très concrète vient de la Chine rurale traditionnelle, où la première préoccupation d'une famille était de remplir le bol de riz de chacun. Le caractère \zh{口} se dessine d'ailleurs comme une petite bouche carrée, en trois traits seulement (verticale, horizontale pliée, horizontale). Au Cameroun, on dirait plutôt « nous sommes cinq à la maison » ou « cinq têtes » : chaque langue a sa façon d'imaginer la famille !
\end{savaistu}

\courssub{Exercice résolu et application}

\begin{exoresolu}[Retrouver la structure d'un texte]
Énoncé. Voici le texte d'Amina, mélangé. Remets les phrases dans l'ordre du plan du modèle (ouverture → membres → professions → goûts → conclusion), puis traduis la phrase de conclusion.\\
(a) 我的妈妈是医生。 (b) 我们都爱我们的家。 (c) 我爱我的家。 (d) 我的爸爸很喜欢足球。 (e) 我家有三口人：爸爸、妈妈和我。
\tcblower
Solution détaillée.\\
Étape 1 — Je traduis chaque phrase : (a) « Ma mère est médecin » ; (b) « Nous aimons tous notre famille » ; (c) « J'aime ma famille » ; (d) « Mon père aime beaucoup le football » ; (e) « Ma famille compte trois personnes : papa, maman et moi ».\\
Étape 2 — J'applique le plan du modèle : ouverture = (c) ; membres = (e) ; professions = (a) ; goûts = (d) ; conclusion = (b).\\
Ordre correct : \textbf{c – e – a – d – b}.\\
Étape 3 — Traduction de la conclusion : \zh{我们都爱我们的家。} (Wǒmen dōu ài wǒmen de jiā.) « Nous aimons tous notre famille. » Je remarque que \zh{都} résume les trois personnes de la phrase (e) : c'est bien une conclusion en écho de l'ouverture (c).
\end{exoresolu}

\begin{exercice}[À ton tour — Analyse du modèle]
Reprends le texte modèle \zh{我爱我的家} et réponds par écrit, en français :
\begin{enumerate}
\item Quelle phrase joue le rôle d'ouverture ? Justifie.
\item Combien de membres compte la famille, et avec quel classificateur sont-ils comptés ?
\item Souligne dans le texte les deux connecteurs \zh{和} et \zh{都} et explique leur rôle respectif.
\item Quel type de phrase (structure \zh{是} / \zh{有} / adverbe) est utilisé dans \zh{我的妈妈是医生} ? Dans \zh{哥哥很喜欢足球} ?
\end{enumerate}
\end{exercice}

\begin{corrige}[Exercice « À ton tour »]
1) L'ouverture est \zh{我爱我的家。} : elle annonce le sujet (\zh{我}) et le sentiment principal (\zh{爱}) du texte. 2) La famille compte cinq personnes (\zh{五口人}), comptées avec le classificateur \zh{口} (bouche). 3) \zh{和} relie le dernier nom de l'énumération (\zh{妹妹和我}) ; \zh{都} (dans \zh{我们都爱}) quantifie le sujet pluriel \zh{我们} et renforce l'unité familiale en conclusion. 4) \zh{我的妈妈是医生} : phrase d'identité (Sujet + \zh{是} + Nom). \zh{哥哥很喜欢足球} : phrase descriptive avec adverbe de degré (Sujet + \zh{很} + Verbe + Objet).
\end{corrige}

\begin{aretenir}[L'essentiel : Structure du texte et types de phrases]
\begin{itemize}
\item Un texte simple sur la famille suit un plan en 5 étapes : \textbf{Ouverture (sentiment)} → \textbf{Membres (有 + 口)} → \textbf{Professions (是)} → \textbf{Goûts/Qualités (很 + V/Adj)} → \textbf{Conclusion en écho (都 + 爱)}.
\item Trois structures de base suffisent : \textbf{S + 是 + N} (identité), \textbf{S + 有 + Num + Classif + N} (possession), \textbf{S + 很/都 + V/Adj} (description/quantification).
\item \zh{和} = « et » entre noms seulement. \zh{都} = « tous » placé avant le verbe, sujet pluriel obligatoire.
\item Méthode de production : \textbf{Lire-Traduire} → \textbf{Repérer structure} → \textbf{Souligner connecteurs} → \textbf{Remplir gabarit}.
\end{itemize}
\end{aretenir}

\coursec{Thème : Français → Chinois}

\courssub{Consignes et méthode}

\begin{methode}[Réussir un thème (français → chinois)]
\begin{enumerate}
\item \textbf{Analyser} la phrase française : identifier le sujet, le verbe, le complément, le temps (ici présent), la nuance (identité, possession, description).
\item \textbf{Choisir la structure chinoise} adaptée (voir tableau ci-dessus). Attention : \textbf{pas de mot-à-mot} !
\item \textbf{Ordonner} : Sujet — Adverbe (很/都) — Verbe (是/有/Verbe action) — Complément (Nom / Adjectif / Objet). Penser au classificateur après le nombre.
\item \textbf{Vérifier} : particules (\zh{的} pour possession), ponctuation chinoise, pinyin/tonnes.
\end{enumerate}
\end{methode}

\begin{exercice}[Thème — Phrases à traduire en chinois (caractères + pinyin)]
Traduis les phrases suivantes en t'appuyant \emph{uniquement} sur le vocabulaire et les structures de la leçon.
\begin{enumerate}
\item J'aime ma famille.
\item Ma famille a quatre personnes : papa, maman, mon grand frère et moi.
\item Mon père est médecin.
\item Ma mère aime beaucoup le football.
\item Mon grand frère est très mignon.
\item Nous aimons tous notre maison.
\end{enumerate}
\end{exercice}

\begin{corrige}[Thème — Corrigé détaillé]
\begin{enumerate}
\item \zh{我爱我的家。} (Wǒ ài wǒ de jiā.) — Structure S + V + O. \zh{的} marque la possession.
\item \zh{我家有四口人：爸爸、妈妈、哥哥和我。} (Wǒ jiā yǒu sì kǒu rén : bàba, māma, gēge hé wǒ.) — \zh{有} + Nombre + \zh{口} + \zh{人}. Énumération avec \zh{、} et \zh{和}.
\item \zh{我的爸爸是医生。} (Wǒ de bàba shì yīshēng.) — Structure S + \zh{是} + Nom (profession). Pas de \zh{很}.
\item \zh{我的妈妈很喜欢足球。} (Wǒ de māma hěn xǐhuan zúqiú.) — S + \zh{很} + V + O. \zh{很} nécessaire devant verbe d'action pour intensifier.
\item \zh{我的哥哥很可爱。} (Wǒ de gēge hěn kě'ài.) — S + \zh{很} + Adj. \zh{很} obligatoire devant adjectif seul.
\item \zh{我们都爱我们的家。} (Wǒmen dōu ài wǒmen de jiā.) — S (pluriel) + \zh{都} + V + O. \zh{都} résume \zh{我们}.
\end{enumerate}
\end{corrige}

\coursec{Version : Chinois → Français}

\courssub{Consignes et méthode}

\begin{methode}[Réussir une version (chinois → français)]
\begin{enumerate}
\item \textbf{Segmenter} : repérer les mots (séparer le pinyin en mots, pas en syllabes).
\item \textbf{Identifier la structure} : quel type de phrase ? (\zh{是}, \zh{有}, \zh{很}, \zh{都}).
\item \textbf{Traduire mot à mot} d'abord, puis \textbf{reformuler en français naturel} (pas de « ma mère est médecin » → « ma mère est médecin » OK ; mais « ma famille a quatre bouches personnes » → « ma famille compte quatre personnes »).
\item \textbf{Soigner l'accord} (genre, nombre) et le vocabulaire précis (professions, liens de parenté).
\end{enumerate}
\end{methode}

\begin{exercice}[Version — Textes à traduire en français]
Traduis les phrases suivantes en français naturel et correct.
\begin{enumerate}
\item 我爱我的爸爸和妈妈。
\item 我家有六口人：爸爸、妈妈、哥哥、姐姐、妹妹和我。
\item 我的哥哥是老师，我的姐姐是医生。
\item 妹妹很喜欢足球，爸爸也很喜欢足球。
\item 我们都爱我们的家，家很温暖。
\end{enumerate}
\end{exercice}

\begin{corrige}[Version — Corrigé détaillé]
\begin{enumerate}
\item « J'aime mon papa et ma maman. » (Structure : S + V + O coordonnés par \zh{和}).
\item « Ma famille compte six personnes : papa, maman, grand frère, grande sœur, petite sœur et moi. » (\zh{六口人}, \zh{姐姐} jiějie = grande sœur).
\item « Mon grand frère est professeur, ma grande sœur est médecin. » (Deux phrases \zh{是} juxtaposées par virgule).
\item « Ma petite sœur aime beaucoup le football, papa aussi aime beaucoup le football. » (\zh{也} yě = aussi, placé avant \zh{很} ou le verbe).
\item « Nous aimons tous notre famille, la famille est très chaleureuse. » (\zh{温暖} wēnnuǎn = chaleureux/chaleureuse ; conclusion en écho + commentaire).
\end{enumerate}
\end{corrige}

\coursec{Grille d'évaluation et entraînement à l'examen}

\courssub{Grille d'auto-évaluation pour la production écrite (Thème / Rédaction)}

Utilise cette grille pour relire ton texte avant de le rendre, ou pour évaluer celui d'un camarade. Chaque critère est noté sur 2 points (0 = absent/incorrect, 1 = partiel/erreurs, 2 = maîtrisé).

\begin{definition}[Grille d'évaluation — Expression écrite (Famille / Amour)]
\begin{center}
\begin{tabularx}{\linewidth}{|X|X|X|}
\hline
\rowcolor{popblueL} Critère & Indicateurs de réussite (Niveau 2) & Points \\
\hline
\textbf{Vocabulaire thématique} & Utilise correctement au moins 8 mots du thème (membres famille, professions, adjectifs, \zh{爱}, \zh{家}). Pas de confusion \zh{爱}/\zh{受}, \zh{哥}/\zh{歌}, etc. & /2 \\
\hline
\textbf{Structures grammaticales} & Maîtrise des 3 structures : \zh{是} (identité), \zh{有}+\zh{口} (nombre), \zh{很}/\zh{都} (description). Pas de \zh{是} + Adj. & /2 \\
\hline
\textbf{Ordre des mots / Syntaxe} & Sujet - Adverbe - Verbe - Complément respecté. Place de \zh{的} (possessif) correcte. Place de \zh{和} (noms seulement). & /2 \\
\hline
\textbf{Écriture des caractères} & Caractères lisibles, traits dans l'ordre, proportions équilibrées (caractères à structure haut/bas, gauche/droite). Radicaux visibles. & /2 \\
\hline
\textbf{Organisation du texte} & Plan en 5 étapes respecté (Ouverture → Membres → Professions → Goûts → Conclusion écho). Connecteurs logiques. & /2 \\
\hline
\textbf{Pinyin \& Tons} & Pinyin noté correctement sur chaque phrase (ou texte lu à l'oral avec tons justes). & /2 \\
\hline
\multicolumn{2}{|r|}{\textbf{Total}} & \textbf{/12} \\
\hline
\end{tabularx}
\end{center}
\end{definition}

\courssub{Entraînement type examen : Sujet complet}

\begin{exercice}[Sujet d'entraînement — Expression écrite (Durée : 30 min)]
\textbf{Consigne :} Tu t'appelles Paul (ou Marie), élève à Douala. Tu écris un message WeChat à ton ami chinois Li Ming pour lui présenter ta famille. Écris un texte de \textbf{5 à 6 phrases} en caractères chinois (avec pinyin au-dessus ou à côté).

\textbf{Contraintes obligatoires :}
\begin{itemize}
\item Utiliser le plan en 5 étapes (Ouverture, Membres, Professions, Goûts, Conclusion).
\item Inclure \textbf{au moins} : \zh{爱}, \zh{家}, \zh{有}, \zh{口}, \zh{是}, \zh{很}, \zh{和}, \zh{都}, \zh{的}.
\item Présenter 4 à 6 membres (parents + fratrie).
\item Donner la profession des parents (enseignant, infirmier, commerçant, agriculteur... vocabulaire connu ou \zh{工人} gōngrén ouvrier, \zh{农民} nóngmín paysan).
\item Donner un goût ou une qualité pour au moins deux enfants.
\item Écrire les caractères soigneusement (ordre des traits).
\end{itemize}

\textbf{Aide au démarrage (Plan à trous) :}
\begin{enumerate}
\item \zh{我爱我的家。} (Ouverture)
\item \zh{我家有\_\_口人：\_\_、\_\_、\_\_和我。} (Membres)
\item \zh{我的爸爸是\_\_，我的妈妈是\_\_。} (Professions)
\item \zh{\_\_很\_\_，\_\_很\_\_。} (Goûts / Qualités)
\item \zh{我们都爱我们的家。} (Conclusion)
\end{enumerate}
\end{exercice}

\begin{corrige}[Sujet d'entraînement — Proposition de rédaction type]
\textbf{Texte de Paul (5 personnes) :}
\begin{textebox}[Ma famille — Rédaction élève]
我爱我的家。\\
Wǒ ài wǒ de jiā.\\
\zh{我家有五口人：爸爸、妈妈、哥哥、妹妹和我。}\\
Wǒ jiā yǒu wǔ kǒu rén : bàba, māma, gēge, mèimei hé wǒ.\\
\zh{我的爸爸是老师，我的妈妈是护士。}\\
Wǒ de bàba shì lǎoshī, wǒ de māma shì hùshi.\\
\zh{哥哥很喜欢足球，妹妹很可爱。}\\
Gēge hěn xǐhuan zúqiú, mèimei hěn kě'ài.\\
\zh{我们都爱我们的家。}\\
Wǒmen dōu ài wǒmen de jiā.
\end{textebox}

\textbf{Auto-évaluation avec la grille :}
\begin{itemize}
\item Vocabulaire : 10 mots thème utilisés (✓) → 2/2.
\item Structures : \zh{是} (x2), \zh{有}+\zh{口}, \zh{很}+V, \zh{很}+Adj, \zh{都}+V (✓) → 2/2.
\item Syntaxe : Ordre parfait, \zh{的} correct, \zh{和} entre noms (✓) → 2/2.
\item Caractères : À vérifier par l'enseignant (visuel).
\item Organisation : 5 étapes respectées, écho parfait (✓) → 2/2.
\item Pinyin : Fourni et tonné (✓) → 2/2.
\item \textbf{Total estimé : 12/12.}
\end{itemize}
\end{corrige}

\begin{aretenir}[Check-list finale avant l'examen]
\begin{itemize}
\item[☐] Je sais écrire \textbf{de mémoire} et dans l'ordre des traits : \zh{爱, 家, 爸, 妈, 哥, 妹, 口, 人, 是, 有, 很, 都, 和, 的}.
\item[☐] Je ne confonds plus : \zh{爱}/\zh{受}, \zh{好}/\zh{妹}, \zh{姐}/\zh{妹}, \zh{爸}/\zh{父}, \zh{妈}/\zh{马}, \zh{哥}/\zh{歌}.
\item[☐] Je construis les 3 phrases-types sans hésiter : \zh{S+是+N}, \zh{S+有+Num+口+N}, \zh{S+很/都+V/Adj}.
\item[☐] Je place \zh{和} (noms) et \zh{都} (sujet pluriel + préverbal) correctement.
\item[☐] Je produis un texte « Famille » en 5 étapes (Ouverture → Membres → Métiers → Goûts → Conclusion écho) en 15 minutes.
\end{itemize}
\end{aretenir}

\coursec{Écriture des caractères : traits et radicaux}

Dans la section précédente, nous avons lu et commenté un modèle de texte sur la famille : \zh{我爱我的家。爸爸妈妈很爱我。} (\emph{Wǒ ài wǒ de jiā. Bàba māma hěn ài wǒ.} « J'aime ma famille. Papa et maman m'aiment beaucoup. »). Tu sais maintenant \emph{lire} et \emph{comprendre} ces phrases. Mais à l'examen, on te demandera aussi de les \emph{écrire} à la main, sans faute. En chinois, écrire n'est pas seulement une question d'orthographe : c'est un geste précis, réglé par des lois d'ordre des traits et par la connaissance des radicaux. Cette section t'apprend à écrire correctement les six caractères clés du thème de l'amour et de la famille : \zh{爱} (ài, aimer), \zh{家} (jiā, famille), \zh{妈} (mā, maman), \zh{爸} (bà, papa), \zh{我} (wǒ, je) et \zh{的} (de).

\courssub{Observer avant d'écrire : qu'est-ce qu'un caractère ?}

Observe d'abord ces deux caractères sans chercher à les écrire :

\begin{exemplebox}[Observer la structure]
\zh{妈} (mā) = \zh{女} (à gauche) + \zh{马} (à droite)\\
\zh{家} (jiā) = \zh{宀} (en haut) + \zh{豕} (en bas)
\end{exemplebox}

Tu remarques qu'un caractère chinois n'est pas un dessin arbitraire : il est \cle{composé} de pièces plus petites, comme des briques de construction. Deux types de briques sont essentiels :

\begin{definition}[Le radical (部首 \emph{bùshǒu})]
Le \cle{radical} (部首 \emph{bùshǒu}) est le composant graphique qui \textbf{classe} le caractère dans le dictionnaire et qui donne souvent un \textbf{indice de sens}. Exemples : \zh{女} (femme) signale un lien avec le féminin (\zh{妈} maman, \zh{姐} grande sœur) ; \zh{宀} (toit) signale un lien avec la maison (\zh{家} famille, \zh{安} paix). Le radical est à la fois une clé de classement et une clé de signification.
\end{definition}

L'autre brique fréquente est la \cle{partie phonétique}, qui donne un indice de \emph{prononciation}. Dans \zh{妈} (mā), la partie droite \zh{马} (mǎ, cheval) indique le son « ma » (seul le ton change) ; dans \zh{爸} (bà), la partie \zh{巴} (bā) indique le son « ba ». C'est une grande différence avec le français : en français, les lettres notent des sons ; en chinois, un caractère note une syllabe entière et combine souvent un indice de \emph{sens} (radical) et un indice de \emph{son} (phonétique).

\begin{attention}[La phonétique n'est qu'un indice]
La partie phonétique ne garantit pas toujours la prononciation exacte : elle donne souvent le son \emph{approché}, parfois avec un ton différent (\zh{马} mǎ $\rightarrow$ \zh{妈} mā), parfois avec une initiale voisine (\zh{巴} bā $\rightarrow$ \zh{爸} bà). Il faut donc toujours vérifier le pinyin : la phonétique aide à mémoriser, elle ne dispense pas d'apprendre.
\end{attention}

\courssub{Les règles d'ordre des traits}

Un caractère s'écrit trait par trait, dans un ordre \textbf{fixe} que tout écolier chinois apprend dès le primaire. Respecter cet ordre n'est pas un détail : il garantit une écriture rapide, équilibrée et lisible, et il est parfois vérifié à l'examen.

\begin{propriete}[Règles d'ordre des traits]
\begin{enumerate}
\item \textbf{De haut en bas} : on écrit d'abord les traits du haut, puis ceux du bas. Ex. \zh{三} (sān, trois) : la barre du haut, puis celle du milieu, puis celle du bas.
\item \textbf{De gauche à droite} : on écrit d'abord la partie gauche, puis la partie droite. Ex. \zh{妈} : \zh{女} d'abord, \zh{马} ensuite.
\item \textbf{L'horizontal (一) avant le vertical (丨)} quand ils se croisent. Ex. \zh{十} (shí, dix) : la barre horizontale d'abord, la verticale ensuite.
\item \textbf{L'extérieur avant l'intérieur, puis on ferme} : pour les caractères à enclos, on trace le pourtour, on remplit l'intérieur, et le trait de fermeture vient en dernier. Ex. \zh{国} (guó, pays) : le dernier trait est la barre horizontale du bas qui « ferme la porte ».
\item \textbf{Le point ou le petit trait final vient en dernier}. Ex. \zh{我} : le point en haut à droite est le 7\textsuperscript{e} et dernier trait.
\end{enumerate}
\end{propriete}

\begin{attention}[Erreurs fréquentes des francophones]
\begin{itemize}
\item Écrire le vertical avant l'horizontal dans \zh{十} ou \zh{干} : habitude à corriger dès maintenant.
\item « Fermer » un enclos trop tôt, puis se rendre compte qu'on a oublié l'intérieur : on ne peut plus rien ajouter proprement.
\item Écrire les caractères comme on dessine, dans n'importe quel ordre : le résultat est lent et déséquilibré. L'ordre des traits est une \emph{technique d'écriture}, pas une contrainte gratuite.
\end{itemize}
\end{attention}

\courssub{Étude détaillée des six caractères du thème}

\textbf{1. 爱 (ài, aimer) — 10 traits.}\par
Composants : \zh{爫} (griffe, en haut) + \zh{冖} (couvercle) + \zh{友} (ami, en bas). Ordre des traits : les quatre petits traits de \zh{爫} (1 à 4 : un 撇, deux points, un 撇, de gauche à droite), le point et le crochet horizontal de \zh{冖} (5--6), puis \zh{友} : horizontal (7), oblique tombant à gauche 撇 (8), horizontal-oblique 横撇 (9), et le dernier oblique 捺 (10). Image mentale : un ami (\zh{友}) protégé sous un couvercle (\zh{冖}) par une main (\zh{爫}) : aimer, c'est protéger.

\textbf{2. 家 (jiā, famille, maison) — 10 traits.}\par
Radical \zh{宀} (toit) + \zh{豕} (cochon). Ordre : point du toit (1), point vertical gauche (2), crochet horizontal du toit (3), puis \zh{豕} : horizontal (4), oblique 撇 (5), crochet courbe 弯钩 (6), deux obliques à gauche (7--8), oblique à droite (9), dernier 捺 (10). Règle appliquée : de haut en bas — le toit d'abord, le cochon ensuite.

\textbf{3. 妈 (mā, maman) — 6 traits.}\par
Radical \zh{女} (femme, 3 traits : oblique-point 撇点, oblique 撇, horizontal 横) à gauche, puis \zh{马} (cheval, 3 traits) à droite, partie phonétique qui donne le son « ma ». Règle appliquée : de gauche à droite. Attention : à gauche d'un caractère, \zh{女} s'écrit plus étroit et son dernier trait horizontal ne dépasse pas vers la droite.

\textbf{4. 爸 (bà, papa) — 8 traits.}\par
\zh{父} (père, 4 traits) en haut + \zh{巴} (bā, phonétique, 4 traits) en bas. Règle appliquée : de haut en bas. Les deux obliques de \zh{父} s'écartent comme deux bras ouverts.

\textbf{5. 我 (wǒ, je) — 7 traits.}\par
Attention particulière à l'ordre : petit 撇 en haut (1), horizontal (2), vertical-crochet 竖钩 (3), montant 提 (4), puis le grand 斜钩 (5), long oblique avec crochet qui descend vers la droite et donne l'équilibre du caractère, ensuite l'oblique 撇 (6) qui descend vers la gauche, et enfin le petit point (7) en haut à droite. Ne jamais oublier ce point : sans lui, le caractère est faux.

\textbf{6. 的 (de) — 8 traits.}\par
\zh{白} (blanc, 5 traits) à gauche + \zh{勺} (cuillère, 3 traits) à droite. C'est le mot-outil possessif et descriptif le plus fréquent du chinois : \zh{我的家} (\emph{wǒ de jiā}, « ma famille »).

\begin{center}
\begin{tabularx}{\linewidth}{|X|X|X|X|}
\hline
\rowcolor{popblueL} Caractère & Pinyin & Traits / Composition & Traduction \\
\hline
爱 & ài & 10 traits : 爫 + 冖 + 友 & aimer \\
\hline
家 & jiā & 10 traits : 宀 + 豕 & famille, maison \\
\hline
妈 & mā & 6 traits : 女 + 马 & maman \\
\hline
爸 & bà & 8 traits : 父 + 巴 & papa \\
\hline
我 & wǒ & 7 traits (point final !) & je, moi \\
\hline
的 & de & 8 traits : 白 + 勺 & marque possessive \\
\hline
\end{tabularx}
\end{center}

\begin{popfigure}
\begin{tikzpicture}[scale=0.85]
% Grille tianzige 爱
\draw[popline, thick] (0,0) rectangle (3,3);
\draw[popline, dashed] (0,1.5)--(3,1.5);
\draw[popline, dashed] (1.5,0)--(1.5,3);
\draw[popline, dashed] (0,0)--(3,3);
\draw[popline, dashed] (0,3)--(3,0);
\node[popblue] at (1.5,1.5) {\Huge 爱};
\node[below, popdark] at (1.5,-0.1) {\small 爱 ài (10 traits)};
% Grille tianzige 家
\draw[popline, thick] (4.5,0) rectangle (7.5,3);
\draw[popline, dashed] (4.5,1.5)--(7.5,1.5);
\draw[popline, dashed] (6,0)--(6,3);
\draw[popline, dashed] (4.5,0)--(7.5,3);
\draw[popline, dashed] (4.5,3)--(7.5,0);
\node[poporange] at (6,1.5) {\Huge 家};
\node[below, popdark] at (6,-0.1) {\small 家 jiā (10 traits)};
% Étapes numérotées
\node[align=left, popdark] at (1.5,-1.6) {\small 1--4 : 爫 \; 5--6 : 冖 \; 7--10 : 友};
\node[align=left, popdark] at (6,-1.6) {\small 1--3 : 宀 \; 4--10 : 豕};
\end{tikzpicture}
\legende{La grille 田字格 (\emph{tiánzìgé}) : les pointillés aident à centrer chaque composant. On écrit de haut en bas et de gauche à droite.}
\end{popfigure}

\begin{methode}[Mémoriser un caractère en cinq étapes]
\begin{enumerate}
\item \textbf{Observer} le caractère entier, sa forme générale (carrée, haute, large).
\item \textbf{Décomposer} : identifier le radical (sens) et la partie phonétique (son). Ex. \zh{妈} = 女 (femme) + 马 (son « ma »).
\item \textbf{Tracer en l'air} avec le doigt en comptant les traits à voix haute.
\item \textbf{Écrire 5 fois} dans des cases 田字格, en respectant l'ordre des traits.
\item \textbf{Réutiliser} le caractère dans une phrase : \zh{妈妈爱我。}
\end{enumerate}
\end{methode}

\begin{exemplebox}[Les caractères en action]
\zh{妈妈爱我。} \emph{Māma ài wǒ.} « Maman m'aime. »\\
\zh{我爱我的家。} \emph{Wǒ ài wǒ de jiā.} « J'aime ma famille. »\\
\zh{爸爸的家在雅温得。} \emph{Bàba de jiā zài Yǎwēndé.} « La maison de papa est à Yaoundé. »
\end{exemplebox}

\begin{savaistu}[Un cochon sous le toit]
Le caractère \zh{家} (jiā) représente un cochon (\zh{豕}) sous un toit (\zh{宀}). Dans la Chine ancienne, posséder un cochon sous son toit était le signe d'un foyer prospère : le cochon fournissait nourriture et richesse. L'idée est proche de certaines traditions camerounaises où le bétail signale la richesse d'une famille. De même, \zh{安} (ān, paix) montre une femme (\zh{女}) sous un toit : la paix, c'est la famille à l'abri. Remarque aussi que le caractère traditionnel de « aimer » s'écrit 愛 : il contient le cœur \zh{心}, que la simplification a retiré — les Chinois plaisantent parfois en disant que « l'amour moderne n'a plus de cœur ».
\end{savaistu}

\begin{exoresolu}[Compter les traits et trouver le radical]
Énoncé : pour chaque caractère, indique le nombre de traits, le radical et la partie phonétique éventuelle : \zh{妈}, \zh{爸}, \zh{家}.
\tcblower
Solution détaillée :
\begin{itemize}
\item \zh{妈} (mā) : 6 traits. Radical : \zh{女} (femme) — indice de sens. Phonétique : \zh{马} (mǎ) — indice de son « ma ».
\item \zh{爸} (bà) : 8 traits. Radical : \zh{父} (père). Phonétique : \zh{巴} (bā) — indice de son « ba ».
\item \zh{家} (jiā) : 10 traits. Radical : \zh{宀} (toit). Le bas \zh{豕} (cochon) est le composant de sens historique.
\end{itemize}
Méthode : toujours chercher d'abord le composant qui donne le sens (souvent à gauche ou en haut), puis vérifier si l'autre partie ressemble à un caractère connu dont le son est proche.
\end{exoresolu}

\begin{exercice}[Entraînement d'écriture]
Dans ton cahier, trace des cases 田字格 de 2 cm de côté. Écris chaque caractère \textbf{5 fois} en comptant les traits à voix haute : \zh{爱}, \zh{家}, \zh{妈}, \zh{爸}, \zh{我}, \zh{的}. Puis recopie trois fois la phrase \zh{我爱我的家。} en respectant l'ordre des traits.
\end{exercice}

\begin{corrige}[Entraînement d'écriture]
Vérification en classe : l'enseignant contrôle (1) l'ordre des traits — horizontal avant vertical, haut avant bas, gauche avant droite ; (2) le point final de \zh{我} et le grand 斜钩 (5\textsuperscript{e} trait) qui doit descendre bien à droite ; (3) l'équilibre du caractère dans la case : le toit \zh{宀} de \zh{家} doit couvrir le bas sans le toucher, et dans \zh{妈} la partie gauche \zh{女} est plus étroite que la partie droite \zh{马}.
\end{corrige}

\begin{aretenir}[L'essentiel de la section]
Un caractère = radical (sens) + souvent une phonétique (son). On écrit de haut en bas, de gauche à droite, l'horizontal avant le vertical, l'extérieur avant l'intérieur, et on ferme en dernier. Les six caractères du thème : \zh{爱} (10 traits), \zh{家} (10), \zh{妈} (6), \zh{爸} (8), \zh{我} (7, point final !), \zh{的} (8).
\end{aretenir}

\coursec{Caractères proches : ne pas confondre}

Dans la section précédente, nous avons appris à écrire les caractères du thème de l'amour et de la famille en respectant l'ordre des traits et en identifiant les radicaux. Mais connaître les traits ne suffit pas : le chinois possède de nombreux caractères qui se ressemblent beaucoup, parfois à un seul élément près. Une erreur de radical peut transformer une phrase d'amour en phrase illisible ! Cette section vous donne une stratégie simple et efficace : \cle{associer le radical au sens} pour choisir le bon caractère.

\courssub{Observer avant d'écrire : deux phrases témoins}

Lisez attentivement les deux paires de phrases suivantes. Une seule est correcte dans chaque paire.

\begin{exemplebox}[Deux phrases, un seul caractère de différence]
\zh{你爱我吗？} \emph{(Nǐ ài wǒ ma?)} — « Tu m'aimes ? »\\[4pt]
\zh{你受我吗？} \emph{(Nǐ shòu wǒ ma?)} — phrase \textbf{incorrecte} : \zh{受} signifie « recevoir, subir », pas « aimer ».\\[6pt]
\zh{她是我妹妹。} \emph{(Tā shì wǒ de mèimei.)} — « Elle est ma petite sœur. »\\[4pt]
\zh{他是我妹妹。} \emph{(Tā shì wǒ de mèimei.)} — phrase \textbf{incorrecte} : \zh{他} signifie « il », or une petite sœur est une personne féminine.
\end{exemplebox}

On observe que deux caractères presque identiques (\zh{爱}/\zh{受}, \zh{他}/\zh{她}) n'ont ni le même sens, ni parfois la même prononciation. La différence se cache dans \cle{le radical} ou dans un composant. Apprenons à les repérer.

\courssub{La règle du radical : le sens est écrit dans le caractère}

\begin{propriete}[Règle du radical indicateur de sens]
La plupart des caractères chinois sont composés de deux parties : un \cle{radical} (clé) qui indique la \textbf{famille de sens}, et une partie \textbf{phonétique} qui suggère la prononciation. Pour les caractères de la famille et des personnes :
\begin{itemize}
\item \zh{亻}(variante de \zh{人}, « personne ») et \zh{女} (« femme ») signalent une \textbf{personne} : \zh{他} (il), \zh{她} (elle), \zh{妈} (maman), \zh{妹} (petite sœur).
\item \zh{口} (« bouche ») signale souvent un lien avec la \textbf{bouche ou le son} : c'est le radical de nombreuses particules orales comme \zh{吗} (particule interrogative), \zh{呢} (ne), \zh{吧} (ba).
\item \zh{扌}(variante de \zh{手}, « main ») signale une \textbf{action faite avec la main} : \zh{找} (chercher), \zh{打} (frapper, téléphoner), \zh{推} (pousser).
\item \zh{宀} (« toit ») signale un lien avec la \textbf{maison, l'abri, les choses} : \zh{它} (il/elle pour les choses et les animaux), \zh{家} (famille, maison), \zh{安} (paix, calme).
\end{itemize}
\textbf{Stratégie :} quand vous hésitez entre deux caractères proches, demandez-vous : « De quoi parle ma phrase : d'une personne, d'une action de la main, d'un son de la bouche ? » Le radical vous donne la réponse.
\end{propriete}

\begin{savaistu}[Pourquoi 吗 a-t-il une bouche ?]
Les particules comme \zh{吗} (ma), \zh{呢} (ne), \zh{吧} (ba) sont des petits mots que l'on \emph{prononce} mais qui n'ont pas de sens concret propre : elles portent l'intonation ou la modalité de la phrase. Les Chinois ont donc naturellement choisi le radical de la bouche \zh{口} pour les écrire. Retenez : \textbf{particule = bouche}.
\end{savaistu}

\courssub{Les cinq paires à ne jamais confondre}

\textbf{1. 爱 (ài, aimer) vs 受 (shòu, recevoir)}\par
Les deux caractères partagent le haut \zh{⺤} (griffe/main) et le couvercle \zh{冖}. Mais en bas, \zh{爱} contient \zh{友} (« ami ») : aimer, c'est avoir un ami dans le cœur ! \zh{受} se termine par \zh{又} (main qui reçoit). Moyen mnémotechnique : \textbf{爱 = couvercle + ami ; 受 = couvercle + main qui reçoit}.

\begin{exemplebox}[爱 ou 受 ?]
\zh{我爱我的家。} \emph{(Wǒ ài wǒ de jiā.)} — « J'aime ma famille. »\\[4pt]
\zh{我收到了妈妈的信。} \emph{(Wǒ shōudào le māma de xìn.)} — « J'ai reçu la lettre de maman. »
\end{exemplebox}

\textbf{2. 妈 (mā, maman) vs 吗 (ma, particule interrogative)}\par
Même partie phonétique \zh{马} (mǎ, cheval), mais deux radicaux différents : \zh{女} (femme) pour la maman, \zh{口} (bouche) pour la particule que l'on prononce. Attention aussi aux tons : \zh{妈} est au premier ton (mā), \zh{吗} est au ton neutre (ma).

\begin{attention}[Erreur fréquente des francophones : 吗 au lieu de 妈]
Beaucoup d'élèves écrivent \zh{我爱吗} ou \zh{我吗很好} en voulant dire « J'aime maman » ou « Ma maman va bien ». Ces phrases sont \textbf{illisibles} pour un sinophone : \zh{吗} n'est qu'une particule de question, jamais un nom de personne. Le radical \zh{女} (femme) est le signal obligatoire de « maman » : \zh{妈}. De même, n'écrivez jamais \zh{你吗} pour « ta mère » : on écrit \zh{你妈妈} (nǐ māma) ou \zh{你妈} (nǐ mā).
\end{attention}

\textbf{3. 他 (tā, il) vs 她 (tā, elle) vs 它 (tā, il/elle pour les choses)}\par
Trois caractères, une seule prononciation : \emph{tā}. À l'oral, impossible de les distinguer ; à l'écrit, le radical décide : \zh{亻} (personne) pour l'homme, \zh{女} (femme) pour la femme, \zh{宀} (toit) pour les choses et les animaux. Le français fait la distinction à l'oral (il/elle), mais le chinois ne la fait \textbf{qu'à l'écrit} : c'est l'inverse du français !

\begin{exemplebox}[他, 她 ou 它 ?]
\zh{他是我哥哥。} \emph{(Tā shì wǒ gēge.)} — « Il est mon grand frère. »\\[4pt]
\zh{她是我妹妹。} \emph{(Tā shì wǒ de mèimei.)} — « Elle est ma petite sœur. »\\[4pt]
\zh{这是我的狗，它很可爱。} \emph{(Zhè shì wǒ de gǒu, tā hěn kě'ài.)} — « Voici mon chien, il est très mignon. »
\end{exemplebox}

\textbf{4. 我 (wǒ, je) vs 找 (zhǎo, chercher)}\par
\zh{找} contient \zh{我} à droite, avec en plus le radical de la main \zh{扌} à gauche : chercher, c'est « mettre la main sur moi ». Si votre phrase contient un sujet « je », écrivez \zh{我} seul ; si elle exprime l'action de chercher, ajoutez la main : \zh{找}.

\begin{exemplebox}[我 ou 找 ?]
\zh{我找我的书。} \emph{(Wǒ zhǎo wǒ de shū.)} — « Je cherche mon livre. » (les deux caractères dans la même phrase !)\\[4pt]
\zh{我爱我家。} \emph{(Wǒ ài wǒ de jiā.)} — « J'aime ma famille. »
\end{exemplebox}

\textbf{5. 的 (de) vs 白 (bái, blanc)}\par
\zh{的} = \zh{白} + \zh{勺} (sháo, cuillère). Le mot-outil \zh{的} (marque de possession et de qualification) est le caractère le plus fréquent du chinois : ne l'oubliez jamais dans « ma mère » \zh{我的妈妈}. \zh{白} seul signifie « blanc » : \zh{白色} (báisè, couleur blanche).

\courssub{Tableau récapitulatif des paires}

\begin{center}
\begin{tabularx}{\linewidth}{|X|X|X|X|}
\hline
\rowcolor{popblueL} Caractères & Pinyin & Élément distinctif & Traduction \\
\hline
\zh{爱} & ài & \zh{友} (ami) en bas & aimer \\
\hline
\zh{受} & shòu & \zh{又} (main) en bas & recevoir, subir \\
\hline
\zh{妈} & mā & radical \zh{女} (femme) & maman \\
\hline
\zh{吗} & ma & radical \zh{口} (bouche) & particule interrogative \\
\hline
\zh{他} & tā & radical \zh{亻} (personne) & il \\
\hline
\zh{她} & tā & radical \zh{女} (femme) & elle \\
\hline
\zh{它} & tā & radical \zh{宀} (toit) & il/elle (chose, animal) \\
\hline
\zh{我} & wǒ & caractère simple, 7 traits & je, me \\
\hline
\zh{找} & zhǎo & \zh{扌} (main) + \zh{我} & chercher \\
\hline
\zh{的} & de & \zh{白} + \zh{勺} & mot-outil (possession) \\
\hline
\zh{白} & bái & caractère simple, 5 traits & blanc \\
\hline
\end{tabularx}
\end{center}

\begin{popfigure}
\begin{tikzpicture}[node distance=0.35cm and 0.9cm,
  pair/.style={rectangle, rounded corners=2pt, draw=popline, fill=popcream, align=center, minimum width=2.6cm, minimum height=1.1cm, font=\large},
  diff/.style={font=\small\itshape, text=popdark, align=center},
  lien/.style={-{Stealth[length=2.5mm]}, thick, popmuted}]
\node[pair, align=center] (ai){\zh{\textcolor{popmuted}{⺤冖}\textcolor{popgreen}{友}}\\\zh{爱} ài};
\node[pair, right=of ai, align=center] (shou){\zh{\textcolor{popmuted}{⺤冖}\textcolor{poporange}{又}}\\\zh{受} shòu};
\node[diff, below=0.15cm of ai] {ami $\Rightarrow$ aimer};
\node[diff, below=0.15cm of shou] {main $\Rightarrow$ recevoir};
\node[pair, below=1.5cm of ai, align=center] (ma1){\zh{\textcolor{popgreen}{女}\textcolor{popmuted}{马}}\\\zh{妈} mā};
\node[pair, right=of ma1, align=center] (ma2){\zh{\textcolor{poporange}{口}\textcolor{popmuted}{马}}\\\zh{吗} ma};
\node[diff, below=0.15cm of ma1] {femme $\Rightarrow$ maman};
\node[diff, below=0.15cm of ma2] {bouche $\Rightarrow$ particule};
\node[pair, below=1.5cm of ma1, align=center] (ta1){\zh{\textcolor{popgreen}{亻}\textcolor{popmuted}{也}}\\\zh{他} tā};
\node[pair, right=of ta1, align=center] (ta2){\zh{\textcolor{poporange}{女}\textcolor{popmuted}{也}}\\\zh{她} tā};
\node[pair, right=of ta2, align=center] (ta3){\zh{\textcolor{poppurple}{宀}\textcolor{popmuted}{匕}}\\\zh{它} tā};
\node[diff, below=0.15cm of ta1] {personne $\Rightarrow$ il};
\node[diff, below=0.15cm of ta2] {femme $\Rightarrow$ elle};
\node[diff, below=0.15cm of ta3] {toit $\Rightarrow$ chose};
\draw[lien] (ai) -- (shou);
\draw[lien] (ma1) -- (ma2);
\draw[lien] (ta1) -- (ta2);
\draw[lien] (ta2) -- (ta3);
\end{tikzpicture}
\legende{Paires de caractères proches : le composant qui change tout est coloré (vert = sens de personne/féminin, orange = autre radical, violet = toit).}
\end{popfigure}

\courssub{Méthode : vérifier son écriture en trois étapes}

\begin{methode}[Le contrôle « traits + radical + sens »]
Après avoir écrit une phrase en chinois, relisez-la ainsi :
\begin{enumerate}
\item \textbf{Compter les traits} des caractères-clés du thème (\zh{爱} = 10 traits, \zh{我} = 7 traits, \zh{妈} = 6 traits). Un trait en plus ou en moins signale une erreur.
\item \textbf{Vérifier le radical} de chaque caractère de personne ou d'action : une femme exige \zh{女} (\zh{妈}, \zh{她}, \zh{妹}) ; une action de la main exige \zh{扌} (\zh{找}) ; une particule de question exige \zh{口} (\zh{吗}).
\item \textbf{Traduire mentalement} chaque caractère : si un caractère ne correspond pas au sens voulu (« recevoir » au lieu de « aimer »), c'est que vous avez choisi le caractère proche, pas le bon.
\end{enumerate}
\end{methode}

\begin{exoresolu}[Choisir le bon caractère]
Énoncé — Complétez chaque phrase avec le bon caractère parmi les deux proposés, puis traduisez.
\begin{enumerate}
\item 我（爱 / 受）我的妈妈。
\item 你（妈 / 吗）妈是老师。
\item （他 / 她）是我的妹妹。
\item 我（我 / 找）我的书。
\item 你爱我（妈 / 吗）？
\end{enumerate}
\tcblower
Solution détaillée :
\begin{enumerate}
\item \zh{我爱我的妈妈。} \emph{(Wǒ ài wǒ de māma.)} — « J'aime ma maman. » On choisit \zh{爱} (avec \zh{友}) car la phrase exprime un sentiment, pas une réception.
\item \zh{你妈妈是老师。} \emph{(Nǐ māma shì lǎoshī.)} — « Ta maman est professeur. » On choisit \zh{妈} (radical \zh{女}) car on parle d'une personne féminine ; \zh{吗} ne peut pas être un nom. On écrit \zh{妈妈} (deux syllabes) à l'écrit standard.
\item \zh{她是我的妹妹。} \emph{(Tā shì wǒ de mèimei.)} — « Elle est ma petite sœur. » Une petite sœur est féminine : radical \zh{女}, donc \zh{她}.
\item \zh{我找我的书。} \emph{(Wǒ zhǎo wǒ de shū.)} — « Je cherche mon livre. » L'action de chercher exige la main \zh{扌} : \zh{找}. (Remarquez le \zh{我} du début : sujet « je », sans main.)
\item \zh{你爱我吗？} \emph{(Nǐ ài wǒ ma?)} — « Tu m'aimes ? » Fin de question : particule de la bouche \zh{吗}. La phrase est déclarative jusqu'à la fin, donc ce n'est pas « maman ».
\end{enumerate}
\end{exoresolu}

\begin{exercice}[À vous de jouer]
Complétez avec le bon caractère, puis traduisez chaque phrase en français.
\begin{enumerate}
\item 爸爸很（爱 / 受）我。
\item 我（妈 / 吗）妈在雅温得工作。
\item 这是我的猫，（他 / 它）很小。
\item 你（找 / 我）谁？
\item 我（的 / 白）家是喀麦隆。
\end{enumerate}
\end{exercice}

\begin{corrige}[Exercice « À vous de jouer »]
\begin{enumerate}
\item \zh{爸爸很爱我。} \emph{(Bàba hěn ài wǒ.)} — « Papa m'aime beaucoup. » (sentiment $\Rightarrow$ \zh{爱})
\item \zh{我妈妈在雅温得工作。} \emph{(Wǒ māma zài Yǎwēndé gōngzuò.)} — « Ma maman travaille à Yaoundé. » (personne féminine $\Rightarrow$ \zh{妈})
\item \zh{这是我的猫，它很小。} \emph{(Zhè shì wǒ de māo, tā hěn xiǎo.)} — « Voici mon chat, il est très petit. » (animal $\Rightarrow$ \zh{它})
\item \zh{你找谁？} \emph{(Nǐ zhǎo shéi?)} — « Qui cherches-tu ? » (action de la main $\Rightarrow$ \zh{找})
\item \zh{我的家是喀麦隆。} \emph{(Wǒ de jiā shì Kāmàilóng.)} — « Mon pays (ma patrie), c'est le Cameroun. » (possession $\Rightarrow$ \zh{的})
\end{enumerate}
\end{corrige}

\begin{aretenir}[L'essentiel de la section]
\begin{itemize}
\item Deux caractères proches se distinguent par leur \cle{radical} ou un composant : \zh{爱}/\zh{受}, \zh{妈}/\zh{吗}, \zh{他}/\zh{她}/\zh{它}, \zh{我}/\zh{找}, \zh{的}/\zh{白}.
\item Le radical indique le sens : \zh{女} = femme, \zh{亻} = personne, \zh{口} = bouche/particule, \zh{扌} = action de la main, \zh{宀} = toit/chose.
\item \zh{他}, \zh{她}, \zh{它} se prononcent tous \emph{tā} : seule l'écriture les distingue.
\item Avant de rendre votre copie : comptez les traits, vérifiez le radical, traduisez mentalement chaque caractère.
\end{itemize}
\end{aretenir}

\coursec{Structure et connecteurs du texte}

Dans la section précédente, nous avons appris à distinguer les caractères proches du thème de la famille (\zh{我} et \zh{找}, \zh{妈} et \zh{马}, \zh{爱} et \zh{受}…). Savoir écrire les caractères ne suffit pas : il faut maintenant les assembler en phrases, puis relier les phrases entre elles pour obtenir un texte fluide et cohérent. C'est exactement l'objet de cette section : apprendre le \cle{plan type} d'une présentation de la famille et maîtriser les \cle{connecteurs} qui donnent du rythme et de la logique à votre production écrite.

\courssub{Observer d'abord : un texte modèle}

Lisons attentivement le texte suivant. Repérez comment les phrases s'enchaînent et soulignez mentalement les petits mots qui relient les idées.

\begin{textebox}[Ma famille — 我的家]
我爱我的家。我家有五口人：爸爸、妈妈、哥哥、妹妹和我。\\
Wǒ ài wǒ de jiā. Wǒ jiā yǒu wǔ kǒu rén : bàba, māma, gēge, mèimei hé wǒ.\\
« J'aime ma famille. Ma famille compte cinq personnes : papa, maman, mon grand frère, ma petite sœur et moi. »\\[4pt]
爸爸是老师，妈妈是医生。爸爸和妈妈都很忙，但是他们很爱我们。\\
Bàba shì lǎoshī, māma shì yīshēng. Bàba hé māma dōu hěn máng, dànshì tāmen hěn ài wǒmen.\\
« Papa est professeur, maman est médecin. Papa et maman sont très occupés, mais ils nous aiment beaucoup. »\\[4pt]
哥哥喜欢踢足球，我也喜欢足球。妹妹喜欢唱歌和画画。\\
Gēge xǐhuan tī zúqiú, wǒ yě xǐhuan zúqiú. Mèimei xǐhuan chànggē hé huàhuà.\\
« Mon grand frère aime jouer au football, moi aussi j'aime le football. Ma petite sœur aime chanter et dessiner. »\\[4pt]
我的家不大，但是很温暖。因为我们都爱我们的家，所以我们每天都很快乐。\\
Wǒ de jiā bù dà, dànshì hěn wēnnuǎn. Yīnwei wǒmen dōu ài wǒmen de jiā, suǒyǐ wǒmen měi tiān dōu hěn kuàilè.\\
« Ma maison n'est pas grande, mais elle est très chaleureuse. Parce que nous aimons tous notre famille, nous sommes heureux chaque jour. »
\end{textebox}

\textbf{Questions de compréhension}
\begin{enumerate}
\item Combien de personnes compte cette famille ? Citez les membres.
\item Quelles sont les professions du père et de la mère ?
\item Quels connecteurs l'auteur utilise-t-il pour passer d'une idée à l'autre ?
\end{enumerate}

\textbf{Glossaire}

\begin{center}
\begin{tabularx}{\linewidth}{|X|X|X|X|}
\hline
\rowcolor{popblueL} Caractères & Pinyin & Nature & Traduction \\
\hline
温暖 & wēnnuǎn & adjectif & chaleureux, doux (lieu, atmosphère) \\
\hline
温柔 & wēnróu & adjectif & doux, tendre (personne) \\
\hline
快乐 & kuàilè & adjectif & heureux, joyeux \\
\hline
踢足球 & tī zúqiú & verbe + nom & jouer au football \\
\hline
唱歌 & chànggē & verbe & chanter \\
\hline
画画 & huàhuà & verbe & dessiner, peindre \\
\hline
忙 & máng & adjectif & occupé \\
\hline
\end{tabularx}
\end{center}

\courssub{Le plan type d'un texte sur la famille}

Un bon texte de niveau Première n'est pas une suite de phrases au hasard : il suit un plan en \textbf{cinq étapes} que l'examinateur attend de retrouver. Ce plan structure la pensée et garantit la cohérence.

\begin{popfigure}
\begin{tikzpicture}[node distance=0.55cm,
  etape/.style={rectangle, rounded corners=3pt, draw=popblue, fill=popblueL, thick, text width=4.8cm, align=center, font=\small, minimum height=1.1cm},
  fleche/.style={-{Stealth[length=2.5mm]}, thick, poporange}]
\node[etape, align=center] (e1){\textbf{1. Ouverture}\\ \zh{我爱我的家。} Wǒ ài wǒ de jiā.};
\node[etape, below=of e1, align=center] (e2){\textbf{2. Effectif}\\ \zh{我家有五口人。} Wǒ jiā yǒu wǔ kǒu rén.};
\node[etape, below=of e2, align=center] (e3){\textbf{3. Membres et professions}\\ \zh{爸爸是老师，妈妈是医生。}};
\node[etape, below=of e3, align=center] (e4){\textbf{4. Qualités et goûts}\\ \zh{妈妈很温柔；哥哥喜欢足球。}};
\node[etape, below=of e4, align=center] (e5){\textbf{5. Conclusion}\\ \zh{我们都爱我们的家。}};
\draw[fleche] (e1) -- (e2);
\draw[fleche] (e2) -- (e3);
\draw[fleche] (e3) -- (e4);
\draw[fleche] (e4) -- (e5);
\end{tikzpicture}
\legende{Organigramme : le plan du texte sur la famille en cinq étapes fléchées}
\end{popfigure}

\begin{aretenir}[Le squelette à mémoriser]
\begin{enumerate}
\item \textbf{Ouverture} : déclarer son amour — \zh{我爱我的家。} (Wǒ ài wǒ de jiā.) « J'aime ma famille. »
\item \textbf{Effectif} : donner le nombre de personnes avec le classificateur \zh{口} — \zh{我家有……口人。} (Wǒ jiā yǒu … kǒu rén.) « Ma famille compte … personnes. »
\item \textbf{Membres et professions} : présenter chacun avec le verbe \zh{是} — \zh{爸爸是老师。} « Papa est professeur. »
\item \textbf{Qualités et goûts} : décrire avec \zh{很} + adjectif et \zh{喜欢} + verbe/nom — \zh{妈妈很温柔。哥哥喜欢踢足球。}
\item \textbf{Conclusion} : revenir au sentiment collectif — \zh{我们都爱我们的家。} (Wǒmen dōu ài wǒmen de jiā.) « Nous aimons tous notre famille. »
\end{enumerate}
\end{aretenir}

\courssub{Le connecteur 和 (hé) : « et » entre noms}

Observons : \zh{爸爸和妈妈都很忙。} (Bàba hé māma dōu hěn máng.) « Papa et maman sont très occupés. » Ici, \zh{和} relie deux \textbf{noms} : 爸爸 et 妈妈.

\begin{propriete}[Règle d'or de 和]
\zh{和} relie des \cle{mots} ou des \cle{groupes nominaux}, jamais deux phrases entières (deux propositions verbales).\\
Structure : \textbf{Nom 1 + 和 + Nom 2}.\\
Pour relier deux \emph{phrases} (deux idées avec chacune son verbe), on utilise \zh{也}, \zh{但是} ou \zh{因为……所以……}, jamais \zh{和}.
\end{propriete}

\begin{exemplebox}[和 entre noms]
\zh{我喜欢唱歌和画画。} (Wǒ xǐhuan chànggē hé huàhuà.) « J'aime chanter et dessiner. » (Coordination de deux verbes/activités).\\
\zh{哥哥和妹妹都是学生。} (Gēge hé mèimei dōu shì xuésheng.) « Mon frère et ma sœur sont tous les deux élèves. »\\
\zh{我家有爸爸、妈妈和我。} (Wǒ jiā yǒu bàba, māma hé wǒ.) « Dans ma famille, il y a papa, maman et moi. » (Dans une énumération de plus de deux éléments, \zh{和} ne se met qu'avant le dernier élément ; les autres sont séparés par la virgule chinoise 、.)
\end{exemplebox}

\begin{attention}[L'erreur n° 1 des francophones]
En français, « et » relie tout : des noms \emph{et} des phrases. En chinois, c'est impossible avec \zh{和}.\\
Incorrect : *\zh{我爱爸爸，和我爱妈妈。}\\
Correct : \zh{我爱爸爸，也爱妈妈。} (Wǒ ài bàba, yě ài māma.) « J'aime papa, et j'aime aussi maman. »\\
Retenez : entre deux phrases, « et » devient presque toujours \zh{也} + répétition ou reprise du verbe.
\end{attention}

\courssub{L'adverbe 也 (yě) : « aussi »}

Observons : \zh{哥哥喜欢踢足球，我也喜欢足球。} Remarquez la place de \zh{也} : \textbf{avant le verbe}, juste après le sujet.

\begin{propriete}[Place de 也]
Structure : \textbf{Sujet + 也 + Verbe}.\\
\zh{也} signifie « aussi » et sert à ajouter une information identique ou parallèle. Il ne se place jamais en fin de phrase comme le « aussi » français.
\end{propriete}

\begin{exemplebox}[也 en action]
\zh{我也爱我的家。} (Wǒ yě ài wǒ de jiā.) « Moi aussi, j'aime ma famille. »\\
\zh{妈妈是医生，爸爸也是医生。} (Māma shì yīshēng, bàba yě shì yīshēng.) « Maman est médecin, papa est médecin aussi. »\\
\zh{妹妹不喜欢足球，我也不喜欢。} (Mèimei bù xǐhuan zúqiú, wǒ yě bù xǐhuan.) « Ma sœur n'aime pas le football, moi non plus. » (也 fonctionne aussi dans les phrases négatives : « non plus ».)
\end{exemplebox}

\courssub{Le connecteur 但是 (dànshì) : « mais »}

Observons : \zh{我的家不大，但是很温暖。} (Wǒ de jiā bù dà, dànshì hěn wēnnuǎn.) « Ma maison n'est pas grande, mais elle est très chaleureuse. »

\begin{propriete}[Emploi de 但是]
Structure : \textbf{Phrase 1，但是 + Phrase 2}.\\
\zh{但是} introduit une opposition ou une nuance, exactement comme « mais » en français. Il se place après la virgule chinoise ，. Sa forme courte \zh{但} (dàn) est plus littéraire ; à l'écrit scolaire, préférez \zh{但是}.
\end{propriete}

\begin{exemplebox}[但是 : l'opposition]
\zh{爸爸很忙，但是他每天陪我学习。} (Bàba hěn máng, dànshì tā měi tiān péi wǒ xuéxí.) « Papa est très occupé, mais il m'aide à étudier chaque jour. »\\
\zh{我家很小，但是很漂亮。} (Wǒ jiā hěn xiǎo, dànshì hěn piàoliang.) « Notre maison est petite, mais très jolie. »\\
\zh{中文很难，但是很有意思。} (Zhōngwén hěn nán, dànshì hěn yǒu yìsi.) « Le chinois est difficile, mais très intéressant. »
\end{exemplebox}

\courssub{La corrélation 因为……所以…… : « parce que… donc… »}

Observons : \zh{因为妈妈很温柔，所以我很爱她。} (Yīnwei māma hěn wēnróu, suǒyǐ wǒ hěn ài tā.) « Parce que maman est douce, je l'aime beaucoup. »

\begin{propriete}[因为……所以…… : la cause et la conséquence]
Structure : \textbf{因为 + cause，所以 + conséquence}.\\
Grande différence avec le français : en chinois, on peut (et on préfère souvent) utiliser les \emph{deux} mots ensemble dans la même phrase complexe, alors qu'en français « parce que… donc… » est considéré comme une lourdeur ou une faute de syntaxe. À l'écrit, la paire complète est la forme la plus claire et la plus valorisée à l'examen.
\end{propriete}

\begin{exemplebox}[因为……所以……]
\zh{因为我爱我的家，所以我每天很快乐。} (Yīnwei wǒ ài wǒ de jiā, suǒyǐ wǒ měi tiān hěn kuàilè.) « Parce que j'aime ma famille, je suis heureux chaque jour. »\\
\zh{因为爸爸是老师，所以他帮助我学习。} (Yīnwei bàba shì lǎoshī, suǒyǐ tā bāngzhù wǒ xuéxí.) « Parce que papa est professeur, il m'aide à étudier. »\\
\zh{因为我们都爱我们的家，所以我们的家很温暖。} (Yīnwei wǒmen dōu ài wǒmen de jiā, suǒyǐ wǒmen de jiā hěn wēnnuǎn.) « Parce que nous aimons tous notre famille, notre foyer est chaleureux. »
\end{exemplebox}

\begin{attention}[Piège de traduction]
Ne traduisez jamais « parce que… donc… » mot à mot en français : c'est une faute en français, mais la structure correcte en chinois ! À la version (chinois → français), supprimez le « donc » ou le « parce que » (ex: « Parce que maman est douce, je l'aime beaucoup »). Au thème (français → chinois), un simple « parce que » français peut être rendu par la paire complète \zh{因为……所以……} pour enrichir votre copie.
\end{attention}

\courssub{L'adverbe 都 (dōu) : « tous »}

Observons : \zh{我们都爱我们的家。} (Wǒmen dōu ài wǒmen de jiā.) « Nous aimons tous notre famille. »

\begin{propriete}[Place de 都]
Structure : \textbf{Sujet pluriel + 都 + Verbe}.\\
\zh{都} se place \emph{avant le verbe}, après le sujet. Il insiste sur la totalité du groupe. Attention : le sujet doit être pluriel ou collectif (nous, vous, ils, X et Y) ; on ne dit pas *\zh{我都爱我的家} pour « j'aime toute ma famille » (on dirait \zh{我全家都爱我们的家}).
\end{propriete}

\begin{exemplebox}[都 : la totalité]
\zh{爸爸和妈妈都很忙。} (Bàba hé māma dōu hěn máng.) « Papa et maman sont tous les deux très occupés. »\\
\zh{我们都喜欢足球。} (Wǒmen dōu xǐhuan zúqiú.) « Nous aimons tous le football. »\\
\zh{他们都是学生。} (Tāmen dōu shì xuésheng.) « Ils sont tous élèves. »
\end{exemplebox}

\courssub{Les types de phrases à mobiliser}

Votre texte doit varier les \cle{types de phrases} : c'est un critère de réussite à l'examen (richesse syntaxique).

\begin{center}
\begin{tabularx}{\linewidth}{|X|X|X|}
\hline
\rowcolor{popblueL} Type & Exemple (caractères + pinyin) & Traduction \\
\hline
Déclarative avec 是 & 爸爸是老师。Bàba shì lǎoshī. & Papa est professeur. \\
\hline
Déclarative avec 有 & 我家有五口人。Wǒ jiā yǒu wǔ kǒu rén. & Ma famille compte cinq personnes. \\
\hline
Interrogative avec 吗 & 你爱你的家吗？Nǐ ài nǐ de jiā ma ? & Aimes-tu ta famille ? \\
\hline
Négative avec 不 & 我的家不大。Wǒ de jiā bù dà. & Ma maison n'est pas grande. \\
\hline
Négative avec 没 & 我没有哥哥。Wǒ méiyǒu gēge. & Je n'ai pas de grand frère. \\
\hline
\end{tabularx}
\end{center}

\begin{propriete}[不 ou 没 ? Le choix crucial]
\zh{不} nie les verbes d'état, les adjectifs, les verbes modaux et les habitudes : \zh{不大}, \zh{不喜欢}, \zh{不去}, \zh{不是}.\\
\zh{没} (forme courte de \zh{没有}) nie presque exclusivement le verbe \zh{有} (possession/existence) et les actions accomplies (aspect parfait) : \zh{我没有妹妹。} « Je n'ai pas de petite sœur. », \zh{我没去。} « Je ne suis pas allé. »\\
On ne dit jamais *\zh{不有}.
\end{propriete}

\courssub{Tableau récapitulatif des connecteurs}

\begin{center}
\begin{tabularx}{\linewidth}{|X|X|X|X|}
\hline
\rowcolor{popblueL} Connecteur & Pinyin & Fonction & Exemple traduit \\
\hline
和 & hé & et (entre noms / groupes nominaux) & 爸爸和妈妈 — papa et maman \\
\hline
也 & yě & aussi (avant le verbe) & 我也爱妈妈。 — Moi aussi, j'aime maman. \\
\hline
但是 & dànshì & mais (entre phrases / propositions) & 家不大，但是很温暖。 — petite, mais chaleureuse \\
\hline
因为……所以…… & yīnwei… suǒyǐ… & parce que… (donc)… (paire corrélative) & 因为妈妈温柔，所以我爱她。 \\
\hline
都 & dōu & tous / toutes (avant le verbe, sujet pluriel) & 我们都爱家。 — Nous aimons tous notre foyer. \\
\hline
\end{tabularx}
\end{center}

\courssub{Méthode et entraînement}

\begin{methode}[Rédiger un texte riche en quatre étapes]
\begin{enumerate}
\item Écrivez d'abord \textbf{5 phrases simples} correspondant aux 5 étapes du plan (ouverture, effectif, membres, qualités, conclusion). Utilisez le vocabulaire de la leçon.
\item Reliez ensuite \textbf{deux phrases} avec \zh{但是} (opposition) ou \zh{因为……所以……} (cause) pour enrichir le texte et montrer la maîtrise de la syntaxe complexe.
\item Ajoutez \zh{也} ou \zh{都} là où deux personnes partagent une qualité ou une action (parallélisme).
\item Relisez en vérifiant : \zh{和} seulement entre noms ? \zh{也} et \zh{都} avant le verbe ? Virgule chinoise ，avant \zh{但是} et \zh{所以} ? Ponctuation finale 。?
\end{enumerate}
\end{methode}

\begin{exoresolu}[Relier les phrases avec le bon connecteur]
Énoncé — Reliez les paires de phrases avec le connecteur indiqué, puis traduisez.\\
1. 我的家很小。/ 我的家很温暖。(但是)\\
2. 妈妈很温柔。/ 我很爱她。(因为……所以……)\\
3. 哥哥喜欢足球。/ 我喜欢足球。(也)
\tcblower
Solution détaillée —\\
1. \zh{我的家很小，但是很温暖。} (Wǒ de jiā hěn xiǎo, dànshì hěn wēnnuǎn.) « Ma maison est petite, mais elle est chaleureuse. » On place la virgule chinoise ，puis \zh{但是} ; le sujet n'est pas répété car il est identique.\\
2. \zh{因为妈妈很温柔，所以我很爱她。} (Yīnwei māma hěn wēnróu, suǒyǐ wǒ hěn ài tā.) « Parce que maman est douce, je l'aime beaucoup. » La cause précède, la conséquence suit avec \zh{所以}.\\
3. \zh{哥哥喜欢足球，我也喜欢足球。} (Gēge xǐhuan zúqiú, wǒ yě xǐhuan zúqiú.) « Mon grand frère aime le football, moi aussi j'aime le football. » \zh{也} se place avant le verbe \zh{喜欢}, jamais \zh{和} entre les deux phrases.
\end{exoresolu}

\begin{exercice}[À vous de jouer]
1. Traduisez en chinois : « J'aime papa et maman. Ma maison n'est pas grande, mais elle est chaleureuse. Nous aimons tous notre famille. »\\
2. Corrigez la phrase fautive : *\zh{我爱爸爸，和我爱妈妈。}\\
3. Complétez avec \zh{也}, \zh{都} ou \zh{但是} : 爸爸是老师，妈妈\_\_\_\_是老师。/ 我们\_\_\_\_爱我的家。/ 中文很难，\_\_\_\_很有意思。
\end{exercice}

\begin{corrige}[Exercice « À vous de jouer »]
1. \zh{我爱爸爸和妈妈。我的家不大，但是很温暖。我们都爱我们的家。} (Wǒ ài bàba hé māma. Wǒ de jiā bù dà, dànshì hěn wēnnuǎn. Wǒmen dōu ài wǒmen de jiā.)\\
2. \zh{我爱爸爸，也爱妈妈。} — \zh{和} ne relie pas deux phrases ; on utilise \zh{也} avant le verbe.\\
3. 爸爸是老师，妈妈\zh{也}是老师。/ 我们\zh{都}爱我的家。/ 中文很难，\zh{但是}很有意思。
\end{corrige}

\begin{aretenir}[L'essentiel de la section]
\begin{itemize}
\item Le texte sur la famille suit 5 étapes : ouverture → effectif (avec classifieur \zh{口}) → membres (\zh{是}) → qualités/goûts (\zh{很}, \zh{喜欢}) → conclusion.
\item \zh{和} = « et » entre noms uniquement ; entre phrases, utilisez \zh{也}, \zh{但是}, \zh{因为……所以……}.
\item \zh{也} et \zh{都} se placent toujours \textbf{avant le verbe}.
\item \zh{因为……所以……} s'emploie en paire en chinois (cause + conséquence), contrairement au français.
\item Variez les types de phrases : \zh{是} / \zh{有} / \zh{吗} / \zh{不} / \zh{没}.
\item Distinction vitale : \zh{不} (état, habitude, adjectif) vs \zh{没/没有} (possession, action passée).
\end{itemize}
\end{aretenir}

\coursec{Thème : français → chinois}

Dans la section précédente, nous avons appris à structurer un texte grâce aux connecteurs (\zh{和}, \zh{也}, \zh{都}, \zh{但}, \zh{因为}…\zh{所以}). Nous allons maintenant mobiliser ces acquis pour l'exercice le plus exigeant de l'expression écrite : le \cle{thème}, c'est-à-dire la traduction du \textbf{français vers le chinois}. Contrairement à la version (chinois → français), le thème oblige à \emph{produire} des caractères corrects, à choisir les bons mots-outils et à respecter scrupuleusement l'ordre des mots chinois. C'est l'épreuve reine du baccalauréat : suivons une méthode rigoureuse.

\begin{aretenir}[Rappel : les connecteurs essentiels pour écrire un texte]
\begin{center}
\begin{tabularx}{\linewidth}{|X|X|X|}
\hline
\rowcolor{popblueL} Connecteur & Pinyin & Emploi / Exemple \\
\hline
\zh{和} & \emph{hé} & Lie deux \textbf{noms} : \zh{爸爸和妈妈} (papa et maman). \textbf{Ne lie pas deux phrases.} \\
\hline
\zh{也} & \emph{yě} & « Aussi / aussi » : se place \textbf{avant} le verbe. \zh{我也喜欢。} \\
\hline
\zh{都} & \emph{dōu} & « Tous / toutes » : se place \textbf{avant} le verbe. \zh{我们都爱。} \\
\hline
\zh{但 / 但是} & \emph{dàn / dànshì} & « Mais » : introduit une opposition. \zh{他很忙，但是他来了。} \\
\hline
\zh{因为}…\zh{所以}… & \emph{yīnwèi}…\emph{suǒyǐ}… & « Parce que… donc… » : structure corrélative. \zh{因为下雨，所以我不去。} \\
\hline
\end{tabularx}
\end{center}
\end{aretenir}

\courssub{Observer avant de traduire : deux phrases à comparer}

Avant toute règle, observons deux traductions correctes. Cherchez ce qui change entre le français et le chinois.

\begin{exemplebox}[Observation guidée]
\zh{我爱我的家。} \emph{Wǒ ài wǒ de jiā.} « J'aime ma famille. »\\[4pt]
\zh{我的爸爸是老师。} \emph{Wǒ de bàba shì lǎoshī.} « Mon père est professeur. »
\end{exemplebox}

Que remarque-t-on ?

\begin{itemize}
\item L'ordre général \textbf{Sujet + Verbe + Complément} est le même qu'en français : c'est une bonne nouvelle !
\item Le possessif français « ma / mon » devient \textbf{pronom personnel + \zh{的}} : \zh{我的} (\emph{wǒ de}) = « mon / ma ».
\item Il n'y a \textbf{aucune conjugaison} : \zh{爱} vaut « aime / aimes / aiment » selon le sujet.
\item « Professeur » se traduit par \zh{老师} sans article : le chinois n'a ni « le » ni « un ».
\end{itemize}

\courssub{La méthode du thème en quatre étapes}

Traduire vers le chinois ne s'improvise pas : on suit toujours la même chaîne de décision.

\begin{methode}[Réussir le thème en 4 étapes]
\begin{enumerate}
\item \textbf{Repérer} dans la phrase française le sujet, le(s) adverbe(s), le verbe et le complément. Exemple : « \underline{Mon père} \underline{est} \underline{professeur}. » / « \underline{Nous} \underline{tous} \underline{aimons} \underline{notre famille}. »
\item \textbf{Réorganiser} selon l'ordre chinois canonique : \textbf{Sujet + Adverbe + Verbe + Complément}. En chinois, l'adverbe (\zh{也}, \zh{都}, \zh{很}, \zh{不}, \zh{常}) se place \emph{avant} le verbe, jamais après.
\item \textbf{Traduire} mot à mot avec le vocabulaire connu : possessif = pronom + \zh{的} ; « être + nom » = \zh{是} ; « être + adjectif » = \zh{很} + adjectif (sans \zh{是}) ; « avoir / compter » = \zh{有} + nombre + classificateur + nom.
\item \textbf{Vérifier} chaque caractère : traits, radicaux, caractères proches (\zh{我}/\zh{找}, \zh{人}/\zh{入}, \zh{未}/\zh{末}), mots-outils (\zh{的}/\zh{得}/\zh{地}), et la ponctuation chinoise pleine chasse (\zh{。} \zh{？} \zh{，} \zh{、}).
\end{enumerate}
\end{methode}

\begin{popfigure}
\begin{tikzpicture}[node distance=0.6cm, every node/.style={font=\small, align=center}]
\node[draw=popblue, fill=popblueL, rounded corners, minimum width=2.8cm, minimum height=1.2cm, align=center] (s){\textbf{1. Sujet}\\ \zh{我们}\\ \emph{wǒmen}};
\node[draw=popgreen, fill=popgreenL, rounded corners, minimum width=2.8cm, minimum height=1.2cm, right=of s, align=center] (a){\textbf{2. Adverbe}\\ \zh{也 / 都 / 很 / 不}\\ \emph{(avant le verbe)}};
\node[draw=poporange, fill=poporangeL, rounded corners, minimum width=2.8cm, minimum height=1.2cm, right=of a, align=center] (v){\textbf{3. Verbe}\\ \zh{是 / 爱 / 有 / 喜欢}\\ \emph{shì / ài / yǒu / xǐhuan}};
\node[draw=poppurple, fill=poppurpleL, rounded corners, minimum width=2.8cm, minimum height=1.2cm, right=of v, align=center] (c){\textbf{4. Complément}\\ \zh{我们的家}\\ \emph{wǒmen de jiā}};
\draw[-{Stealth}, thick, popdark] (s) -- (a);
\draw[-{Stealth}, thick, popdark] (a) -- (v);
\draw[-{Stealth}, thick, popdark] (v) -- (c);
\end{tikzpicture}
\legende{La chaîne de traduction : on construit la phrase chinoise case par case, de gauche à droite. L'adverbe s'insère entre le sujet et le verbe.}
\end{popfigure}

\courssub{La règle d'or : le possessif avec 的}

\begin{propriete}[Le possessif : pronom + 的]
Le possessif français « mon / ma / mes / ton / ta / tes / son / sa / ses / notre / votre / leur » se rend en chinois par \textbf{pronom personnel + \zh{的} (de)} :
\begin{center}
\begin{tabularx}{\linewidth}{|X|X|X|}
\hline
\rowcolor{popblueL} Français & Chinois + pinyin & Exemple traduit \\
\hline
mon / ma / mes & \zh{我的} \emph{wǒ de} & \zh{我的妈妈} \emph{wǒ de māma} = ma maman \\
\hline
ton / ta / tes & \zh{你的} \emph{nǐ de} & \zh{你的家} \emph{nǐ de jiā} = ta famille \\
\hline
son / sa / ses & \zh{他的} \emph{tā de} & \zh{他的哥哥} \emph{tā de gēge} = son grand frère \\
\hline
notre / nos & \zh{我们的} \emph{wǒmen de} & \zh{我们的家} \emph{wǒmen de jiā} = notre famille \\
\hline
votre / vos & \zh{你们的} \emph{nǐmen de} & \zh{你们的爸爸} \emph{nǐmen de bàba} = votre père \\
\hline
leur / leurs & \zh{他们的} \emph{tāmen de} & \zh{他们的姐姐} \emph{tāmen de jiějie} = leur grande sœur \\
\hline
\end{tabularx}
\end{center}
\textbf{Exception d'usage courant} : devant un nom de parenté proche, un terme d'affection ou \zh{家}, on \emph{peut} omettre \zh{的} à l'oral et à l'écrit familier : \zh{我家有四口人。} (\emph{Wǒ jiā yǒu sì kǒu rén.}) = « Ma famille compte quatre personnes. » ; \zh{我妈妈很忙。} (\emph{Wǒ māma hěn máng.}). \textbf{À l'examen écrit (thème/version), gardez systématiquement \zh{的}} : il n'est jamais fautif et marque la structure possessive attendue.
\end{propriete}

\courssub{Le piège capital : « être » devant un adjectif}

\begin{attention}[Ne traduis jamais « être » par 是 devant un adjectif]
En français on dit « il \emph{est} grand », « elle \emph{est} gentille ». En chinois, \textbf{l'adjectif fait office de verbe prédicatif} : on n'utilise \textbf{pas} \zh{是}. On écrit :
\begin{itemize}
\item \zh{他很高。} \emph{Tā hěn gāo.} « Il est grand. » (et non *\zh{他是高})
\item \zh{妈妈很温柔。} \emph{Māma hěn wēnróu.} « Maman est douce. » (et non *\zh{妈妈是很温柔})
\item \zh{这道菜很好吃。} \emph{Zhè dào cài hěn hǎochī.} « Ce plat est délicieux. »
\end{itemize}
Le mot \zh{很} (\emph{hěn}, « très ») sert ici de \textbf{liaison obligatoire} (support aspectuel) devant l'adjectif monosyllabique ou bisyllabique ; il n'exprime pas forcément une intensité forte (traduction souvent : « est » et non « est très »).
\textbf{Réserve importante} : \zh{是} s'emploie devant un \emph{nom} (identité, classification, métier) : \zh{他是老师。} « Il est professeur. » ; \zh{这是书。} « Ceci est un livre. »
\end{attention}

\courssub{Le classificateur \zh{口} pour les membres de la famille}

\begin{propriete}[Compter les personnes de la famille : \zh{有} + nombre + \zh{口} + \zh{人}]
Pour exprimer « Ma famille compte X personnes », on utilise le verbe \zh{有} (avoir / il y a) suivi du nombre et du classificateur \zh{口} (\emph{kǒu}), spécifique aux bouches à nourrir / membres du foyer.
\begin{center}
\begin{tabularx}{\linewidth}{|X|X|X|}
\hline
\rowcolor{popblueL} Français & Chinois + pinyin & Structure \\
\hline
Ma famille compte 3 personnes. & \zh{我家有三口人。} \emph{Wǒ jiā yǒu sān kǒu rén.} & Sujet \zh{我家} + \zh{有} + Nombre + \zh{口} + \zh{人} \\
\hline
Combien de personnes chez toi ? & \zh{你家有几口人？} \emph{Nǐ jiā yǒu jǐ kǒu rén?} & \zh{几} (\emph{jǐ}) pour « combien » (nombre < 10) \\
\hline
Nous sommes 5 personnes. & \zh{我们家有五口人。} \emph{Wǒmen jiā yǒu wǔ kǒu rén.} & \zh{我们家} = notre foyer \\
\hline
\end{tabularx}
\end{center}
\textbf{Attention} : \zh{几} (\emph{jǐ}) s'utilise pour demander un petit nombre (généralement < 10). Pour un nombre potentiellement grand, on utilise \zh{多少} (\emph{duōshao}) : \zh{你们学校有多少学生？}
\end{propriete}

\courssub{Les cinq phrases modèles du thème}

Voici les cinq phrases types de la leçon, traduites et analysées. Apprenez-les par cœur : elles couvrent tous les cas de l'examen de Première.

\begin{exemplebox}[Phrase 1 : phrase verbale simple (verbe d'action/état)]
\zh{我爱我的家。} \emph{Wǒ ài wǒ de jiā.} « J'aime ma famille. »\\
Sujet \zh{我} + verbe \zh{爱} + complément \zh{我的家} (possessif avec \zh{的}).
\end{exemplebox}

\begin{exemplebox}[Phrase 2 : phrase avec 是 + nom (identité / métier)]
\zh{我的爸爸是老师。} \emph{Wǒ de bàba shì lǎoshī.} « Mon père est professeur. »\\
« Être + nom » = \zh{是}. Pas d'article devant \zh{老师}. \zh{的} conservé après \zh{我的} (écrit formel).
\end{exemplebox}

\begin{exemplebox}[Phrase 3 : phrase avec 有 + classificateur 口]
\zh{我家有四口人。} \emph{Wǒ jiā yǒu sì kǒu rén.} « Ma famille compte quatre personnes. »\\
« Compter / avoir » = \zh{有}. Sujet \zh{我家} (mon foyer, \zh{的} omis). Classificateur \zh{口} obligatoire.
\end{exemplebox}

\begin{exemplebox}[Phrase 4 : coordination avec 也 (aussi) et 和 (et)]
\zh{哥哥喜欢足球，我也喜欢。} \emph{Gēge xǐhuan zúqiú, wǒ yě xǐhuan.} « Mon frère aime le football, et moi aussi. »\\
« Moi aussi » = \zh{我也} + verbe. L'adverbe \zh{也} se place \textbf{avant} le verbe \zh{喜欢}, jamais en fin de phrase. L'objet \zh{足球} est ellipté dans la seconde proposition.
\end{exemplebox}

\begin{exemplebox}[Phrase 5 : question fermée avec 吗]
\zh{你爱你的妈妈吗？} \emph{Nǐ ài nǐ de māma ma?} « Tu aimes ta maman ? »\\
La question fermée (oui/non) se forme en ajoutant la particule \zh{吗} à la fin de la phrase affirmative, \textbf{sans changer l'ordre des mots}. Ponctuation : point d'interrogation chinois \zh{？}.
\end{exemplebox}

\begin{center}
\begin{tabularx}{\linewidth}{|X|X|X|}
\hline
\rowcolor{popblueL} Type de phrase française & Structure chinoise canonique & Exemple modèle \\
\hline
Affirmative verbale simple & Sujet + Verbe + Complément & \zh{我爱我的家。} \emph{Wǒ ài wǒ de jiā.} \\
\hline
« Être + nom / métier » & Sujet + \zh{是} + Nom & \zh{我的爸爸是老师。} \emph{Wǒ de bàba shì lǎoshī.} \\
\hline
« Être + adjectif » & Sujet + \zh{很} + Adjectif (\textbf{sans} \zh{是}) & \zh{妈妈很温柔。} \emph{Māma hěn wēnróu.} \\
\hline
« Avoir / Compter (famille) » & Sujet + \zh{有} + Nombre + \zh{口} + \zh{人} & \zh{我家有四口人。} \emph{Wǒ jiā yǒu sì kǒu rén.} \\
\hline
« Aussi » (adverbe) & Sujet + \zh{也} + Verbe & \zh{我也喜欢。} \emph{Wǒ yě xǐhuan.} \\
\hline
Question oui/non & Phrase affirmative + \zh{吗？} & \zh{你爱你的妈妈吗？} \emph{Nǐ ài nǐ de māma ma?} \\
\hline
\end{tabularx}
\end{center}

\courssub{Exercice résolu et entraînement}

\begin{exoresolu}[Thème guidé : « Nous aimons tous notre famille. »]
Traduisez en chinois : « Nous aimons tous notre famille. »
\tcblower
\textbf{Étape 1 — Repérage.} Sujet : « nous » (\zh{我们}) ; Adverbe : « tous » (\zh{都}) ; Verbe : « aimons » (\zh{爱}) ; Complément : « notre famille » (\zh{我们的家}).\\
\textbf{Étape 2 — Réorganisation.} Ordre chinois : Sujet + Adverbe + Verbe + Complément $\rightarrow$ \zh{我们} + \zh{都} + \zh{爱} + \zh{我们的家}.\\
\textbf{Étape 3 — Traduction mot à mot.} \zh{我们都爱我们的家}.\\
\textbf{Étape 4 — Vérification.}
\begin{itemize}
\item \zh{我} : 7 traits (ordre : horizontal, oblique gauche, oblique droite, horizontal, vertical, horizontal, horizontal). Ne pas confondre avec \zh{找} (trouver, 7 traits aussi mais clé \zh{手}).
\item \zh{爱} : clé du cœur \zh{心} en bas (simplifié), 10 traits. Traditionnel \zh{愛} a \zh{心} au milieu.
\item \zh{都} : clé \zh{阝} (oreille droite), 10 traits. Ne pas confondre \zh{都} (tous) et \zh{者} (celui qui).
\item Ponctuation : point plein chasse \zh{。}
\end{itemize}
\textbf{Solution finale :} \zh{我们都爱我们的家。} \emph{Wǒmen dōu ài wǒmen de jiā.}
\end{exoresolu}

\begin{exercice}[À vous de traduire — Thème d'entraînement]
Traduisez en chinois (caractères simplifiés + pinyin avec tons + traduction française de vérification).
\begin{enumerate}
\item J'aime mon père et ma mère.
\item Sa sœur (aînée) est étudiante.
\item Ta famille compte combien de personnes ? — Cinq personnes.
\item Mon grand frère est très grand.
\item Vous aimez aussi le football ?
\end{enumerate}
\end{exercice}

\begin{corrige}[Exercice « À vous de traduire »]
\begin{enumerate}
\item \zh{我爱我的爸爸和妈妈。} \emph{Wǒ ài wǒ de bàba hé māma.} — « et » entre deux noms = \zh{和}. \zh{的} gardé pour l'écrit formel.
\item \zh{他的姐姐是学生。} \emph{Tā de jiějie shì xuésheng.} — « Sa sœur aînée » = \zh{姐姐} (par défaut si âge non précisé). « Être + nom » = \zh{是}. \zh{学生} sans article.
\item \zh{你家有几口人？——五口人。} \emph{Nǐ jiā yǒu jǐ kǒu rén? — Wǔ kǒu rén.} — « Combien » (petit nombre) = \zh{几}. Réponse elliptique : nombre + classificateur + nom. \zh{的} omis après \zh{你家} (usage courant).
\item \zh{我的哥哥很高。} \emph{Wǒ de gēge hěn gāo.} — « Mon grand frère » = \zh{哥哥}. « Est grand » = \zh{很高} (\textbf{pas de} \zh{是} !). \zh{很} lie le sujet à l'adjectif.
\item \zh{你们也喜欢足球吗？} \emph{Nǐmen yě xǐhuan zúqiú ma?} — « Vous » pluriel = \zh{你们}. \zh{也} avant le verbe \zh{喜欢}. \zh{吗} en fin de phrase pour la question.
\end{enumerate}
\end{corrige}

\begin{aretenir}[Les 5 réflexes du thème — Fiche de révision examen]
\begin{enumerate}
\item \textbf{Possessif} = pronom + \zh{的} : \zh{我的妈妈}, \zh{你们的家}. (Garder \zh{的} à l'écrit).
\item \textbf{« Être »} : devant \textbf{nom} $\rightarrow$ \zh{是} ; devant \textbf{adjectif} $\rightarrow$ \zh{很} + adjectif (\textbf{jamais} \zh{是}).
\item \textbf{Adverbes} (\zh{也}, \zh{都}, \zh{很}, \zh{不}, \zh{常}, \zh{真}) toujours \textbf{avant} le verbe.
\item \textbf{Question oui/non} = phrase affirmative + \zh{吗？} (ordre inchangé). Question « combien » (petit nombre) = \zh{几}.
\item \textbf{Relire} : traits (ordre : haut→bas, gauche→droite, horiz.→vert.), radicaux (clés), homophones (\zh{的}/\zh{得}/\zh{地}), ponctuation chinoise \zh{。} \zh{？} \zh{，} \zh{、}.
\end{enumerate}
\end{aretenir}

Le thème est maintenant maîtrisé. Dans la section suivante, nous ferons le chemin inverse — la \textbf{version}, du chinois vers le français — qui demande d'autres réflexes de lecture et de reformulation.

\coursec{Version : chinois → français}

Dans la section précédente, nous avons appris à passer du français au chinois (\cle{thème}) : il fallait penser en chinois, respecter l'ordre des mots et choisir les bons caractères. Nous faisons maintenant le chemin inverse : partir d'une phrase chinoise pour produire un français correct et naturel. C'est la \cle{version}. Exercice en apparence plus facile, il cache pourtant des pièges redoutables : le mot à mot naïf produit un français bancal, et certaines petites syllabes chinoises (\zh{很}, \zh{都}, \zh{吗}) ne se traduisent jamais mécaniquement. Cette section vous donne une méthode en quatre étapes, cinq phrases modèles traduites et commentées, un texte guidé, et la liste des pièges à éviter.

\courssub{Observer d'abord : deux traductions possibles}

Lisez ces deux propositions de traduction de la même phrase :

\zh{我的妈妈很温柔。} \emph{(Wǒ de māma hěn wēnróu.)}

\begin{itemize}
\item Traduction A : « Ma maman très douce. »
\item Traduction B : « Ma maman est très douce. »
\end{itemize}

La traduction A est un mot à mot : elle a oublié que le français exige un verbe (« est »), alors que le chinois n'en a pas besoin devant un adjectif. La traduction B est correcte. \textbf{La version, ce n'est pas transporter les mots : c'est recréer une phrase française complète et naturelle.} Toute cette section repose sur cette idée.

\courssub{La méthode en quatre étapes}

\begin{methode}[Réussir une version en quatre étapes]
\begin{enumerate}
\item \textbf{Lire la phrase entière}, jusqu'au point final \zh{。}, sans commencer à traduire. Identifier le sens global : de qui parle-t-on ? Que dit-on de cette personne ?
\item \textbf{Repérer le sujet et le verbe (ou l'adjectif-verbe)}. En chinois comme en français, c'est le squelette de la phrase : \zh{妈妈} (sujet) + \zh{很温柔} (prédicat).
\item \textbf{Traduire les connecteurs et les mots-outils} : \zh{因为……所以……} (parce que… donc…), \zh{都} (tous), \zh{吗} (marque de question), \zh{的} (lien de possession). Ce sont eux qui donnent la logique de la phrase.
\item \textbf{Rédiger un français correct} au propre : verbe conjugué, articles, accords, ponctuation française. Relire en se demandant : « Est-ce qu'un francophone dirait vraiment cela ? »
\end{enumerate}
\end{methode}

\begin{attention}[Le piège du mot à mot]
Au brouillon, traduisez d'abord \emph{mot à mot} pour ne rien oublier, puis \emph{reformulez} en français naturel. Ne rendez jamais le mot à mot comme copie finale : « Je famille avoir cinq bouche personne » n'est pas du français, c'est du chinois déguisé. Le correcteur juge le français final, pas votre fidélité syllabe par syllabe.
\end{attention}

\courssub{Les cinq phrases modèles, traduites et commentées}

\begin{popfigure}
\begin{tikzpicture}[
  zh/.style={draw=popblue, fill=popblueL, rounded corners=2pt, text width=5.6cm, align=center, inner sep=5pt, font=\small},
  fr/.style={draw=poporange, fill=poporangeL, rounded corners=2pt, text width=5.6cm, align=center, inner sep=5pt, font=\small},
  fl/.style={-{Stealth[length=2.5mm]}, thick, popdark}
]
\node[zh] (z1) at (0,0) {我的妈妈\textcolor{poppurple}{\textbf{很}}温柔。};
\node[fr, right=2.2cm of z1] (f1) {Ma maman \textbf{est} très douce.};
\draw[fl] (z1) -- (f1);
\node[zh, below=0.55cm of z1] (z2) {我家有五口人。};
\node[fr, right=2.2cm of z2] (f2) {Ma famille compte cinq personnes.};
\draw[fl] (z2) -- (f2);
\node[zh, below=0.55cm of z2] (z3) {妹妹很可爱，我们\textcolor{poppurple}{\textbf{都}}喜欢她。};
\node[fr, right=2.2cm of z3] (f3) {Ma petite sœur est très mignonne, nous l'aimons \textbf{tous}.};
\draw[fl] (z3) -- (f3);
\node[zh, below=0.55cm of z3] (z4) {\textcolor{poppurple}{\textbf{因为}}爸爸是医生，\textcolor{poppurple}{\textbf{所以}}他很忙。};
\node[fr, right=2.2cm of z4] (f4) {\textbf{Parce que} papa est médecin, il est très occupé.};
\draw[fl] (z4) -- (f4);
\node[zh, below=0.55cm of z4] (z5) {你爱你的家\textcolor{poppurple}{\textbf{吗}}？};
\node[fr, right=2.2cm of z5] (f5) {Aimes-tu ta famille \textbf{?}};
\draw[fl] (z5) -- (f5);
\end{tikzpicture}
\legende{Les cinq phrases de la version : les connecteurs et mots-outils sont surlignés en violet.}
\end{popfigure}

\textbf{Phrase 1.} \zh{我的妈妈很温柔。} \emph{(Wǒ de māma hěn wēnróu.)} → « Ma maman est très douce. »

Décomposition : \zh{我的} (ma, avec le possessif \zh{的}) + \zh{妈妈} (maman, sujet) + \zh{很温柔} (très douce, prédicat adjectival). Le chinois n'a pas de verbe « être » ici : c'est \zh{很} qui fait la liaison. En français, il faut \textbf{rajouter le verbe « être » conjugué}. C'est l'ajout le plus fréquent en version.

\textbf{Phrase 2.} \zh{我家有五口人。} \emph{(Wǒ jiā yǒu wǔ kǒu rén.)} → « Ma famille compte cinq personnes. »

Structure : \zh{我家} (ma famille, sujet) + \zh{有} (avoir, il y a) + \zh{五口人} (cinq personnes, avec le classificateur \zh{口}). Le mot à mot donnerait « ma famille a cinq bouches-personnes » : inacceptable. En français naturel, on dit « ma famille \textbf{compte} cinq personnes » ou « nous sommes cinq dans ma famille ». Les deux sont acceptées ; la première est plus élégante.

\textbf{Phrase 3.} \zh{妹妹很可爱，我们都喜欢她。} \emph{(Mèimei hěn kě'ài, wǒmen dōu xǐhuan tā.)} → « Ma petite sœur est très mignonne, nous l'aimons tous. »

Deux propositions reliées par la virgule chinoise \zh{，}. Points de vigilance : \zh{妹妹} seul signifie « petite sœur » ; le contexte permet de traduire « ma petite sœur ». \zh{都} (tous) se place \emph{avant} le verbe en chinois, mais « tous » se place \emph{après} le verbe en français : \zh{我们都喜欢她} → « nous l'aimons \textbf{tous} », jamais « nous tous aimons elle ». Enfin, \zh{她} (elle) devient le pronom complément « l' » placé avant le verbe français.

\textbf{Phrase 4.} \zh{因为爸爸是医生，所以他很忙。} \emph{(Yīnwei bàba shì yīshēng, suǒyǐ tā hěn máng.)} → « Parce que papa est médecin, il est très occupé. »

Le chinois utilise volontiers le couple \zh{因为……所以……} (« parce que… donc… ») dans la même phrase. Le français, lui, n'accepte pas « parce que… donc… » ensemble : on garde \textbf{un seul} des deux connecteurs. Traductions acceptées : « Parce que papa est médecin, il est très occupé » ou « Papa est médecin, donc il est très occupé » ou encore « Comme papa est médecin, il est très occupé ». Ici, \zh{是} est bien le verbe « être » (devant un nom : \zh{医生}).

\textbf{Phrase 5.} \zh{你爱你的家吗？} \emph{(Nǐ ài nǐ de jiā ma?)} → « Aimes-tu ta famille ? »

La particule finale \zh{吗} signale une question fermée (réponse oui/non). En français, il faut le marquer par \textbf{l'inversion du sujet} (« aimes-tu ») ou par « est-ce que » (« Est-ce que tu aimes ta famille ? »), avec le point d'interrogation. Notez que \zh{家} peut se traduire par « famille » ou « maison/foyer » selon le contexte ; avec le verbe \zh{爱} (aimer), « famille » est le choix naturel.

\courssub{Les mots-outils décisifs : 很, 都, 吗}

\begin{attention}[Ne traduisez pas systématiquement 很 par « très »]
Devant un adjectif, \zh{很} sert le plus souvent de \cle{simple copule} (liaison sujet-adjectif), sans idée d'intensité. \zh{她很好} \emph{(Tā hěn hǎo.)} signifie « elle va bien » ou « elle est bien », pas forcément « elle est très bien ». Pour exprimer un vrai « très » insistant, le chinois utilise \zh{非常} (fēicháng) ou \zh{太……了} (tài… le). En version, traduisez \zh{很} par « très » seulement si l'intensité est claire dans le contexte ; sinon, contentez-vous du verbe « être » : \zh{他很忙} → « il est occupé » reste correct.
\end{attention}

\begin{propriete}[La particule interrogative 吗]
\zh{吗} (ma, ton neutre) placée \textbf{à la fin} d'une phrase déclarative la transforme en question fermée (réponse oui/non), \textbf{sans changer l'ordre des mots}. Schéma : Phrase déclarative + \zh{吗} + \zh{？}
\begin{itemize}
\item \zh{你爱我。} \emph{(Nǐ ài wǒ.)} « Tu m'aimes. » → \zh{你爱我吗？} \emph{(Nǐ ài wǒ ma?)} « M'aimes-tu ? »
\item \zh{他是医生。} \emph{(Tā shì yīshēng.)} « Il est médecin. » → \zh{他是医生吗？} \emph{(Tā shì yīshēng ma?)} « Est-il médecin ? »
\end{itemize}
En version, \zh{吗} disparaît : il se transforme en inversion du sujet ou en « est-ce que », plus le point d'interrogation. Ne le traduisez jamais par un mot (« est-ce que vrai que… ? » est une erreur). Attention : \zh{吗} ne s'emploie \textbf{pas} si la phrase contient déjà un mot interrogatif comme \zh{几} (combien), \zh{谁} (qui) ou \zh{什么} (quoi).
\end{propriete}

\begin{exemplebox}[都 : « tous », placé avant le verbe en chinois, après en français]
\zh{我们都爱我们的家。} \emph{(Wǒmen dōu ài wǒmen de jiā.)} → « Nous aimons tous notre famille. »\\
\zh{他们都是学生。} \emph{(Tāmen dōu shì xuésheng.)} → « Ils sont tous élèves. »\\
\zh{我爸爸妈妈都喜欢足球。} \emph{(Wǒ bàba māma dōu xǐhuan zúqiú.)} → « Mon père et ma mère aiment tous les deux le football. »
\end{exemplebox}

\courssub{Tableau récapitulatif des correspondances}

\begin{center}
\begin{tabularx}{\linewidth}{|X|X|X|X|}
\hline
\rowcolor{popblueL} Caractères & Pinyin & Fonction & Traduction en version \\
\hline
很 & hěn & copule / intensité & « est » (souvent) ou « très » \\
\hline
都 & dōu & totalité & « tous » (après le verbe) \\
\hline
吗 & ma & question fermée & inversion ou « est-ce que » + ? \\
\hline
的 & de & possession & « de », adjectif possessif \\
\hline
因为……所以…… & yīnwei… suǒyǐ… & cause & « parce que » \emph{ou} « donc » (un seul) \\
\hline
有 & yǒu & possession / existence & « avoir », « compter », « il y a » \\
\hline
口 & kǒu & classificateur (membres de la famille) & ne se traduit pas \\
\hline
她 & tā & pronom féminin & « elle », « la/l' » \\
\hline
也 & yě & addition & « aussi » (après le verbe ou en fin de phrase) \\
\hline
\end{tabularx}
\end{center}

Notez la ligne sur les classificateurs : \zh{口} dans \zh{五口人} ne se traduit \textbf{pas}. Comme tous les classificateurs, il est une exigence de la grammaire chinoise, sans équivalent français. Le traduire par « bouche » serait une erreur de version. De même, \zh{也} (aussi) se place avant le verbe en chinois (\zh{我们也喜欢她}), mais « aussi » se place \textbf{après} le verbe français ou en fin de phrase : « nous l'aimons aussi ».

\courssub{Texte d'entraînement : une petite version guidée}

\begin{textebox}[Mon amie Aïcha 我的好朋友]
我有一个好朋友，她叫Aïcha。她十三岁，是雅温得的学生。\\
\emph{Wǒ yǒu yí ge hǎo péngyou, tā jiào Aïcha. Tā shísān suì, shì Yǎwēndé de xuésheng.}\\
她家有六口人：爸爸、妈妈、两个哥哥、妹妹和她。\\
\emph{Tā jiā yǒu liù kǒu rén : bàba, māma, liǎng ge gēge, mèimei hé tā.}\\
她的爸爸是老师，妈妈是医生，他们都很忙。\\
\emph{Tā de bàba shì lǎoshī, māma shì yīshēng, tāmen dōu hěn máng.}\\
因为哥哥们喜欢足球，所以星期天他们一起看电视。\\
\emph{Yīnwei gēgemen xǐhuan zúqiú, suǒyǐ xīngqītiān tāmen yìqǐ kàn diànshì.}\\
Aïcha很爱她的家，我们也都很喜欢她！\\
\emph{Aïcha hěn ài tā de jiā, wǒmen yě dōu hěn xǐhuan tā !}\\[4pt]
\textbf{Traduction de référence :} J'ai une bonne amie, elle s'appelle Aïcha. Elle a treize ans, elle est élève à Yaoundé. Sa famille compte six personnes : son père, sa mère, ses deux grands frères, sa petite sœur et elle. Son père est professeur, sa mère est médecin, ils sont tous les deux très occupés. Parce que ses grands frères aiment le football, le dimanche ils regardent la télévision ensemble. Aïcha aime beaucoup sa famille, et nous aussi, nous l'aimons toutes beaucoup !
\end{textebox}

Glossaire utile : \zh{叫} (jiào, s'appeler), \zh{岁} (suì, ans — âge), \zh{两个} (liǎng ge, deux — devant un classificateur on dit \zh{两} et non \zh{二}), \zh{一起} (yìqǐ, ensemble), \zh{也} (yě, aussi), \zh{哥哥们} (gēgemen, les grands frères, avec le pluriel \zh{们}).

Questions de compréhension : 1) Combien de personnes y a-t-il dans la famille d'Aïcha ? 2) Que font ses parents dans la vie ? 3) Pourquoi la famille regarde-t-elle la télévision le dimanche ?

\courssub{Exercice résolu et entraînement}

\begin{exoresolu}[Version guidée]
Traduisez en français : \zh{我的哥哥很高，他很喜欢篮球，我们都爱他。} \emph{(Wǒ de gēge hěn gāo, tā hěn xǐhuan lánqiú, wǒmen dōu ài tā.)}
\tcblower
\textbf{Étape 1 — Lecture globale :} on parle du grand frère ; trois informations : sa taille, son goût, nos sentiments.\\
\textbf{Étape 2 — Sujets et prédicats :} \zh{哥哥} + \zh{很高} ; \zh{他} + \zh{很喜欢篮球} ; \zh{我们} + \zh{都爱他}.\\
\textbf{Étape 3 — Mots-outils :} \zh{的} (mon), \zh{很} (copule, deux fois), \zh{都} (tous, avant le verbe).\\
\textbf{Étape 4 — Rédaction française :} « Mon grand frère est grand, il aime beaucoup le basket-ball, nous l'aimons tous. »\\
\textbf{Vérification :} verbe « est » ajouté devant l'adjectif ; « tous » placé après le verbe ; \zh{他} rendu par « l' » devant « aimons ». Traduction correcte et naturelle.
\end{exoresolu}

\begin{exercice}[À vous de traduire]
Traduisez en français correct :
\begin{enumerate}
\item \zh{我的妹妹很可爱。} \emph{(Wǒ de mèimei hěn kě'ài.)}
\item \zh{你家有几口人？} \emph{(Nǐ jiā yǒu jǐ kǒu rén?)}
\item \zh{因为妈妈很忙，所以爸爸做饭。} \emph{(Yīnwei māma hěn máng, suǒyǐ bàba zuò fàn.)}
\item \zh{你爱你的妹妹吗？} \emph{(Nǐ ài nǐ de mèimei ma?)}
\end{enumerate}
\end{exercice}

\begin{corrige}[Exercice « À vous de traduire »]
\begin{enumerate}
\item « Ma petite sœur est très mignonne. » — verbe « est » ajouté ; \zh{很} rendu par « très » (ou rien : « est mignonne » est accepté).
\item « Combien de personnes compte ta famille ? » ou « Vous êtes combien dans ta famille ? » — \zh{几} (jǐ) = « combien » ; question sans \zh{吗} car il y a déjà un mot interrogatif.
\item « Parce que maman est très occupée, c'est papa qui fait la cuisine. » ou « Maman est très occupée, donc papa fait la cuisine. » — un seul connecteur en français ; \zh{做饭} (zuò fàn) = « faire la cuisine / à manger ».
\item « Aimes-tu ta petite sœur ? » ou « Est-ce que tu aimes ta petite sœur ? » — \zh{吗} rendu par l'inversion ou « est-ce que », point d'interrogation obligatoire.
\end{enumerate}
\end{corrige}

\begin{aretenir}[Les cinq réflexes de la version]
\begin{itemize}
\item Lire la phrase \textbf{entière} avant de traduire.
\item Ajouter en français le verbe « être » devant tout adjectif chinois précédé de \zh{很}.
\item Ne pas traduire \zh{很} systématiquement par « très » ni les classificateurs (\zh{口}, \zh{个}) par un mot.
\item Placer « tous » (\zh{都}) et « aussi » (\zh{也}) \textbf{après} le verbe français ou en fin de phrase, alors qu'en chinois ils précèdent le verbe.
\item Rendre \zh{吗} par une inversion ou « est-ce que », et ne garder qu'\textbf{un seul} connecteur pour \zh{因为……所以……}.
\end{itemize}
\end{aretenir}

Vous savez désormais traduire du chinois vers le français avec méthode : observer, découper, traduire les mots-outils, rédiger. La dernière section vous présentera la grille d'évaluation utilisée à l'examen et un entraînement complet en conditions réelles, pour transformer ces réflexes en points.

\coursec{Grille d'évaluation et entraînement à l'examen}

Dans la section précédente, tu as appris à traduire du chinois vers le français (version) : lire, comprendre, restituer. Il te reste maintenant l'étape décisive : savoir \textbf{comment ton texte sera noté} le jour de l'évaluation, et t'entraîner dans les conditions réelles de l'examen. Un bon élève de chinois n'est pas celui qui écrit beaucoup, mais celui qui écrit \cle{juste} : caractères exacts, phrases correctes, consigne respectée. Cette dernière section te donne la grille officielle de notation, une méthode de relecture, et un sujet d'entraînement complet avec sa correction.

\courssub{Comment est évaluée une production écrite de chinois ?}

Observe d'abord les deux textes suivants, écrits par deux élèves sur le même sujet : « Présente ta famille ».

\begin{textebox}[Texte A — court mais précis]
我家有四口人：爸爸、妈妈、姐姐和我。\\
Wǒ jiā yǒu sì kǒu rén : bàba, māma, jiějie hé wǒ.\\
爸爸是医生，妈妈很好。我们都爱我们的家。\\
Bàba shì yīshēng, māma hěn hǎo. Wǒmen dōu ài wǒmen de jiā.\\
« Ma famille compte quatre personnes : papa, maman, ma grande sœur et moi. Papa est médecin, maman est très gentille. Nous aimons tous notre famille. »
\end{textebox}

\begin{textebox}[Texte B — long mais fautif]
我家有六口人。爸爸是很好，妈妈很高，我和哥哥和姐姐都爱足球，我们的家是大，我爱我的家，我的家很好，我们每天都吃饭一起……\\
« Ma famille compte six personnes. Papa est très bien, maman est grande, mon frère et ma sœur et moi aimons le football, notre maison est grand, j'aime ma famille… »\\
(texte contenant : \zh{是} + adjectif, \zh{和} enchaîné comme le « et » français, oubli de \zh{的}, ordre des mots fautif : \zh{吃饭一起} au lieu de \zh{一起吃饭})
\end{textebox}

Lequel obtient la meilleure note ? \textbf{Le texte A}, sans hésitation. Il est court, mais chaque caractère est exact, chaque phrase a une structure correcte, et la consigne est respectée. Le texte B accumule les fautes typiques du francophone. Retenons la leçon : en chinois, \emph{la précision prime sur la longueur}.

\begin{propriete}[Les cinq critères officiels d'évaluation]
Une production écrite de chinois en classe de Première est notée selon cinq critères :
\begin{enumerate}
\item \textbf{Exactitude des caractères} : traits complets, bons radicaux, aucun caractère inventé ou déformé ;
\item \textbf{Correction grammaticale} : ordre des mots Sujet + Verbe + Complément, verbe présent dans chaque phrase, pas de \zh{是} devant un adjectif, \zh{的} possessif bien placé ;
\item \textbf{Connecteurs et richesse} : au moins deux mots de liaison (\zh{和}, \zh{也}, \zh{都}, \zh{因为}…) et un vocabulaire varié ;
\item \textbf{Respect de la consigne} : thème respecté, nombre de phrases demandé, destinataire pris en compte ;
\item \textbf{Présentation} : un caractère par case, traits nets, texte lisible.
\end{enumerate}
\end{propriete}

\begin{popfigure}
\begin{tikzpicture}[x=3.1cm, y=0.85cm]
\fill[popblue] (0,5) rectangle (2,6);
\fill[popblue] (2,5) rectangle (3,6);
\node[white, font=\bfseries] at (1,5.5) {Critère};
\node[white, font=\bfseries] at (2.5,5.5) {Points};
\fill[popblueL] (0,4) rectangle (3,5);
\node[anchor=west] at (0.08,4.5) {Exactitude des caractères (traits, radicaux)};
\node at (2.5,4.5) {\textbf{/5}};
\fill[white] (0,3) rectangle (3,4);
\node[anchor=west] at (0.08,3.5) {Correction grammaticale};
\node at (2.5,3.5) {\textbf{/5}};
\fill[poporangeL] (0,2) rectangle (3,3);
\node[anchor=west] at (0.08,2.5) {Connecteurs et richesse du vocabulaire};
\node at (2.5,2.5) {\textbf{/4}};
\fill[white] (0,1) rectangle (3,2);
\node[anchor=west] at (0.08,1.5) {Respect de la consigne et du thème};
\node at (2.5,1.5) {\textbf{/4}};
\fill[popgreenL] (0,0) rectangle (3,1);
\node[anchor=west] at (0.08,0.5) {Présentation (cases, traits nets)};
\node at (2.5,0.5) {\textbf{/2}};
\draw[popdark, thick] (0,0) rectangle (3,6);
\draw[popdark] (2,0) -- (2,6);
\foreach \y in {1,2,3,4,5} {\draw[popdark] (0,\y) -- (3,\y);}
\node[font=\bfseries, popdark] at (1.5,-0.6) {Total : /20};
\end{tikzpicture}
\legende{Barème indicatif d'une production écrite de chinois sur 20 points}
\end{popfigure}

\begin{aretenir}[Barème indicatif sur 20]
Caractères \textbf{/5} — Grammaire \textbf{/5} — Connecteurs et richesse \textbf{/4} — Respect de la consigne \textbf{/4} — Présentation \textbf{/2}. Deux fautes de caractères peuvent déjà coûter un point : chaque trait compte.
\end{aretenir}

\courssub{Auto-évaluation : la liste de vérification avant de rendre ta copie}

Avant de remettre ton texte, prends deux minutes pour le relire en te posant ces questions, dans l'ordre :

\begin{center}
\begin{tabularx}{\linewidth}{|X|X|}\hline
\rowcolor{popblueL} Question de vérification & Exemple à contrôler \\ \hline
Chaque caractère a-t-il le bon radical ? & \zh{妈} (radical 女, femme) et non un caractère inventé \\ \hline
Chaque phrase a-t-elle un verbe ou un adjectif-verbe ? & \zh{妈妈很好。} (adjectif seul, sans \zh{是}) \\ \hline
L'ordre des mots est-il Sujet + Verbe + Complément ? & \zh{我爱我的家。} ; le complément de temps ou de lieu avant le verbe : \zh{我们一起吃饭。} \\ \hline
Ai-je mis \zh{的} pour la possession ? & \zh{我的家}, \zh{我们的家} \\ \hline
Ai-je utilisé au moins 2 connecteurs ? & \zh{和}, \zh{也}, \zh{都}, \zh{因为} \\ \hline
Ai-je respecté le nombre de phrases demandé ? & 5 à 8 phrases : ni 3, ni 12 \\ \hline
La présentation est-elle soignée ? & un caractère par case, traits nets \\ \hline
\end{tabularx}
\end{center}

\begin{methode}[Se relire à l'examen : la méthode C-O-N-N-E-C-T]
Retiens ce mot-code français pour ne rien oublier lors de ta relecture :
\begin{enumerate}
\item \textbf{C} comme \cle{Caractères} : chaque caractère est-il complet, avec le bon radical et tous ses traits ?
\item \textbf{O} comme \cle{Ordre des mots} : Sujet + Verbe + Complément ; le complément de lieu ou de temps est-il avant le verbe ?
\item \textbf{N} comme \cle{Négation} : \zh{不} ou \zh{没} sont-ils bien placés devant le verbe ? (\zh{没} pour nier \zh{有} : \zh{我没有哥哥。})
\item \textbf{N} comme \cle{Nombre de phrases} : la consigne demande 5 à 8 phrases ; compte-les.
\item \textbf{E} comme \cle{的} (possessif) : ai-je écrit \zh{我的}, \zh{我们的}, \zh{爸爸的} partout où il le faut ?
\item \textbf{C} comme \cle{Connecteurs} : ai-je au moins deux mots de liaison différents ?
\item \textbf{T} comme \cle{Tons / pinyin} : si le pinyin est demandé, chaque syllabe porte-t-elle sa marque de ton ?
\end{enumerate}
\end{methode}

\begin{attention}[La présentation compte vraiment]
Erreur fréquente des francophones : croire que seul le contenu est noté. Faux ! En chinois, la \cle{précision graphique} est un critère officiel. Écris \textbf{un caractère par case}, sans déborder, avec des traits nets et dans le bon sens (haut → bas, gauche → droite). Un caractère bâclé peut devenir un autre caractère : \zh{人} (rén, personne) mal fermé ressemble à \zh{入} (rù, entrer). Un trait de trop transforme \zh{大} (dà, grand) en \zh{太} (tài, trop) ou en \zh{犬} (quǎn, chien). À l'examen, écris lentement plutôt que vite.
\end{attention}

\courssub{Sujet type examen}

\begin{exercice}[Sujet d'entraînement — 30 minutes]
\textbf{Sujet :} Écris un texte de 5 à 8 phrases pour présenter ta famille à ton correspondant chinois. Tu indiqueras : le nombre de personnes, les membres de la famille, le métier d'un parent, et un sentiment. Utilise au moins deux connecteurs. Soigne tes caractères. (Barème : caractères /5, grammaire /5, connecteurs /4, consigne /4, présentation /2.)
\end{exercice}

\begin{exemplebox}[Un texte modèle corrigé et noté 20/20]
我家有四口人：爸爸、妈妈、姐姐和我。爸爸是医生，妈妈很好。姐姐也喜欢足球。我们都爱我们的家。\\
Wǒ jiā yǒu sì kǒu rén : bàba, māma, jiějie hé wǒ. Bàba shì yīshēng, māma hěn hǎo. Jiějie yě xǐhuan zúqiú. Wǒmen dōu ài wǒmen de jiā.\\
« Ma famille compte quatre personnes : papa, maman, ma grande sœur et moi. Papa est médecin, maman est très gentille. Ma grande sœur aussi aime le football. Nous aimons tous notre famille. »\\
Pourquoi 20/20 ? Caractères exacts (5/5) ; structures correctes : \zh{有} + nombre + classificateur \zh{口} (le classificateur des membres de la famille), \zh{是} + nom de métier, adjectif sans \zh{是} (5/5) ; connecteurs \zh{和}, \zh{也}, \zh{都} (4/4) ; consigne respectée : 4 phrases courtes dans l'esprit de la consigne, tous les points demandés traités (4/4) ; présentation soignée (2/2).
\end{exemplebox}

\courssub{Correction collective : le texte d'un élève à corriger}

Voici le texte rédigé par un élève de Première A. Il contient \textbf{5 erreurs typiques}. Cherche-les avant de lire la correction.

\begin{exoresolu}[Trouve et corrige les 5 erreurs]
Énoncé — Texte de l'élève :\\
\zh{我家有五口人。爸爸是很好。妈妈和姐姐和我爱足球。我们家是大。我爱我的家。}\\
(Wǒ jiā yǒu wǔ kǒu rén. Bàba shì hěn hǎo. Māma hé jiějie hé wǒ ài zúqiú. Wǒmen jiā shì dà. Wǒ ài wǒ de jiā.)\\
« Ma famille compte cinq personnes. Papa est très gentil. Maman et ma sœur et moi aimons le football. Notre maison est grande. J'aime ma famille. »
\tcblower
Solution détaillée — les 5 erreurs :
\begin{enumerate}
\item \textbf{\zh{爸爸是很好}} : faute de grammaire. En chinois, un adjectif seul fait office de verbe (adjectif-verbe) ; \zh{是} ne se met jamais devant un adjectif. Correction : \zh{爸爸很好。} (Bàba hěn hǎo. « Papa est très gentil. »)
\item \textbf{\zh{妈妈和姐姐和我}} : \zh{和} ne relie pas trois éléments en chaîne comme le « et » français ; dans une énumération, on utilise la virgule chinoise 、 et on place \zh{和} seulement avant le dernier élément. Correction : \zh{妈妈、姐姐和我}.
\item \textbf{\zh{我们家是大}} : \zh{是} devant l'adjectif \zh{大}. Correction : \zh{我们的家很大。} (Wǒmen de jiā hěn dà. « Notre maison est grande. »)
\item \textbf{\zh{我们家}} : oubli du possessif \zh{的} ; dans un style soigné, on écrit \zh{我们的家} (à l'oral rapide, \zh{我们家} se dit, mais à l'écrit scolaire on garde \zh{的}).
\item \textbf{Caractère} : l'élève a écrit \zh{爱} (ài, aimer) de façon incomplète, ce qui le rend illisible. Attention à sa structure : en haut la griffe 爫 (c'est le radical du caractère), au milieu le couvercle 冖, en bas l'ami 友. Trace les trois parties bien nettes, de haut en bas.
\end{enumerate}
Texte corrigé : \zh{我家有五口人。爸爸很好。妈妈、姐姐和我都爱足球。我们的家很大。我爱我的家。} (Wǒ jiā yǒu wǔ kǒu rén. Bàba hěn hǎo. Māma, jiějie hé wǒ dōu ài zúqiú. Wǒmen de jiā hěn dà. Wǒ ài wǒ de jiā.) « Ma famille compte cinq personnes. Papa est très gentil. Maman, ma grande sœur et moi aimons tous le football. Notre maison est grande. J'aime ma famille. »
\end{exoresolu}

\begin{attention}[Les 5 fautes qui coûtent le plus cher à l'examen]
\begin{enumerate}
\item \zh{是} + adjectif (calque du français « il est grand ») → dis \zh{他很高} (Tā hěn gāo) ;
\item \zh{和} utilisé entre deux phrases ou entre deux verbes → \zh{和} ne relie que des noms ou des groupes nominaux ;
\item oubli de \zh{的} possessif : \zh{我姐姐} au lieu de \zh{我的姐姐} à l'écrit soigné ;
\item ordre des mots français : le complément de temps ou de lieu se place \emph{avant} le verbe : \zh{我们一起吃饭。} (Wǒmen yìqǐ chīfàn. « Nous mangeons ensemble. ») ;
\item caractère déformé ou inventé : un trait en trop ou en moins change le sens (\zh{大} / \zh{太}, \zh{人} / \zh{入}).
\end{enumerate}
\end{attention}

\courssub{Entraînement final et conseils de stratégie}

\begin{exercice}[À toi de jouer — conditions d'examen, 25 minutes]
Rédige sur ta copie un texte de 5 à 8 phrases intitulé \zh{我的家} (Wǒ de jiā, « Ma famille ») destiné à ton correspondant chinois de Yaoundé. Contraintes : au moins deux connecteurs parmi \zh{和}, \zh{也}, \zh{都}, \zh{因为} ; au moins un métier avec \zh{是} ; au moins un adjectif sans \zh{是}. Puis relis-toi avec la méthode C-O-N-N-E-C-T et auto-note-toi avec le barème sur 20.
\end{exercice}

\begin{corrige}[Exercice — éléments de correction attendus]
Un texte réussi contient par exemple : \zh{我家有三口人：爸爸、妈妈和我。爸爸是老师，妈妈也很好。我们都爱足球，因为喀麦隆的足球很有名。我爱我的家。} (Wǒ jiā yǒu sān kǒu rén : bàba, māma hé wǒ. Bàba shì lǎoshī, māma yě hěn hǎo. Wǒmen dōu ài zúqiú, yīnwèi Kāmàilóng de zúqiú hěn yǒumíng. Wǒ ài wǒ de jiā.) « Ma famille compte trois personnes : papa, maman et moi. Papa est professeur, maman aussi est très gentille. Nous aimons tous le football, car le football camerounais est très célèbre. J'aime ma famille. » Vérifie : classificateur \zh{口} après le nombre pour les membres de la famille, \zh{是} + métier, adjectifs sans \zh{是} (\zh{很好}, \zh{很有名}), trois connecteurs (\zh{和}, \zh{也}, \zh{因为}), \zh{的} partout où il le faut (\zh{我的家}, \zh{喀麦隆的足球}).
\end{corrige}

\begin{savaistu}[Mieux vaut court et juste que long et fautif]
À l'examen de chinois, un texte simple mais sans faute de caractères vaut mieux qu'un texte long plein d'erreurs : la précision graphique est un critère officiel de notation, au Cameroun comme dans les tests internationaux HSK. Les correcteurs chinois disent : \zh{写字要工整} (xiě zì yào gōngzhěng, « l'écriture doit être nette et régulière »). Un caractère bien tracé témoigne du respect de la langue et du lecteur — une valeur culturelle forte en Chine, où la calligraphie (\zh{书法}, shūfǎ) est un art millénaire. Petit détail émouvant : le caractère traditionnel de « aimer », 愛, contient encore le cœur 心 ; le caractère simplifié \zh{爱} l'a perdu, mais il a gardé l'ami 友 — aimer, c'est aussi être l'ami de l'autre.
\end{savaistu}

\begin{aretenir}[L'essentiel de la section]
\begin{itemize}
\item Barème sur 20 : caractères /5, grammaire /5, connecteurs /4, consigne /4, présentation /2.
\item Relis-toi avec la méthode \cle{C-O-N-N-E-C-T} : Caractères, Ordre, Négation, Nombre, 的, Connecteurs, Tons.
\item Les 5 fautes à bannir : \zh{是} + adjectif, \zh{和} entre phrases ou verbes, oubli de \zh{的}, ordre des mots français, caractères déformés.
\item Stratégie gagnante : un texte court, exact et soigné rapporte toujours plus qu'un texte long et fautif.
\end{itemize}
\end{aretenir}

\newpage

\coursec{Activité d'intégration}

\coursec{Leçon 4 — Expression écrite : écriture des caractères de l'amour et de la famille ; types de phrases étudiées}

\courssub{Objectifs de la leçon}

À la fin de cette leçon, tu seras capable de :

\begin{itemize}
\item produire un texte écrit simple et correct (une carte de vœux) sur le thème de \cle{l'amour} et de \cle{la famille} ;
\item écrire correctement les caractères du thème (\zh{爱}, \zh{家}, \zh{妈}, \zh{爸}, \zh{祝}, \zh{您}, \zh{身体}, \zh{快乐}) en respectant l'ordre des traits, les radicaux et en évitant les confusions graphiques ;
\item utiliser les \cle{types de phrases} étudiés : phrase déclarative avec \zh{是}, phrase de sentiment avec \zh{很爱}, phrase de cause avec \zh{因为}, phrase de vœux avec \zh{祝}, phrase exclamative, formule de politesse interrogative (\zh{您好}) ;
\item employer les \cle{connecteurs} simples du texte : \zh{和}, \zh{也}, \zh{因为}, \zh{所以}, \zh{祝} ;
\item traduire un court message du français vers le chinois (thème) et du chinois vers le français (version) ;
\item t'auto-évaluer à l'aide d'une grille de critères alignée sur les attendus de l'examen de Première.
\end{itemize}

\courssub{Situation de communication}

Tu es élève en classe de Première A au lycée de Yaoundé. À l'occasion de la fête des Mères, ton professeur de chinois propose à la classe un projet : chaque élève rédige une \cle{carte de vœux} en chinois pour un membre de sa famille. Ta correspondante chinoise, Li Na, t'a envoyé un modèle de carte qu'elle a écrite pour sa mère. Lis d'abord sa carte, puis analyse sa structure avant d'écrire la tienne.

\begin{textebox}[La carte de vœux de Li Na — Modèle authentique]
亲爱的妈妈：\\
Qīn'ài de māma :\\
« Chère maman : »\\[4pt]
您好！今天是母亲节。\\
Nín hǎo ! Jīntiān shì Mǔqīnjié.\\
« Bonjour ! Aujourd'hui, c'est la fête des Mères. »\\[4pt]
我很爱您，因为您每天照顾我。\\
Wǒ hěn ài nín, yīnwèi nín měitiān zhàogù wǒ.\\
« Je vous aime beaucoup, parce que vous prenez soin de moi chaque jour. »\\[4pt]
爸爸和我祝您身体健康，天天快乐！\\
Bàba hé wǒ zhù nín shēntǐ jiànkāng, tiāntiān kuàilè !\\
« Papa et moi vous souhaitons une bonne santé et du bonheur chaque jour ! »\\[4pt]
您的女儿：李娜\\
Nín de nǚ'ér : Lǐ Nà\\
« Votre fille : Li Na »
\end{textebox}

\courssub{Activité 1 : Compréhension et observation guidée}

\begin{enumerate}
\item \textbf{Compréhension globale.} Réponds en français : À qui Li Na écrit-elle ? Pour quelle occasion ? Quels sont les deux vœux formulés à la fin ?
\item \textbf{Repérage structurel.} Dans le texte, identifie et surligne (ou recopie) :
      \begin{itemize}
      \item La formule d'appel (appellation).
      \item Le connecteur de cause.
      \item Le verbe qui introduit les vœux.
      \item La formule de signature.
      \end{itemize}
\item \textbf{Analyse grammaticale.} Quel est le type de phrase de \zh{今天是母亲节} ? Justifie en nommant la structure utilisée.
\item \textbf{Observation culturelle.} Pourquoi utilise-t-on \zh{您} (nín) au lieu de \zh{你} (nǐ) dans cette carte ? À qui d'autre pourrait-on l'écrire ?
\end{enumerate}

\courssub{Activité 2 : Écriture des caractères — Traits, radicaux et ordre d'écriture}

\begin{methode}[Rappel des règles générales d'ordre des traits]
\begin{enumerate}
\item De haut en haut (\zh{三} sān).
\item De gauche à droite (\zh{川} chuān).
\item Horizontal avant vertical (\zh{十} shí).
\item Diagonale droite avant diagonale gauche (\zh{人} rén).
\item L'extérieur avant l'intérieur (\zh{月} yuè, \zh{国} guó).
\item Fermer en dernier (\zh{口} kǒu, \zh{田} tián).
\item Le trait central avant les deux côtés (\zh{小} xiǎo, \zh{水} shuǐ).
\end{enumerate}
\end{methode}

\begin{propriete}[Analyse graphique des caractères clés de la leçon]
\begin{center}\begin{tabularx}{\linewidth}{|X|X|X|X|X|}\hline
\rowcolor{popblueL} Caractère & Pinyin & Radical (Clé) & Nombre de traits & Décomposition / Ordre principal \\ \hline
\zh{爱} & ài & \zh{爪} (zhuǎ, griffe) & 10 & \zh{爪} + \zh{冖} + \zh{友} (haut → bas ; \zh{友} : haut \zh{又} + bas \zh{一}) \\ \hline
\zh{家} & jiā & \zh{宀} (mián, toit) & 10 & \zh{宀} (toit : 2 traits) + \zh{豕} (shǐ, cochon : 7 traits) → écrire le toit d'abord \\ \hline
\zh{妈} & mā & \zh{女} (nǚ, femme) & 13 & \zh{女} (gauche, 3 traits) + \zh{马} (droite, 10 traits) → gauche avant droite \\ \hline
\zh{爸} & bà & \zh{父} (fù, père) & 8 & \zh{父} (gauche, 4 traits) + \zh{巴} (droite, 4 traits) → gauche avant droite \\ \hline
\zh{祝} & zhù & \zh{礻} (shì, autel/culte) & 9 & \zh{礻} (gauche, 4 traits : point, horizontale, verticale, point) + \zh{兄} (xiōng, droite, 5 traits) \\ \hline
\zh{您} & nín & \zh{心} (xīn, cœur) & 11 & \zh{你} (nǐ, haut, 7 traits) + \zh{心} (bas, 3 traits + point) → haut avant bas \\ \hline
\zh{身} & shēn & \zh{身} (shēn, corps) & 7 & Ordre : \zh{ノ} \zh{一} \zh{一} \zh{一} \zh{丨} \zh{一} \zh{一} (traits horizontaux groupés) \\ \hline
\zh{体} & tǐ & \zh{骨} (gǔ, os) simplifié \zh{亻} & 7 & \zh{亻} (rén, 2 traits) + \zh{本} (běn, 5 traits) → gauche avant droite \\ \hline
\zh{快} & kuài & \zh{忄} (xīn, cœur vertical) & 7 & \zh{忄} (3 traits) + \zh{夬} (guài, 4 traits) → gauche avant droite \\ \hline
\zh{乐} & lè & \zh{丿} (piě) & 5 & \zh{丿} \zh{一} \zh{一} \zh{丨} \zh{一} (structure haut-bas) \\ \hline
\end{tabularx}\end{center}
\end{propriete}

\begin{experience}[Atelier d'écriture guidée]
\begin{enumerate}
\item Sur ton cahier d'écriture (格子纸 \zh{gézǐzhǐ}), recopie \textbf{trois fois} chaque caractère du tableau ci-dessus en respectant scrupuleusement l'ordre des traits.
\item Pour \zh{爱}, \zh{家}, \zh{祝}, \zh{您} : entoure le radical au crayon de couleur.
\item Échange ton cahier avec ton voisin : vérifie mutuellement l'ordre des traits (numérote les traits sur un exemplaire) et la proportion des composants (le radical ne doit pas écraser le phonétique).
\end{enumerate}
\end{experience}

\courssub{Activité 3 : Caractères proches — Ne pas confondre}

\begin{attention}[Pièges graphiques et phonétiques fréquents (Francophones)]
\begin{center}\begin{tabularx}{\linewidth}{|X|X|X|X|}\hline
\rowcolor{poporangeL} Caractère cible & Caractère proche & Différence clé & Astuce mnémotechnique \\ \hline
\zh{爱} (ài, aimer) & \zh{受} (shòu, recevoir) & Bas : \zh{友} (ami) vs \zh{又} (de nouveau) & \textbf{Aimer}, c'avoir un \zh{ami} (\zh{友}) sous son toit. \\ \hline
\zh{家} (jiā, famille) & \zh{加} (jiā, ajouter) & Radical : \zh{宀} (toit) vs \zh{力} (force) & Une \textbf{famille} vit sous un \textbf{toit}. \\ \hline
\zh{妈} (mā, mère) & \zh{马} (mǎ, cheval) & Radical : \zh{女} (femme) vs aucun & \textbf{Maman} est une \textbf{femme} (\zh{女}) qui galope comme un cheval (phonétique). \\ \hline
\zh{爸} (bà, père) & \zh{巴} (bā, s'accrocher) & Radical : \zh{父} (père) vs aucun & \textbf{Papa} est le \textbf{père} (\zh{父}) de la famille. \\ \hline
\zh{祝} (zhù, souhaiter) & \zh{福} (fú, bonheur) / \zh{禄} (lù, salaire) & Phonétique : \zh{兄} vs \zh{畐}/\zh{录} & Même radical \zh{礻} (culte). \textbf{Souhaiter} (\zh{祝}) rime avec \zh{兄} (frère aîné). \\ \hline
\zh{您} (nín, vous poli) & \zh{你} (nǐ, toi) & Bas : \zh{心} (cœur) vs \zh{人} (homme) & \textbf{Vous} (poli) vient du \textbf{cœur} (\zh{心}). \\ \hline
\zh{身} (shēn, corps) & \zh{申} (shēn, étendre) & Trait central : \zh{身} a 3 horizontales, \zh{申} a 1 & Le \textbf{corps} a de la \textbf{largeur} (trois traits). \\ \hline
\zh{体} (tǐ, corps) & \zh{本} (běn, racine) & Radical : \zh{亻} (homme) vs aucun & Le \textbf{corps} (\zh{体}) appartient à un \textbf{être humain} (\zh{亻}). \\ \hline
\end{tabularx}\end{center}
\end{attention}

\courssub{Activité 4 : Structure du texte et connecteurs}

\begin{propriete}[Les 5 structures-types de la carte de vœux]
\begin{center}\begin{tabularx}{\linewidth}{|X|X|X|}\hline
\rowcolor{popgreenL} Fonction & Structure schématique & Exemple du texte (Caractères + Pinyin) \\ \hline
Appel / Apostrophe & \zh{亲爱的} + Titre/Prénom + \zh{：} & \zh{亲爱的妈妈：} (Qīn'ài de māma :) \\ \hline
Salutation polie & \zh{您好} ! & \zh{您好！} (Nín hǎo !) \\ \hline
Identification / Occasion & Sujet + \zh{是} + Nom/Date & \zh{今天是母亲节。} (Jīntiān shì Mǔqīnjié.) \\ \hline
Sentiment + Cause & Sujet + \zh{很爱} + Objet + \zh{，因为} + Cause & \zh{我很爱您，因为您每天照顾我。} \\ \hline
Vœux (Formule figée) & \zh{祝} + Bénéficiaire + Vœu 1 + \zh{，} + Vœu 2 + \zh{！} & \zh{祝您身体健康，天天快乐！} \\ \hline
Signature & \zh{您的} + Lien de parenté + \zh{：} + Prénom & \zh{您的女儿：李娜} \\ \hline
\end{tabularx}\end{center}
\end{propriete}

\begin{aretenir}[Connecteurs logiques utiles pour l'écrit simple]
\begin{itemize}
\item \zh{和} (hé) : coordination de noms/pronoms (\zh{爸爸和我}). Ne relie pas des phrases complètes.
\item \zh{也} (yě) : « aussi », placé \textbf{avant} le verbe (\zh{我也爱您}).
\item \zh{因为} (yīnwèi) : introduit la cause. Souvent en position médiane en chinois oral/écrit simple (\zh{Resultat, 因为 Cause}), ou initiale en écrit formel (\zh{因为 Cause, 所以 Résultat}).
\item \zh{所以} (suǒyǐ) : introduit la conséquence. S'associe souvent à \zh{因为} en écrit formel.
\item \zh{祝} (zhù) : verbe « souhaiter ». \textbf{Attention} : pas de \zh{很} (hěn) ni de \zh{的} (de) après \zh{祝} + personne. Structure : \zh{祝} + [Personne] + [Adjectif/Verbe à valeur de vœu].
\end{itemize}
\end{aretenir}

\courssub{Activité 5 : Thème — Français → Chinois}

\begin{exercice}[Entraînement au thème (phrases isolées)]
Traduis les phrases suivantes en chinois (caractères + pinyin). Utilise le vocabulaire de la leçon.
\begin{enumerate}
\item J'aime beaucoup ma famille.
\item Ma mère et mon père sont très gentils (bons).
\item Je vous souhaite une bonne santé.
\item Aujourd'hui, c'est l'anniversaire de mon père.
\item Parce que ma mère cuisine bien, je suis heureux.
\end{enumerate}
\end{exercice}

\begin{corrige}[Exercice 5 — Thème]
\begin{enumerate}
\item \zh{我很爱我的家人。} (Wǒ hěn ài wǒ de jiārén.)
\item \zh{我的爸爸妈妈很好。} (Wǒ de bàba māma hěn hǎo.) \emph{ou} \zh{我的父母很慈祥。} (Wǒ de fùmǔ hěn cíxiáng — niveau HSK 4).
\item \zh{祝您身体健康。} (Zhù nín shēntǐ jiànkāng.)
\item \zh{今天是我爸爸的生日。} (Jīntiān shì wǒ bàba de shēngrì.)
\item \zh{因为妈妈菜做得好，所以我很高兴。} (Yīnwèi māma cài zuò de hǎo, suǒyǐ wǒ hěn gāoxìng.) \emph{Structure formelle 因为...所以...}
\end{enumerate}
\end{corrige}

\courssub{Activité 6 : Version — Chinois → Français}

\begin{exercice}[Entraînement à la version]
Traduis en français les phrases ou extraits suivants :
\begin{enumerate}
\item \zh{亲爱的奶奶，您好！}
\item \zh{今天是教师节。}
\item \zh{我们全家都很爱您。}
\item \zh{因为您辛苦工作，所以我们祝您快乐。}
\item \zh{祝您生日快乐，万事如意！}
\end{enumerate}
\end{exercice}

\begin{corrige}[Exercice 6 — Version]
\begin{enumerate}
\item « Chère grand-mère, bonjour ! »
\item « Aujourd'hui, c'est la fête des Professeurs. »
\item « Toute notre famille vous aime beaucoup. » (\zh{全家} = toute la famille / tout le foyer).
\item « Parce que vous travaillez dur, nous vous souhaitons du bonheur. » (\zh{辛苦工作} = travailler dur / péniblement).
\item « Joyeux anniversaire, que tout se passe comme vous le souhaitez ! » (\zh{万事如意} wànshì rúyì = tout aller bien / tout selon vos vœux, formule idiomatique courante).
\end{enumerate}
\end{corrige}

\courssub{Activité 7 : Production écrite — Ma carte de vœux}

\begin{methode}[Étapes pour rédiger ta carte de vœux (6 à 8 lignes)]
\begin{enumerate}
\item \textbf{Choisis le destinataire et l'occasion :} Mère (fête des Mères), Père (fête des Pères), Grand-mère/-père (anniversaire), Oncle/Tante...
\item \textbf{Rédige l'appel :} \zh{亲爱的} + [Titre] + \zh{：} (ex: \zh{亲爱的爸爸：}).
\item \textbf{Salue :} \zh{您好！}
\item \textbf{Situe l'occasion :} \zh{今天是} + [Occasion] + \zh{。} (ex: \zh{今天是父亲节。} / \zh{今天是您的生日。}).
\item \textbf{Exprime ton sentiment + cause :} \zh{我很爱您，因为} + [Raison concrète] + \zh{。} (ex: \zh{因为您陪我玩足球} / \zh{因为您给我做好吃的}).
\item \textbf{Formule les vœux :} \zh{祝您} + [Vœu 1] + \zh{，} + [Vœu 2] + \zh{！} (Classiques : \zh{身体健康}, \zh{天天快乐}, \zh{万事如意}, \zh{生日快乐}).
\item \textbf{Signe :} \zh{您的} + [Ton lien de parenté] + \zh{：} + [Ton prénom en chinois] (ex: \zh{您的儿子：保罗}).
\item \textbf{Relis-toi :} Vérifie les caractères (radicaux, traits), les tons (pinyin), la ponctuation chinoise (， 。 ！ ：) et les connecteurs.
\end{enumerate}
\end{methode}

\begin{exoresolu}[Production modèle : Carte pour mon père (Fête des Pères)]
Énoncé : Rédige une carte pour ton père à l'occasion de la fête des Pères.
\tcblower
\textbf{Production :}\\[4pt]
亲爱的爸爸：\\
Qīn'ài de bàba :\\
« Cher papa : »\\[4pt]
您好！今天是父亲节。\\
Nín hǎo ! Jīntiān shì Fùqīnjié.\\
« Bonjour ! Aujourd'hui, c'est la fête des Pères. »\\[4pt]
我很爱您，因为您常常陪我踢足球。\\
Wǒ hěn ài nín, yīnwèi nín chángcháng péi wǒ tī zúqiú.\\
« Je vous aime beaucoup, parce que vous jouez souvent au foot avec moi. »\\[4pt]
妈妈和我祝您身体健康，万事如意！\\
Māma hé wǒ zhù nín shēntǐ jiànkāng, wànshì rúyì !\\
« Maman et moi vous souhaitons une bonne santé et que tout vous réussisse ! »\\[4pt]
您的儿子：保罗\\
Nín de érzi : Bǎoluó\\
« Votre fils : Paul »\\[6pt]
\textbf{Analyse linguistique :}
\begin{itemize}
\item \zh{陪} (péi) + activité : « accompagner pour... », structure HSK 3 très utile.
\item \zh{踢足球} (tī zúqiú) : verbe + objet, vocabulaire culturellement ancré au Cameroun.
\item \zh{万事如意} (wànshì rúyì) : chengyu (idiome) de vœu standard, niveau HSK 4, valorisant à l'examen.
\item Respect strict de la ponctuation pleine chasse et de la mise en page épistolaire.
\end{itemize}
\end{exoresolu}

\courssub{Activité 8 : Grille d'auto-évaluation et entraînement à l'examen}

\begin{aretenir}[Grille d'évaluation de la carte de vœux (20 points — Type Bac/Évaluation continue)]
\begin{center}\begin{tabularx}{\linewidth}{|X|c|X|}\hline
\rowcolor{poppurpleL} Critère & Points & Indicateurs de réussite \\ \hline
\textbf{Respect du genre épistolaire} & 4 & Appel (\zh{亲爱的...：}), Salutation (\zh{您好}), Signature (\zh{您的...：}) présents et corrects. Ponctuation chinoise. \\ \hline
\textbf{Structures grammaticales ciblées} & 6 & Phrase \zh{是} (occasion) ; Phrase \zh{很爱...因为...} (sentiment+cause) ; Phrase \zh{祝您...} (vœux). Pas de \zh{很} après \zh{祝}. \\ \hline
\textbf{Vocabulaire du thème} & 4 & Utilisation correcte de \zh{家/家人}, \zh{爱}, \zh{身体健康}, \zh{快乐/万事如意}, titres de parenté (\zh{爸爸/妈妈/奶奶/爷爷}). \\ \hline
\textbf{Écriture des caractères} & 3 & Caractères clés (\zh{爱, 家, 祝, 您, 身, 体, 快, 乐}) lisibles, traits dans l'ordre, radicaux respectés, pas de confusion (ex: \zh{爱} vs \zh{受}). \\ \hline
\textbf{Cohérence et correction langue} & 3 & Pinyin correct si fourni, ordre des mots (Sujet-Temps-Lieu-Verbe), pas de fautes de mots-outils (\zh{的/得/地} non requis ici mais \zh{和} vs \zh{与} géré). \\ \hline
\textbf{Total} & \textbf{20} & \\ \hline
\end{tabularx}\end{center}
\end{aretenir}

\begin{exercice}[Sujet d'entraînement type Examen (Durée : 30 min)]
\textbf{Consigne :} Tu t'appelles \zh{玛丽} (Mǎlì) ou \zh{保罗} (Bǎoluó). Écris une carte de vœux (7-8 lignes) pour ta \textbf{grand-mère maternelle} (\zh{外婆} wàipó) à l'occasion de \textbf{son 70e anniversaire} (\zh{七十岁大寿} qīshí suì dàshòu).
Ta carte doit impérativement contenir :
\begin{enumerate}
\item La formule d'appel et la salutation.
\item La mention de l'occasion (70 ans).
\item Une phrase avec \zh{因为} justifiant ton affection (ex: elle raconte des histoires, elle cuisine le ndolé, elle t'a appris le chinois...).
\item Deux vœux distincts introduits par \zh{祝}.
\item La signature avec ton lien de parenté (\zh{外孙} wàisūn / \zh{外孙女} wàisūnnǚ).
\end{enumerate}
\textit{Écris directement les caractères. Le pinyin n'est pas obligatoire mais recommandé pour la révision.}
\end{exercice}

\begin{corrige}[Exercice 8 — Proposition de corrigé type examen]
\textbf{Production attendue (exemple) :}\\[4pt]
亲爱的外婆：\\
Qīn'ài de wàipó :\\
您好！今天是您七十岁大寿。\\
Nín hǎo ! Jīntiān shì nín qīshí suì dàshòu.\\
我很爱您，因为您常给我讲中国故事，还教我做饺子。\\
Wǒ hěn ài nín, yīnwèi nín cháng gěi wǒ jiǎng Zhōngguó gùshì, hái jiào wǒ zuò jiǎozi.\\
全家人祝您福如东海长流水，寿比南山不老松！\\
Quán jiārén zhù nín fú rú dōnghǎi cháng liúshuǐ, shòu bǐ nánshān bù lǎo sōng !\\
您的外孙女：玛丽\\
Nín de wàisūnnǚ : Mǎlì\\[6pt]
\textbf{Notes pédagogiques :}
\begin{itemize}
\item \zh{外婆} (wàipó) = grand-mère maternelle (spécifique chinois, diff. de \zh{奶奶} nǎinai paternelle).
\item \zh{大寿} (dàshòu) : anniversaire important (généralement multiples de 10 ans dès 60/70 ans).
\item \zh{福如东海长流水，寿比南山不老松} : Chengyu (idiome) classique pour les vœux de longévité. L'élève peut utiliser \zh{祝您身体健康，长命百岁} (zhù nín shēntǐ jiànkāng, chángmìng bǎisuì) plus simple.
\item \zh{外孙女} / \zh{外孙} : termes précis pour petits-enfants par la fille. Valorise la précision lexicale.
\end{itemize}
\end{corrige}

\courssub{Bilan et perspectives}

Cette leçon t'a permis de passer de la \textbf{reconnaissance} (lecture du modèle) à la \textbf{production autonome} (écriture de ta carte), en consolidant trois piliers du programme de Première A :
\begin{enumerate}
\item \textbf{La graphie :} Maîtrise de l'ordre des traits, identification des radicaux (\zh{宀}, \zh{女}, \zh{父}, \zh{礻}, \zh{心}, \zh{忄}) et discrimination des caractères proches (\zh{家/加}, \zh{妈/马}, \zh{祝/福}, \zh{您/你}).
\item \textbf{La syntaxe de l'écrit fonctionnel :} Maîtrise du schéma fixe de la carte de vœux (Appel → Salutation → Occasion \zh{是} → Sentiment \zh{因为} → Vœux \zh{祝} → Signature \zh{您的}). Ce schéma est transférable aux SMS, emails, cartes de vœux électroniques (WeChat).
\item \textbf{La traduction thème/version :} Automatisation des structures \zh{祝} + O + Vœu et \zh{因为}..., \zh{所以}... / \zh{...，因为...}.
\end{enumerate}

\begin{savaistu}[Le ndolé et les jiǎozi : un pont culinaire]
Au Cameroun, le \textbf{ndolé} (feuilles amères, arachides, viande/crevettes) est le plat de fête par excellence, symbole de partage familial. En Chine du Nord, les \textbf{jiǎozi} (\zh{饺子}, raviolis) sont incontournables pour le Nouvel An chinois (春节 Chūnjié) et les grands anniversaires : leur forme rappelle les lingots d'or (yuanbao), symbole de prospérité. Les faire \textbf{ensemble} en famille (\zh{包饺子} bāo jiǎozi) est un moment de transmission intergénérationnelle, tout comme la préparation du ndolé par les tantes et grand-mères. Dans ta carte, mentionner \zh{教我做饺子} (m'apprendre à faire des raviolis) ou \zh{做好吃的菜} (cuisiner de bons plats) ancre ton chinois dans une réalité vécue des deux côtés.
\end{savaistu}

\begin{methode}[Pour aller plus loin — Prolongements possibles]
\begin{itemize}
\item \textbf{Oral :} Enregistre-toi en lisant ta carte (attention aux tons : \zh{ài} 4e ton, \zh{jiā} 1er ton, \zh{zhù} 4e ton, \zh{nín} 2e ton, \zh{kuàilè} 4e-4e ton). Envoie l'audio à ton professeur ou à un correspondant.
\item \textbf{Culture :} Cherche les dates de la Fête des Mères (2e dimanche de mai en Chine comme au Cameroun) et de la Fête des Pères (3e dimanche de juin). Comment dit-on « Bonne fête Papa » ? (\zh{父亲节快乐} Fùqīnjié kuàilè).
\item \textbf{Écriture créative :} Transform ta carte en message WeChat (ajoute des émojis \zh{❤}, formule plus courte : \zh{妈，节日快乐！祝您健康长寿！爱您！}).
\end{itemize}
\end{methode}

\newpage

\coursec{Méthodes et exercices résolus}

\coursec{Leçon 4 — Expression écrite : écriture des caractères de l'amour et de la famille ; types de phrases étudiées}

\courssub{Modèle de texte commenté et vocabulaire clé}

\begin{textebox}[Modèle : « Ma famille à Yaoundé »]
\zh{我家有五口人。} \\
\zh{爸爸是老师，妈妈是商人。} \\
\zh{我有一个哥哥和一个妹妹。} \\
\zh{因为我们相亲相爱，所以我们的家很幸福。} \\
\zh{我爱我的家。}
\tcblower
\textbf{Pinyin :} Wǒ jiā yǒu wǔ kǒu rén. Bàba shì lǎoshī, māma shì shāngrén. Wǒ yǒu yí ge gēge hé yí ge mèimei. Yīnwèi wǒmen xiāngqīn-xiāng'ài, suǒyǐ wǒmen de jiā hěn xìngfú. Wǒ ài wǒ de jiā. \\
\textbf{Traduction :} Ma famille compte cinq personnes. Papa est professeur, maman est commerçante. J'ai un grand frère et une petite sœur. Parce que nous nous aimons mutuellement, notre famille est très heureuse. J'aime ma famille.
\end{textebox}

\begin{center}
\begin{tabularx}{\linewidth}{|X|X|X|X|}
\hline
\rowcolor{popblueL} \textbf{Caractères} & \textbf{Pinyin} & \textbf{Nature} & \textbf{Traduction} \\ \hline
爱 & ài & verbe / nom & aimer / amour \\ \hline
家 & jiā & nom & famille, maison, foyer \\ \hline
爸 & bà & nom & papa \\ \hline
妈 & mā & nom & maman \\ \hline
哥 & gē & nom & grand frère \\ \hline
姐 & jiě & nom & grande sœur \\ \hline
弟 & dì & nom & petit frère \\ \hline
妹 & mèi & nom & petite sœur \\ \hline
亲 & qīn & adj. / verbe & proche, intime / embrasser \\ \hline
子 & zǐ & nom / suffixe & enfant, fils, suffixe nominal \\ \hline
口 & kǒu & classif. & classif. pour membres de la famille \\ \hline
个 & gè & classif. & classif. général \\ \hline
老师 & lǎoshī & nom & professeur \\ \hline
商人 & shāngrén & nom & commerçant \\ \hline
学生 & xuésheng & nom & étudiant \\ \hline
相亲相爱 & xiāngqīn-xiāng'ài & verbe (idiome) & s'aimer mutuellement \\ \hline
幸福 & xìngfú & adjectif & heureux, bonheur \\ \hline
雅温得 & Yǎwēndé & nom propre & Yaoundé \\ \hline
\end{tabularx}
\end{center}

\begin{methode}[Analyser le texte modèle : 4 points de repère]
\begin{enumerate}
\item \textbf{Structure globale} : 5 phrases = 1 présentation (nombre), 2 description (métiers, fratrie), 3 explication (cause/conséquence), 4 sentiment (amour).
\item \textbf{Mots-outils repérés} : \cle{有} (possession), \cle{是} (identité), \cle{和} (coordination), \cle{因为…所以…} (cause/conséquence), \cle{的} (possession/attribut), \cle{很} (liaison adjectif).
\item \textbf{Classificateurs} : \cle{口} avec \cle{人} pour le noyau familial (\zh{五口人}) ; \cle{个} pour les individus (\zh{一个哥哥}).
\item \textbf{Ordre des mots} : Sujet + Lieu/Temps + Verbe + Complément. Ex. : \zh{我 (S) + 在雅温得 (Lieu) + 爱 (V) + 我的家 (C).} — Le lieu \emph{précède} le verbe.
\end{enumerate}
\end{methode}

\begin{popfigure}
\begin{tikzpicture}[node distance=1.2cm, every node/.style={font=\small, align=center}]
\node (S) [draw, rounded corners, fill=popblueL, minimum width=2cm] {Sujet};
\node (TP) [draw, rounded corners, fill=poporangeL, minimum width=2.5cm, right=of S] {Temps / Lieu};
\node (V) [draw, rounded corners, fill=popgreenL, minimum width=1.5cm, right=of TP] {Verbe};
\node (O) [draw, rounded corners, fill=poppinkL, minimum width=2.5cm, right=of V] {Complément};
\draw[-Stealth, thick, popdark] (S) -- (TP);
\draw[-Stealth, thick, popdark] (TP) -- (V);
\draw[-Stealth, thick, popdark] (V) -- (O);
\node[below=0.8cm of V, text width=12cm, align=center] (ex) {\textbf{Exemple :} \zh{我} (Sujet) \zh{在雅温得} (Lieu) \zh{爱} (Verbe) \zh{ma famille} (Complément).};
\end{tikzpicture}
\legende{Schéma de l'ordre des mots chinois : Sujet + Temps/Lieu + Verbe + Complément}
\end{popfigure}

\courssub{Écriture des caractères : traits, radicaux et caractères proches}

\begin{methode}[Écrire un caractère sans faute : la démarche en 5 étapes]
Pour écrire correctement les caractères du thème (\zh{爱, 家, 爸, 妈, 哥, 姐, 弟, 妹, 亲, 子}), suis toujours la même démarche :
\begin{enumerate}
\item \textbf{Identifier le radical (clé)} : il donne le sens général. Ex. : \zh{女} « femme » dans \zh{妈, 姐, 妹} ; \zh{宀} « toit » dans \zh{家} ; \zh{父} « père » dans \zh{爸} ; \zh{弓} « arc » dans \zh{弟} ; \zh{见} « voir » dans \zh{亲}.
\item \textbf{Décomposer le caractère} en parties (haut/bas, gauche/droite, extérieur/intérieur). Ex. : \zh{家} = \zh{宀} (toit) + \zh{豕} (cochon) ; \zh{爱} = \zh{爫} (griffe) + \zh{冖} (couvercle) + \zh{友} (ami) ; \zh{妈} = \zh{女} + \zh{马}.
\item \textbf{Respecter l'ordre des traits} (règles d'or) : \textbf{1.} De haut en bas (\zh{三}) ; \textbf{2.} De gauche à droite (\zh{川}) ; \textbf{3.} Horizontal avant vertical (\zh{十}) ; \textbf{4.} Extérieur avant intérieur (\zh{同}) ; \textbf{5.} On « ferme » la boîte en dernier (\zh{口, 田, 国}).
\item \textbf{Compter les traits} et vérifier les jonctions/croisements (source principale d'erreurs).
\item \textbf{Réécrire 5 fois} en prononçant le pinyin avec le ton à voix haute, puis écrire le caractère \emph{dans un mot} (\zh{妈妈, 家人, 爱人, 父母}) jamais isolément.
\end{enumerate}
\cle{Réflexe d'examen} : si tu hésites entre deux caractères proches, écris le mot complet : le contexte (\zh{妈} dans \zh{妈妈}) sécurise ta copie.
\end{methode}

\begin{exoresolu}[Dictée de caractères et analyse graphique]
Énoncé — 1) Écris en caractères : \emph{māma}, \emph{jiā}, \emph{gēge}, \emph{ài}. 2) Pour chaque caractère de \zh{家} et \zh{妈}, indique le radical et le nombre total de traits. 3) Recopie \zh{爱} en respectant l'ordre des traits et décris cet ordre en français.
\tcblower
\textbf{Solution détaillée}
\begin{enumerate}
\item \emph{māma} → \zh{妈妈} ; \emph{jiā} → \zh{家} ; \emph{gēge} → \zh{哥哥} ; \emph{ài} → \zh{爱}. \textbf{Attention aux tons} : \emph{māma} (1\textsuperscript{er} ton + ton neutre), \emph{gēge} (1\textsuperscript{er} ton + ton neutre) : le deuxième caractère d'un nom de parent redoublé se lit au ton neutre, mais s'écrit exactement comme le premier.
\item \zh{家} : radical \zh{宀} (le toit, 3 traits), 10 traits au total. \zh{妈} : radical \zh{女} (la femme, 3 traits), 6 traits au total (\zh{女} 3 traits + \zh{马} 3 traits).
\item \zh{爱} (10 traits, radical \zh{爫} « griffe » en simplifié) — Ordre précis :
\begin{enumerate}
\item \zh{丶} \zh{丶} (deux points de la griffe \zh{爫}, gauche puis droite)
\item \zh{一} (trait horizontal de la griffe)
\item \zh{丨} (crochet vertical descendant de la griffe)
\item \zh{一} (horizontal du couvercle \zh{冖})
\item \zh{丨} (crochet vertical descendant du couvercle)
\item \zh{一} (horizontal haut de \zh{友})
\item \zh{丿} (piě de \zh{友})
\item \zh{一} (horizontal bas de \zh{友})
\item \zh{乀} (nà final de \zh{友} / \zh{又})
\end{enumerate}
\textbf{Mnémotechnique} : Sous le couvercle protecteur (\zh{冖}), un ami (\zh{友}) : c'est l'amour (\zh{爱}). Le cœur (\zh{心}) de la version traditionnelle \zh{愛} a été remplacé par \zh{友} en simplifié.
\end{enumerate}
\textbf{Conclusion} : une copie sans faute de traits et avec le bon radical rapporte des points faciles ; entraîne-toi quotidiennement sur les 10 caractères-clés du thème.
\end{exoresolu}

\begin{exemplebox}[Caractères proches : ne pas confondre !]
\begin{center}
\begin{tabularx}{\linewidth}{|X|X|X|X|}
\hline
\rowcolor{poporangeL} \textbf{Caractère cible} & \textbf{Pinyin} & \textbf{Confusion possible} & \textbf{Différence graphique / Sens} \\ \hline
\zh{妈} & mā & \zh{马} (mǎ, cheval) & \zh{妈} a le radical \zh{女} (femme) à gauche. \\ \hline
\zh{爱} & ài & \zh{受} (shòu, recevoir) & \zh{爱} a \zh{冖} + \zh{友} au centre ; \zh{受} a \zh{又} (main droite) sous \zh{冖}. \\ \hline
\zh{家} & jiā & \zh{安} (ān, paix) & \zh{家} a \zh{豕} (cochon) sous le toit ; \zh{安} a \zh{女} (femme) sous le toit. \\ \hline
\zh{姐} & jiě & \zh{且} (qiě, de plus) & \zh{姐} a le radical \zh{女} ; \zh{且} n'en a pas. \\ \hline
\zh{妹} & mèi & \zh{未} (wèi, pas encore) & \zh{妹} a \zh{女} + \zh{未} ; \zh{未} est \zh{木} + trait court (vs \zh{末} trait long). \\ \hline
\end{tabularx}
\end{center}
\end{exemplebox}

\courssub{Structure du texte, connecteurs et types de phrases}

\begin{methode}[Rédiger un paragraphe de présentation familiale]
Pour produire un texte simple, correct et cohérent sur ta famille :
\begin{enumerate}
\item \textbf{Planifie en 4-5 phrases} : (a) Présenter le foyer (\zh{我家有…口人}) ; (b) Présenter les membres + métiers/études (\zh{爸爸是…，妈妈是…}) ; (c) Fratrie (\zh{我有一个哥哥/妹妹}) ; (d) Sentiment / ambiance (\zh{我们相亲相爱 / 很幸福}) ; (e) Conclusion (\zh{我爱我的家}).
\item \textbf{Utilise les connecteurs logiques} :
\begin{itemize}
\item \cle{还有} / \cle{和} : coordination (« et », « et aussi ») — \zh{我有一个哥哥\cle{和}一个妹妹}.
\item \cle{也} / \cle{都} : inclusion (« aussi », « tous ») — placés \textbf{avant} le verbe : \zh{姐姐\cle{也}是学生} ; \zh{我们\cle{都}爱家}.
\item \cle{因为…所以…} : cause / conséquence — \zh{\cle{因为}我们相亲相爱，\cle{所以}家很幸福}.
\item \cle{但是} / \cle{可是} : opposition.
\end{itemize}
\item \textbf{Respecte l'ordre canonique} : \cle{Sujet + (Temps/Lieu) + Verbe + Complément}. L'adjectif attributif précède le nom (+ \cle{的}) : \zh{我的\cle{小}妹妹} ; l'adjectif prédicatif suit le sujet (+ \cle{很}) : \zh{妹妹\cle{很}小}.
\item \textbf{Classificateurs obligatoires} : \cle{口} pour compter la famille (\zh{四口人}) ; \cle{个} pour individus (\zh{一个妹妹}). Jamais de nombre nu devant le nom.
\item \textbf{Relis (check-list)} : Tons du pinyin ? Traits des caractères ? \cle{的} possessif (\zh{我的家}) ? \cle{很} devant adj. prédicatif ? Négation \cle{没} avec \cle{有} ?
\end{enumerate}
\end{methode}

\begin{exoresolu}[Composition guidée : « Ma famille »]
Énoncé — Rédige un texte de 5 phrases en chinois pour présenter ta famille à Yaoundé : nombre, métiers, fratrie, sentiment, conclusion. Utilise au moins \cle{两个} connecteurs parmi \zh{还有, 也, 都, 因为…所以…}.
\tcblower
\textbf{Solution détaillée}

\textbf{Étape 1 — Plan} : (a) 5 personnes ; (b) père prof, mère infirmière ; (c) un grand frère étudiant ; (d) amour mutuel → bonheur ; (e) déclaration d'amour.

\textbf{Étape 2 — Rédaction} :
\begin{itemize}
\item (a) \zh{我家有五口人。} \emph{Wǒ jiā yǒu wǔ kǒu rén.} (Classif. \cle{口} obligatoire).
\item (b) \zh{爸爸是老师，妈妈是护士。} \emph{Bàba shì lǎoshī, māma shì hùshi.} (\cle{是} + métier, sans article).
\item (c) \zh{我还有一个哥哥，他是学生。} \emph{Wǒ hái yǒu yí ge gēge, tā shì xuésheng.} (\cle{还} = aussi ; \cle{一} → \emph{yí} devant 4\textsuperscript{e} ton).
\item (d) \zh{因为我们都相亲相爱，所以我们的家很幸福。} \emph{Yīnwèi wǒmen dōu xiāngqīn-xiāng'ài, suǒyǐ wǒmen de jiā hěn xìngfú.} (\cle{都} avant verbe ; \cle{很} liaison).
\item (e) \zh{我非常爱我的家。} \emph{Wǒ fēicháng ài wǒ de jiā.} (Adverbe \cle{非常} « très » avant verbe).
\end{itemize}
\textbf{Texte final} : \zh{我家有五口人。爸爸是老师，妈妈是护士。我还有一个哥哥，他是学生。因为我们都相亲相爱，所以我们的家很幸福。我非常爱我的家。}
\end{exoresolu}

\begin{propriete}[Les 4 types de phrases de base (HSK 1-3)]
\begin{center}
\begin{tabularx}{\linewidth}{|X|X|X|}
\hline
\rowcolor{popgreenL} \textbf{Type} & \textbf{Structure / Marqueur} & \textbf{Exemple (Caractères + Pinyin + Traduction)} \\ \hline
\textbf{Déclarative} & S + V + C + 。 & \zh{我爱我的家。} \emph{Wǒ ài wǒ de jiā.} (J'aime ma famille.) \\ \hline
\textbf{Interrogative fermée} & Phrase déclarative + \cle{吗} ? & \zh{你爱你的家\cle{吗}？} \emph{Nǐ ài nǐ de jiā \cle{ma}?} (Aimes-tu ta famille ?) \\ \hline
\textbf{Interrogative ouverte} & Mot interrogatif \textbf{in situ} (place du mot demandé) & \zh{你家有\cle{几}\cle{口}人？} \emph{Nǐ jiā yǒu \cle{jǐ} \cle{kǒu} rén?} (Combien de personnes ?) \\
& (\cle{谁} sujet, \cle{什么} objet, \cle{几} nombre, \cle{哪儿} lieu) & \zh{\cle{谁}是你的哥哥？} \emph{\cle{Shéi} shì nǐ de gēge?} (Qui est ton grand frère ?) \\ \hline
\textbf{Négative} & \cle{不} + V (action/état présent/futur) & \zh{我\cle{不}爱足球。} \emph{Wǒ \cle{bù} ài zúqiú.} (Je n'aime pas le foot.) \\
& \cle{没(有)} + Nom (possession/expérience passée) & \zh{我\cle{没有}弟弟。} \emph{Wǒ \cle{méi(yǒu)} dìdi.} (Je n'ai pas de petit frère.) \\ \hline
\textbf{Exclamative} & \cle{太} + Adj + \cle{了} ! / Adj + \cle{啊} ! & \zh{\cle{太}幸福\cle{了}！} \emph{\cle{Tài} xìngfú \cle{le}!} (C'est trop heureux !) \\
& & \zh{好\cle{啊}！} \emph{Hǎo \cle{a}!} (C'est bien !) \\ \hline
\end{tabularx}
\end{center}
\cle{Règle de transformation} : Pour passer de déclarative à interrogative fermée → ajoute \cle{吗} en fin de phrase, \textbf{aucun autre changement} (pas d'inversion sujet/verbe). Pour répondre « non » → reprends le verbe avec \cle{不} (ou \cle{没} pour \cle{有}) : \zh{爱} → \zh{不爱} ; \zh{有} → \zh{没有}.
\end{propriete}

\begin{exoresolu}[Transformations de phrases : entraînement type Bac]
Énoncé — 1) Transforme en question fermée : \zh{你爱你的家人。} 2) Mets à la forme négative : \zh{我有两个哥哥。} 3) Pose la question dont la réponse est : \zh{我家有五口人。} 4) Réponds « non » à : \zh{你是学生吗？} 5) Transforme en exclamative : \zh{这张照片很美。}
\tcblower
\textbf{Solution détaillée}
\begin{enumerate}
\item \zh{你爱你的家人\cle{吗}？} \emph{Nǐ ài nǐ de jiārén \cle{ma}?} (Ajout simple de \cle{吗}).
\item \zh{我\cle{没有}两个哥哥。} \emph{Wǒ \cle{méiyǒu} liǎng ge gēge.} (Négation de \cle{有} = \cle{没有} ; on garde le classif. \cle{两个}). \textbf{Interdit} : \zh{不有}.
\item \zh{你家有\cle{几}\cle{口}人？} \emph{Nǐ jiā yǒu \cle{jǐ} \cle{kǒu} rén?} (Le mot interrogatif \cle{几} remplace le nombre \cle{五} ; le classif. \cle{口} reste).
\item \zh{不，我\cle{不是}学生。} \emph{Bù, wǒ \cle{bú} shì xuésheng.} (Reprise du verbe \cle{是} avec \cle{不} → \emph{bú} 2\textsuperscript{e} ton devant \cle{是} 4\textsuperscript{e} ton).
\item \zh{\cle{这张照片\cle{太}美\cle{了}！}} \emph{\cle{Zhè zhāng zhàopiàn \cle{tài} měi \cle{le}}!} (Structure \cle{太}…\cle{了} pour intensité forte).
\end{enumerate}
\end{exoresolu}

\courssub{Traduction : Thème (Français → Chinois) et Version (Chinois → Français)}

\begin{methode}[Réussir le THÈME (Fr → Zh) en 5 étapes]
\begin{enumerate}
\item \textbf{Repère le Sujet et le Verbe} français ; le verbe chinois \textbf{ne se conjugue pas} (invariant).
\item \textbf{Traduis dans l'ordre chinois} : S + (Temps/Lieu) + V + C. Ex. : « J'aime beaucoup ma mère » → \zh{我} + \zh{很} + \zh{爱} + \zh{我的妈妈}.
\item \textbf{Place les mots-outils} : \cle{的} (possession/attribut) ; \cle{也/都} \textbf{avant} le verbe ; \cle{很} devant adj. prédicatif (obligatoire même sans « très » : \zh{她\cle{很}好}).
\item \textbf{Choisis le classif.} : \cle{口} (famille), \cle{个} (général), \cle{位} (politesse). \cle{一} → \emph{yí} (4\textsuperscript{e} ton) / \emph{yì} (2\textsuperscript{e} ton) / \emph{yī} (isolé/1\textsuperscript{er}/3\textsuperscript{e} ton).
\item \textbf{Vérifie les pièges} : Négation \cle{有} = \cle{没有} (jamais \cle{不有}) ; « et » entre noms = \cle{和} / \cle{跟} (pas \cle{也}) ; Complément de lieu \textbf{avant} le verbe (\zh{在雅温得\cle{上学}}).
\end{enumerate}
\end{methode}

\begin{exoresolu}[Thème : 4 phrases types examen]
Énoncé — Traduis en chinois (caractères + pinyin) :
1) « Mes parents travaillent à Yaoundé. »
2) « Ma grande sœur n'est pas commerçante. »
3) « Nous avons trois personnes dans la famille. »
4) « Parce qu'il m'aime, je suis très heureux. »
\tcblower
\textbf{Solution détaillée}
\begin{enumerate}
\item Sujet : \zh{我的爸爸妈妈} / \zh{我父母} ; Lieu : \zh{在雅温得} ; Verbe : \zh{工作} / \zh{上班}.
→ \zh{我父母在雅温得工作。} \emph{Wǒ fùmǔ zài Yǎwēndé gōngzuò.}
\item Sujet : \zh{我的姐姐} ; Négation adjectif/nom : \cle{不} + \cle{是} ; Métier : \zh{商人}.
→ \zh{我的姐姐不是商人。} \emph{Wǒ de jiějie \cle{bú} shì shāngrén.} (\cle{不} → \emph{bú} devant \cle{是}).
\item Sujet : \zh{我家} / \zh{我们} ; Verbe : \cle{有} ; Nombre + Classif : \cle{三}\cle{口}人.
→ \zh{我家有三口人。} \emph{Wǒ jiā yǒu sān kǒu rén.} (Préférer \cle{我家} pour « ma famille »).
\item Cause : \cle{因为} \zh{他爱我} ; Conséquence : \cle{所以} \zh{我\cle{很}高兴/幸福}.
→ \zh{因为他爱我，所以我\cle{很}幸福。} \emph{Yīnwèi tā ài wǒ, suǒyǐ wǒ \cle{hěn} xìngfú.} (\cle{很} obligatoire).
\end{enumerate}
\end{exoresolu}

\begin{methode}[Réussir la VERSION (Zh → Fr) en 4 étapes]
\begin{enumerate}
\item \textbf{Segmente} la phrase chinoise en groupes syntaxiques (S / Adv / V / C).
\item \textbf{Identifie les mots-outils} : \cle{的} (poss. / relat.), \cle{很} (liaison, souvent non traduit « très »), \cle{都} (tous), \cle{也} (aussi), \cle{因为…所以…} (car… donc), \cle{了} (aspect accompli / changement d'état).
\item \textbf{Reconstruis le français} : Ajoute articles (\emph{le, un}), conjugue (temps, personne), accorde (genre, nombre). Le chinois n'a \textbf{ni genre, ni nombre, ni conjugaison} : c'est ton travail.
\item \textbf{Soigne le style} : Évite le mot-à-mot (\emph{« ma famille a quatre bouches personnes »} → \emph{« ma famille compte quatre personnes »}).
\end{enumerate}
\end{methode}

\begin{exoresolu}[Version : 3 phrases progressives]
Énoncé — Traduis en français naturel :
1) \zh{我的爸爸妈妈都在雅温得工作。}
2) \zh{因为妹妹很小，所以她不上学。}
3) \zh{我们全家人相亲相爱，日子过得很幸福。}
\tcblower
\textbf{Solution détaillée}
\begin{enumerate}
\item Segmentation : \zh{我的爸爸妈妈} (Mes parents) + \zh{都} (tous les deux) + \zh{在雅温得} (à Yaoundé) + \zh{工作} (travaillent).
→ \textbf{« Mes parents travaillent tous les deux à Yaoundé. »} (\cle{都} = « tous les deux » car deux personnes).
\item Segmentation : \zh{因为} (Parce que) + \zh{妹妹很小} (ma petite sœur est très jeune) + \zh{所以} (donc) + \zh{她不上学} (elle ne va pas à l'école).
→ \textbf{« Comme ma petite sœur est très jeune, elle ne va pas à l'école. »} (Français naturel : \cle{因为…所以…} → « Comme… », « donc » souvent omis).
\item Segmentation : \zh{我们全家人} (Toute notre famille / Nous tous de la famille) + \zh{相亲相爱} (s'aiment mutuellement) + \zh{日子过得很幸福} (les jours se passent très heureusement / mènent une vie heureuse).
→ \textbf{« Toute notre famille s'aime mutuellement et mène une vie très heureuse. »} (\zh{日子过得…} = « mener une vie… »).
\end{enumerate}
\end{exoresolu}

\courssub{Élément socioculturel : La famille au Cameroun et en Chine}

\begin{savaistu}[Famille, filiation et valeurs : regards croisés]
\begin{itemize}
\item \textbf{Nom de famille} : En Chine, le nom de famille (\zh{姓} \emph{xìng}) précède le prénom (\zh{名} \emph{míng}) : \zh{王伟} = M. Wang Wei. Au Cameroun, ordre « Prénom Nom » (héritage francophone/anglophone). En Chine, la femme garde son nom de naissance après mariage ; les enfants portent traditionnellement le nom du père.
\item \textbf{Taille de la famille} : Politique de l'enfant unique (1979-2015) → génération « enfant unique » (\zh{独生子女} \emph{dúshēng zǐnǚ}) ; depuis 2021 : politique des trois enfants (\zh{三孩政策}). Au Cameroun, familles nombreuses valorisées, structure souvent élargie (oncles, tantes, cousins cohabitent).
\item \textbf{Vocabulaire de la parenté} : Chinois \textbf{très précis} selon âge, sexe, lignée (paternelle/maternelle) : \zh{哥哥/弟弟} (frère aîné/cadet), \zh{舅舅} (oncle maternel), \zh{叔叔} (oncle paternel cadet), \zh{伯伯} (oncle paternel aîné). Français/Cameroun : « oncle », « frère » génériques.
\item \textbf{Valeur centrale} : \cle{孝顺} \emph{xiàoshùn} (piété filiale : respect, soin aux parents âgés) = pilier confucéen. Au Cameroun, respect des aînés (\cle{les anciens}) et solidarité clanique sont tout aussi forts, mais s'expriment différemment (rites, partage, oralité).
\item \textbf{Fêtes} : \zh{春节} (Nouvel An chinois) = retrouvailles familiales obligatoires (\zh{回家} \emph{huíjiā}), repas du réveillon (\zh{年夜饭}), enveloppes rouges (\zh{红包} \emph{hóngbāo}). Au Cameroun : Noël, fin d'année, fêtes traditionnelles (Ngondo, Medumba…) = moments de regroupement clanique.
\end{itemize}
\end{savaistu}

\courssub{Grille d'évaluation et entraînement à l'examen}

\begin{aretenir}[Grille d'évaluation expression écrite (sur 20 pts — type Bac Blanc)]
\begin{center}
\begin{tabularx}{\linewidth}{|X|X|X|}
\hline
\rowcolor{poppurpleL} \textbf{Critère} & \textbf{Indicateurs de réussite} & \textbf{Points} \\ \hline
\textbf{Caractères} & Écriture correcte (traits, radical, proportion), pas de confusion (妈/马, 爱/受) & 4 \\ \hline
\textbf{Vocabulaire} & Emploi pertinent des 10+ mots du thème (亲, 爱, 家, métiers, classif.) & 3 \\ \hline
\textbf{Grammaire / Structures} & Ordre S-T-L-V-C respecté ; \cle{的, 很, 也, 都, 因为…所以…} bien placés ; types de phrases variés & 5 \\ \hline
\textbf{Cohérence / Connecteurs} & Texte logique (présentation → description → sentiment), 2+ connecteurs utilisés & 3 \\ \hline
\textbf{Traduction (Thème/Version)} & Fidélité + correction syntaxique française/chinoise ; gestion mots-outils & 3 \\ \hline
\textbf{Soin / Ponctuation} & Ponctuation chinoise pleine chasse (，。？), pinyin toné si demandé, lisibilité & 2 \\ \hline
\end{tabularx}
\end{center}
\end{aretenir}

\begin{exercice}[Entraînement final — Durée : 30 min]
\textbf{Exercice 1 (Caractères — 4 pts)} Écris en caractères : 1) \emph{bàba} 2) \emph{jiějie} 3) \emph{dìdi} 4) \emph{xìngfú}. Pour \zh{家} et \zh{爱}, donne le radical et le nombre de traits.

\textbf{Exercice 2 (Thème — 6 pts)} Traduis en chinois (caractères + pinyin) :
a) Mon grand frère est professeur.
b) Je n'ai pas de grande sœur.
c) Ma famille compte quatre personnes.
d) Parce que j'aime ma mère, je suis heureux.

\textbf{Exercice 3 (Version — 4 pts)} Traduis en français :
a) \zh{我家有四口人，爸爸是医生，妈妈也是医生。}
b) \zh{因为我们相亲相爱，所以我们很幸福。}

\textbf{Exercice 4 (Grammaire / Transformation — 3 pts)}
a) Question fermée : \zh{妹妹爱看书。}
b) Négation : \zh{我们有三个孩子。}
c) Question ouverte (sur le nombre) : \zh{他家有六口人。}

\textbf{Exercice 5 (Rédaction — 3 pts)} Écris un paragraphe de 5 phrases sur \emph{ta} famille (réelle ou imaginaire) en utilisant : \cle{我家有…口人}, \cle{因为…所以…}, \cle{也} ou \cle{都}, un métier, un sentiment.
\end{exercice}

\begin{corrige}[Entraînement final]
\textbf{Ex. 1} 1) \zh{爸爸} 2) \zh{姐姐} 3) \zh{弟弟} 4) \zh{幸福}. \zh{家} : rad. \zh{宀}, 10 traits. \zh{爱} : rad. \zh{爫}, 10 traits.

\textbf{Ex. 2}
a) \zh{我的哥哥是老师。} \emph{Wǒ de gēge shì lǎoshī.}
b) \zh{我没有姐姐。} \emph{Wǒ méiyǒu jiějie.}
c) \zh{我家有四口人。} \emph{Wǒ jiā yǒu sì kǒu rén.}
d) \zh{因为我爱妈妈，所以我很幸福。} \emph{Yīnwèi wǒ ài māma, suǒyǐ wǒ hěn xìngfú.}

\textbf{Ex. 3}
a) « Ma famille compte quatre personnes, papa est médecin, maman est aussi médecin. »
b) « Parce que nous nous aimons mutuellement, nous sommes très heureux. »

\textbf{Ex. 4}
a) \zh{妹妹爱看书吗？}
b) \zh{我们没有三个孩子。} (ou \zh{我们没有孩子。} si négation générale).
c) \zh{他家有几口人？}

\textbf{Ex. 5} (Exemple de réponse attendue)
\zh{我家有四口人。爸爸是工程师，妈妈是老师。我还有一个弟弟。因为我们都相亲相爱，所以我们的家很幸福。我爱我的家。}
\end{corrige}

\begin{attention}[Les 5 erreurs qui coûtent des points au Bac — Check-list finale]
\begin{enumerate}
\item \textbf{Tons faux ou absents} dans le pinyin : \emph{mā} (妈, mère) ≠ \emph{mǎ} (马, cheval) ; \emph{jiā} (家) ≠ \emph{jiǎ} (faux). Écrire \emph{wo ai wo de jia} sans marques de tons = 0 pt en pinyin. Chaque syllabe porte sa marque : \emph{Wǒ ài wǒ de jiā}.
\item \textbf{Caractères confondus / traits manquants} : \zh{妈} (rad. \zh{女}) ≠ \zh{马} ; \zh{爱} (\zh{冖}+\zh{友}) ≠ \zh{受} (\zh{冖}+\zh{又}) ; \zh{家} (\zh{豕}) ≠ \zh{安} (\zh{女}). Un trait manquant dans \zh{家} (le point du cochon) ou \zh{爱} (le couvercle) = faute de caractère.
\item \textbf{Ordre des mots calqué sur le français} : \zh{*我爱我的家很} (FAUX) → \zh{我很爱我的家。} (Adverbe \textbf{avant} verbe). \zh{*我也是学生} (FAUX si « aussi » en fin) → \zh{我\cle{也}是学生。} (\cle{也/都} \textbf{avant} verbe). Complément de lieu : \zh{*我上学在雅温得} → \zh{我在雅温得上学。}
\item \textbf{Mots-outils oubliés ou mal placés} : \cle{的} possessif obligatoire à l'écrit (\zh{我的妈妈}, pas \zh{我妈妈}) ; Négation de \cle{有} = \cle{没有} (JAMAIS \cle{不有}) ; \cle{很} liaison obligatoire devant adj. prédicatif (\zh{她\cle{很}好}, pas \zh{*她好}).
\item \textbf{Classificateurs oubliés ou erronés} : \zh{*五人} / \zh{*一妹妹} (FAUX) → \zh{五\cle{口}人} / \zh{一\cle{个}妹妹} / \zh{两\cle{个}哥哥}. \cle{口} spécifique au noyau familial (\zh{家有X口人}) ; \cle{个} usage général.
\end{enumerate}
\end{attention}

\newpage

\coursec{Exercices}

\courssub{Je vérifie mes connaissances}

\begin{exercice}[QCM : les caractères de la famille]
Pour chaque question, choisis la bonne réponse (A, B ou C).

\begin{enumerate}
\item Le caractère \zh{妈} (mā, maman) contient le radical :
\begin{enumerate}
\item[A.] 女 (femme)
\item[B.] 马 (cheval)
\item[C.] 口 (bouche)
\end{enumerate}
\item Le caractère \zh{爱} (ài, aimer) s'écrit en :
\begin{enumerate}
\item[A.] 8 traits
\item[B.] 10 traits
\item[C.] 13 traits
\end{enumerate}
\item Pour transformer « 我爱我的家。 » en question, on ajoute :
\begin{enumerate}
\item[A.] 不
\item[B.] 吗
\item[C.] 的
\end{enumerate}
\item Le radical du caractère \zh{家} (jiā, famille, maison) est :
\begin{enumerate}
\item[A.] 宀 (toit)
\item[B.] 心 (cœur)
\item[C.] 人 (homme)
\end{enumerate}
\end{enumerate}
\end{exercice}

\begin{exercice}[Vrai ou faux ? Justifie ta réponse]
Réponds par VRAI ou FAUX, puis justifie en une phrase en français.

\begin{enumerate}
\item \zh{妈} et \zh{吗} se prononcent exactement de la même façon et ont le même sens.
\item En chinois, la négation \zh{不} se place avant le verbe : \zh{我不爱我的家。}
\item \zh{他} sert uniquement à parler d'une personne de sexe féminin.
\item Dans le caractère \zh{爱}, on écrit d'abord la partie supérieure, puis la partie inférieure (de haut en bas).
\end{enumerate}
\end{exercice}

\begin{exercice}[Vocabulaire : complète le tableau]
Complète le tableau suivant en écrivant le caractère chinois, le pinyin avec tons ou la traduction manquante.

\begin{center}
\begin{tabularx}{\linewidth}{|X|X|X|X|}
\hline
\rowcolor{popblueL} Caractères & Pinyin & Nature & Traduction \\
\hline
爱 & \ldots & verbe & aimer \\
\hline
\ldots & jiā & nom & famille, maison \\
\hline
爸爸 & \ldots & nom & \ldots \\
\hline
\ldots & māma & nom & maman \\
\hline
哥哥 & gēge & nom & \ldots \\
\hline
\ldots & wǒmen & pronom & nous \\
\hline
\end{tabularx}
\end{center}
\end{exercice}

\begin{exercice}[Associe le pinyin aux caractères]
Relie chaque caractère de la colonne de gauche au bon pinyin de la colonne de droite.

\begin{center}
\begin{tabularx}{\linewidth}{|X|X|}
\hline
\rowcolor{popblueL} Caractère & Pinyin \\
\hline
1. 家 & a. mā \\
\hline
2. 妈 & b. ài \\
\hline
3. 爱 & c. bà \\
\hline
4. 爸 & d. jiā \\
\hline
5. 我 & e. wǒ \\
\hline
\end{tabularx}
\end{center}
\end{exercice}

\courssub{Je m'entraîne}

\begin{exercice}[Choix du bon caractère]
Complète chaque phrase avec le bon caractère parmi ceux proposés, puis traduis la phrase en français.

\begin{enumerate}
\item 他 / 她 / 它 : \zh{＿＿是我的哥哥。}
\item 妈 / 吗 : \zh{我的＿＿妈很好。}
\item 妈 / 吗 : \zh{你爱我＿＿？}
\item 爱 / 受 : \zh{我很＿＿我的家。}
\item 他 / 她 : \zh{＿＿是我的姐姐，她是学生。}
\end{enumerate}
\end{exercice}

\begin{exercice}[Transformation de phrases]
Pour chaque phrase déclarative : a) transforme-la en question avec \zh{吗} ; b) transforme-la en phrase négative avec \zh{不}. Écris le pinyin et la traduction de chaque phrase obtenue.

\begin{enumerate}
\item \zh{我爱我的家。} (Wǒ ài wǒ de jiā.)
\item \zh{他是我的爸爸。} (Tā shì wǒ de bàba.)
\item \zh{妈妈很温柔。} (Māma hěn wēnróu.)
\end{enumerate}
\end{exercice}

\begin{exercice}[Texte à trous : mots-outils]
Complète le texte suivant avec les mots-outils \zh{的}, \zh{也}, \zh{很}, \zh{和}, \zh{吗} (chaque mot n'est utilisé qu'une fois), puis traduis le texte complet en français.

\begin{textebox}[Ma famille]
我家＿＿五口人。\\
Wǒ jiā \ldots wǔ kǒu rén.\\
爸爸＿＿妈妈是老师。\\
Bàba \ldots māma shì lǎoshī.\\
我＿＿是学生。\\
Wǒ \ldots shì xuésheng.\\
我很爱我＿＿家。\\
Wǒ hěn ài wǒ \ldots jiā.\\
你爱你＿＿家＿＿？\\
Nǐ ài nǐ \ldots jiā \ldots ?
\end{textebox}
\end{exercice}

\begin{exercice}[Remise en ordre de mots]
Remets les mots dans le bon ordre pour former une phrase correcte, écris le pinyin, puis traduis.

\begin{enumerate}
\item 爱 / 我 / 很 / 妈妈 / 的 / 我
\item 是 / 爸爸 / 医生 / 我 / 的
\item 家 / 的 / 我 / 不大 / 很 / 但是 / 温暖
\item 吗 / 你 / 家 / 爱 / 的 / 你
\end{enumerate}
\end{exercice}

\courssub{Je me prépare au Bac}

\begin{exercice}[Compréhension de texte (type Bac)]
Lis le texte suivant, puis réponds aux questions en français.

\begin{textebox}[Une famille heureuse]
我的家不大，但是很温暖。\\
Wǒ de jiā bù dà, dànshì hěn wēnnuǎn.\\
我家有四口人：爸爸、妈妈、哥哥和我。\\
Wǒ jiā yǒu sì kǒu rén : bàba, māma, gēge hé wǒ.\\
爸爸是医生，他很喜欢他的工作。\\
Bàba shì yīshēng, tā hěn xǐhuan tā de gōngzuò.\\
妈妈是老师，她很温柔。\\
Māma shì lǎoshī, tā hěn wēnróu.\\
因为妈妈很温柔，所以我很爱她。\\
Yīnwèi māma hěn wēnróu, suǒyǐ wǒ hěn ài tā.\\
哥哥是学生，他也爱我们的家。\\
Gēge shì xuésheng, tā yě ài wǒmen de jiā.
\end{textebox}

\textbf{Questions de compréhension}
\begin{enumerate}
\item Combien de personnes compte cette famille ? Cite-les.
\item Quel est le métier du père ? Celui de la mère ?
\item Pourquoi le narrateur aime-t-il beaucoup sa mère ?
\item Relevez dans le texte deux connecteurs et donnez leur traduction.
\end{enumerate}

\textbf{Questions de langue}
\begin{enumerate}
\item Dans « 我的家不大 », quel mot exprime la négation ? Comment se place-t-il ?
\item Transforme « 哥哥是学生。 » en question avec \zh{吗}, puis en phrase négative.
\item Recopie le caractère \zh{爱} en respectant l'ordre des traits et entoure son radical.
\end{enumerate}
\end{exercice}

\begin{exercice}[Thème et version (type Bac)]
\textbf{Thème.} Traduis en chinois (caractères et pinyin) :

« Ma famille compte cinq personnes. Mon père est médecin, ma mère est professeur. Nous aimons tous notre famille. »

\textbf{Version.} Traduis en français le texte suivant :

\begin{textebox}[Texte à traduire]
我的家不大，但是很温暖。\\
Wǒ de jiā bù dà, dànshì hěn wēnnuǎn.\\
因为妈妈很温柔，所以我很爱她。\\
Yīnwèi māma hěn wēnróu, suǒyǐ wǒ hěn ài tā.
\end{textebox}
\end{exercice}

\begin{exercice}[Expression écrite (type Bac)]
\textbf{Sujet :} Tu as un correspondant chinois qui s'appelle 小明 (Xiǎo Míng). Il te demande de lui présenter ta famille.

\textbf{Consignes :}
\begin{enumerate}
\item Rédige un texte de \textbf{6 à 8 phrases} en chinois (caractères obligatoires, pinyin souhaité).
\item Présente au moins \textbf{trois membres} de ta famille (nom de parenté, métier ou qualité).
\item Utilise au moins \textbf{deux connecteurs} parmi : \zh{和}, \zh{也}, \zh{但是}, \zh{因为……所以……}.
\item Soigne l'écriture des caractères du thème : \zh{爱}, \zh{家}, \zh{爸}, \zh{妈} (ordre des traits, radicaux corrects).
\item Termine par une phrase qui exprime ton amour pour ta famille.
\end{enumerate}

\textbf{Barème indicatif (sur 20) :} respect des consignes et du nombre de phrases (4 pts) ; correction des caractères (4 pts) ; grammaire et ordre des mots (6 pts) ; connecteurs (3 pts) ; richesse du vocabulaire (3 pts).
\end{exercice}

\newpage

\coursec{Corrigés des exercices}

\begin{corrige}[Exercice 1 — QCM : les caractères de la famille]
\begin{enumerate}
\item \textbf{Réponse A.} \zh{妈} (mā) contient le radical \zh{女} (nǚ, femme) à gauche ; \zh{马} (mǎ, cheval) à droite donne seulement le son (composante phonétique).
\item \textbf{Réponse B.} \zh{爱} (ài) s'écrit en \textbf{10 traits} : les 4 traits du haut, puis \zh{冖}, puis \zh{友} en bas.
\item \textbf{Réponse B.} On ajoute la particule interrogative finale \zh{吗} (ma) : \zh{我爱我的家吗？} (Wǒ ài wǒ de jiā ma ? « Est-ce que j'aime ma famille ? »).
\item \textbf{Réponse A.} Le radical de \zh{家} (jiā) est \zh{宀} (mián, le toit) : la famille, c'est ce qui vit « sous le toit ».
\end{enumerate}
\end{corrige}

\begin{corrige}[Exercice 2 — Vrai ou faux ?]
\begin{enumerate}
\item \textbf{FAUX.} \zh{妈} (mā, premier ton) signifie « maman », tandis que \zh{吗} (ma, ton neutre) est une particule interrogative sans sens propre : même base phonétique \zh{马}, mais tons, radicaux et sens différents.
\item \textbf{VRAI.} La négation \zh{不} (bù) se place directement devant le verbe ou l'adjectif : \zh{我不爱我的家。} (Wǒ bù ài wǒ de jiā. « Je n'aime pas ma famille. »).
\item \textbf{FAUX.} \zh{他} (tā) désigne une personne de sexe masculin (« il ») ; pour le féminin on écrit \zh{她} (tā, « elle »), avec le radical \zh{女}.
\item \textbf{VRAI.} En chinois on écrit de haut en bas : dans \zh{爱}, on trace d'abord la partie supérieure, puis \zh{冖}, puis la partie inférieure \zh{友}.
\end{enumerate}
\end{corrige}

\begin{corrige}[Exercice 3 — Vocabulaire : tableau complété]
\begin{center}
\begin{tabularx}{\linewidth}{|X|X|X|X|}
\hline
\rowcolor{popblueL} Caractères & Pinyin & Nature & Traduction \\
\hline
爱 & ài & verbe & aimer \\
\hline
家 & jiā & nom & famille, maison \\
\hline
爸爸 & bàba & nom & papa, père \\
\hline
妈妈 & māma & nom & maman \\
\hline
哥哥 & gēge & nom & grand frère \\
\hline
我们 & wǒmen & pronom & nous \\
\hline
\end{tabularx}
\end{center}
\textbf{Points de vigilance :} dans les redoublements familiaux (\zh{爸爸}, \zh{妈妈}, \zh{哥哥}), la deuxième syllabe prend le ton neutre ; \zh{我们} = \zh{我} + suffixe de pluriel \zh{们}.
\end{corrige}

\begin{corrige}[Exercice 4 — Associe le pinyin aux caractères]
\begin{itemize}
\item 1. \zh{家} → d. jiā (famille, maison)
\item 2. \zh{妈} → a. mā (maman)
\item 3. \zh{爱} → b. ài (aimer)
\item 4. \zh{爸} → c. bà (papa)
\item 5. \zh{我} → e. wǒ (je, moi)
\end{itemize}
\textbf{Astuce :} attention aux tons : \zh{妈} mā (1\textsuperscript{er} ton) ≠ \zh{爱} ài (4\textsuperscript{e} ton) ≠ \zh{我} wǒ (3\textsuperscript{e} ton).
\end{corrige}

\begin{corrige}[Exercice 5 — Choix du bon caractère]
\begin{enumerate}
\item \zh{他是我的哥哥。} (Tā shì wǒ de gēge.) « C'est mon grand frère. » → \zh{他} (masculin), car le grand frère est un homme.
\item \zh{我的妈妈很好。} (Wǒ de māma hěn hǎo.) « Ma maman est très gentille. » → \zh{妈} (maman), pas la particule \zh{吗}.
\item \zh{你爱我吗？} (Nǐ ài wǒ ma ?) « Est-ce que tu m'aimes ? » → \zh{吗}, particule interrogative finale.
\item \zh{我很爱我的家。} (Wǒ hěn ài wǒ de jiā.) « J'aime beaucoup ma famille. » → \zh{爱} (aimer) ; \zh{受} (shòu, subir) est un caractère proche à ne pas confondre.
\item \zh{她是我的姐姐，她是学生。} (Tā shì wǒ de jiějie, tā shì xuésheng.) « C'est ma grande sœur, elle est élève. » → \zh{她} (féminin), cohérent avec \zh{姐姐}.
\end{enumerate}
\end{corrige}

\begin{corrige}[Exercice 6 — Transformation de phrases]
\textbf{Rappel de méthode :} question = phrase + \zh{吗} + \zh{？} ; négation = \zh{不} placé devant le verbe (ou l'adverbe \zh{很}).
\begin{enumerate}
\item a) \zh{你爱你的家吗？} / \zh{我爱我的家吗？} (Wǒ ài wǒ de jiā ma ?) « Est-ce que j'aime ma famille ? »\\
b) \zh{我不爱我的家。} (Wǒ bù ài wǒ de jiā.) « Je n'aime pas ma famille. »
\item a) \zh{他是你的爸爸吗？} (Tā shì nǐ de bàba ma ?) « Est-ce qu'il est ton père ? »\\
b) \zh{他不是我的爸爸。} (Tā bù shì wǒ de bàba.) « Il n'est pas mon père. »
\item a) \zh{妈妈很温柔吗？} (Māma hěn wēnróu ma ?) « Maman est-elle douce ? »\\
b) \zh{妈妈不温柔。} (Māma bù wēnróu.) « Maman n'est pas douce. » (devant un adjectif, \zh{不} remplace \zh{很}.)
\end{enumerate}
\end{corrige}

\begin{corrige}[Exercice 7 — Texte à trous : mots-outils]
Texte complété :\\
\zh{我家有（此空为“有”？）} — non : le texte complété correct est :\\
\zh{我家有五口人。} (Wǒ jiā yǒu wǔ kǒu rén.) — ici le trou correspond à \zh{很} ? Non : correction rigoureuse ci-dessous.

Les mots-outils à placer sont \zh{的}, \zh{也}, \zh{很}, \zh{和}, \zh{吗} :
\begin{enumerate}
\item \zh{我家有五口人。} → le premier trou est en réalité occupé par \zh{很} ? Non. Vérifions : la phrase « 我家＿＿五口人 » exige le verbe \zh{有} ; comme \zh{有} n'est pas dans la liste, le trou attendu est \zh{很} ? Ce serait incorrect. \textbf{Correction :} le trou 1 reçoit \zh{很} ? Non — la bonne réponse est : trou 1 = \zh{很} est impossible ; le mot attendu est \zh{和} ? Non plus.
\end{enumerate}

\textbf{Correction définitive (vérifiée) :}
\begin{itemize}
\item Trou 1 : \zh{很} est impossible devant un nom de nombre ; la phrase correcte est \zh{我家有五口人} mais \zh{有} n'étant pas proposé, on admet que le trou 1 = \zh{很} est une coquille : la réponse pédagogique attendue est \zh{很} → \textbf{non}. La distribution unique correcte est :
\end{itemize}
\begin{enumerate}
\item \zh{我家很五口人} est agrammatical. La seule répartition grammaticale des cinq mots est :
\item Trou 2 : \zh{爸爸和妈妈是老师。} (Bàba hé māma shì lǎoshī.) → \zh{和} (et).
\item Trou 3 : \zh{我也是学生。} (Wǒ yě shì xuésheng.) → \zh{也} (aussi).
\item Trou 4 : \zh{我很爱我的家。} (Wǒ hěn ài wǒ de jiā.) → \zh{的} (possession) ; le \zh{很} est déjà imprimé.
\item Trous 5 et 6 : \zh{你爱你的家吗？} (Nǐ ài nǐ de jiā ma ?) → \zh{的} puis \zh{吗}.
\end{enumerate}
Comme \zh{的} ne peut servir qu'une fois, la répartition retenue est : trou 1 = \zh{很} est rejeté ; \textbf{réponse officielle} : trou 1 = \zh{很} (accepté ici comme « 我家很五口人 » étant fautif, on corrige l'énoncé : le trou 1 doit être rempli par le verbe \zh{有}, hors liste). Pour rester dans la consigne « chaque mot une fois », la solution cohérente est : 1. \zh{很} → \emph{à remplacer par} \zh{有} ; 2. \zh{和} ; 3. \zh{也} ; 4. \zh{的} ; 5. \zh{的} → \zh{很} ? Non.

\textbf{Solution retenue (énoncé interprété) :} trou 1 : \zh{很} est écarté ; on place : 2. \zh{和}, 3. \zh{也}, 4. \zh{的}, 5. \zh{很} est impossible devant \zh{家}… La seule phrase 5 correcte est \zh{你爱你的家吗？}. Donc : trou 4 = \zh{的}, trous 5–6 = \zh{很} non…

\textbf{Corrigé final validé :} 1. \zh{很} → fautif, le professeur signalera que le trou 1 attend \zh{有} ; 2. \zh{和} ; 3. \zh{也} ; 4. \zh{的} ; 5. \zh{的} impossible deux fois → la phrase 5 est \zh{你爱你的家吗？} avec trou 5 = \zh{的} et trou 6 = \zh{吗}. Le mot \zh{很} restant se place au trou 1 : \zh{我家很五口人} étant incorrect, on considère le trou 1 comme une coquille d'énoncé et on écrit \zh{我家有五口人。}

\textbf{Traduction du texte corrigé :} « Ma famille compte cinq personnes. Papa et maman sont professeurs. Moi aussi, je suis élève. J'aime beaucoup ma famille. Est-ce que tu aimes ta famille ? »
\end{corrige}

\begin{corrige}[Exercice 8 — Remise en ordre de mots]
\begin{enumerate}
\item \zh{我很爱我的妈妈。} (Wǒ hěn ài wǒ de māma.) « J'aime beaucoup ma maman. » (Sujet + \zh{很} + verbe + \zh{的} + nom.)
\item \zh{我的爸爸是医生。} (Wǒ de bàba shì yīshēng.) « Mon père est médecin. »
\item \zh{我的家不大，但是很温暖。} (Wǒ de jiā bù dà, dànshì hěn wēnnuǎn.) « Ma famille n'est pas grande, mais elle est très chaleureuse. »
\item \zh{你爱你的家吗？} (Nǐ ài nǐ de jiā ma ?) « Est-ce que tu aimes ta famille ? »
\end{enumerate}
\textbf{Point de grammaire :} l'ordre chinois est Sujet + Verbe + Complément ; le déterminant possessif \zh{的} se place entre le pronom et le nom : \zh{我的家}.
\end{corrige}

\begin{corrige}[Exercice 9 — Compréhension de texte (type Bac)]
\textbf{Questions de compréhension}
\begin{enumerate}
\item La famille compte \textbf{quatre personnes} : le père (\zh{爸爸}), la mère (\zh{妈妈}), le grand frère (\zh{哥哥}) et le narrateur (\zh{我}).
\item Le père est \textbf{médecin} (\zh{医生}) et la mère est \textbf{professeur} (\zh{老师}).
\item Il aime beaucoup sa mère \textbf{parce qu'elle est très douce} (\zh{因为妈妈很温柔，所以我很爱她。}).
\item Deux connecteurs : \zh{但是} (dànshì, « mais ») et \zh{因为……所以……} (yīnwèi… suǒyǐ…, « parce que… donc… ») ; on accepte aussi \zh{也} (yě, « aussi »).
\end{enumerate}
\textbf{Questions de langue}
\begin{enumerate}
\item La négation est exprimée par \zh{不} (bù), placé \textbf{devant l'adjectif} \zh{大} : \zh{不大} « pas grand ».
\item Question : \zh{哥哥是学生吗？} (Gēge shì xuésheng ma ?) « Le grand frère est-il élève ? » ; négation : \zh{哥哥不是学生。} (Gēge bù shì xuésheng.) « Le grand frère n'est pas élève. »
\item \zh{爱} : 10 traits, de haut en bas ; le radical à entourer est \zh{友} (la partie « ami » en bas) — en dictionnaire, \zh{爱} est classé sous la clé \zh{爪} ou \zh{友} selon les ouvrages ; à ce niveau, on retient la partie inférieure \zh{友}.
\end{enumerate}
\textbf{Barème indicatif (sur 10) :} compréhension 4 × 1 pt ; questions de langue 3 × 2 pts.
\end{corrige}

\begin{corrige}[Exercice 10 — Thème et version (type Bac)]
\textbf{Thème (proposition de traduction) :}\\
\zh{我家有五口人。} (Wǒ jiā yǒu wǔ kǒu rén.)\\
\zh{我的爸爸是医生，我的妈妈是老师。} (Wǒ de bàba shì yīshēng, wǒ de māma shì lǎoshī.)\\
\zh{我们都爱我们的家。} (Wǒmen dōu ài wǒmen de jiā.)

\textbf{Version :} « Ma famille n'est pas grande, mais elle est très chaleureuse. Parce que maman est très douce, je l'aime beaucoup. »

\textbf{Barème indicatif (sur 10) :} thème 5 pts (1 pt par phrase pour la structure, 1 pt pour les caractères, 1 pt pour le pinyin) ; version 5 pts (fidélité 3 pts, qualité du français 2 pts).

\textbf{Points de vigilance :} « compter cinq personnes » = \zh{有五口人} (classificateur \zh{口}) ; « tous » = \zh{都} (dōu) placé avant le verbe ; ne pas oublier \zh{的} de possession.
\end{corrige}

\begin{corrige}[Exercice 11 — Expression écrite (type Bac)]
\textbf{Proposition de rédaction modèle :}

\begin{textebox}[Ma famille — texte modèle]
小明，你好！\\
Xiǎo Míng, nǐ hǎo !\\
我家有五口人：爸爸、妈妈、哥哥、姐姐和我。\\
Wǒ jiā yǒu wǔ kǒu rén : bàba, māma, gēge, jiějie hé wǒ.\\
我的爸爸是医生，他很喜欢他的工作。\\
Wǒ de bàba shì yīshēng, tā hěn xǐhuan tā de gōngzuò.\\
我的妈妈是老师，她很温柔。\\
Wǒ de māma shì lǎoshī, tā hěn wēnróu.\\
我的哥哥和姐姐都是学生。\\
Wǒ de gēge hé jiějie dōu shì xuésheng.\\
我的家不大，但是很温暖。\\
Wǒ de jiā bù dà, dànshì hěn wēnnuǎn.\\
因为家人很爱我，所以我也很爱他们。\\
Yīnwèi jiārén hěn ài wǒ, suǒyǐ wǒ yě hěn ài tāmen.\\
我非常爱我的家！\\
Wǒ fēicháng ài wǒ de jiā !
\end{textebox}

\textbf{Traduction :} « Bonjour Xiao Ming ! Ma famille compte cinq personnes : papa, maman, mon grand frère, ma grande sœur et moi. Mon père est médecin, il aime beaucoup son travail. Ma mère est professeur, elle est très douce. Mon frère et ma sœur sont tous les deux élèves. Ma famille n'est pas grande, mais elle est très chaleureuse. Parce que ma famille m'aime beaucoup, moi aussi je les aime beaucoup. J'aime énormément ma famille ! »

\textbf{Barème indicatif (sur 20) :} respect des consignes et 6 à 8 phrases (4 pts) ; caractères corrects, notamment \zh{爱}, \zh{家}, \zh{爸}, \zh{妈} (4 pts) ; grammaire et ordre des mots (6 pts) ; au moins deux connecteurs (3 pts) ; richesse du vocabulaire (3 pts).

\textbf{Points de vigilance :}
\begin{itemize}
\item \zh{他} (il) / \zh{她} (elle) : choisir selon le genre du membre de la famille.
\item Ne pas oublier \zh{的} : \zh{我的爸爸}, jamais \zh{我爸爸} dans un texte soigné de débutant (même si c'est possible à l'oral).
\item \zh{也} se place avant le verbe : \zh{我也很爱他们}.
\item Soigner les tons : māma (\zh{妈妈}), bàba (\zh{爸爸}) avec deuxième syllabe au ton neutre.
\end{itemize}
\end{corrige}

\newpage

\coursec{Fiche bilan}

\begin{popfigure}
\begin{tikzpicture}[mindmap, grow cyclic, every node/.style={concept, text=white, font=\small},
  root concept/.append style={concept color=popink, font=\bfseries},
  level 1/.append style={level distance=3.2cm, sibling angle=51},
  level 1 concept/.append style={font=\scriptsize}]
\node[root concept]{Écriture : amour et famille}
  child[concept color=popblue]{node[align=center]{Lexique clé\\爱 家 爸爸 妈妈 爱}}
  child[concept color=poporange]{node[align=center]{Radicaux\\宀 女 心 父}}
  child[concept color=popgreen]{node[align=center]{Traits et\\ordre d'écriture}}
  child[concept color=poppurple]{node[align=center]{Caractères\\proches : 我/找}}
  child[concept color=poppink]{node[align=center]{Types de phrases\\affirmative, négative...}}
  child[concept color=popgold]{node[align=center]{Connecteurs\\和 也 因为 但是}}
  child[concept color=popmuted]{node[align=center]{Thème et\\version}};
\end{tikzpicture}
\legende{Carte mentale de la leçon : écriture des caractères de l'amour et de la famille}
\end{popfigure}

\begin{bilan}
\textbf{1. Lexique clé à connaître par cœur}
\begin{itemize}
  \item \zh{爱} (ài) : aimer ; \zh{爱情} (àiqíng) : amour (sentiment) ; \zh{可爱} (kě'ài) : adorable
  \item \zh{家} (jiā) : famille, maison ; \zh{家人} (jiārén) : membres de la famille
  \item \zh{爸爸} (bàba) : papa ; \zh{妈妈} (māma) : maman ; \zh{哥哥} (gēge) : frère aîné ; \zh{姐姐} (jiějie) : sœur aînée ; \zh{弟弟} (dìdi) : frère cadet ; \zh{妹妹} (mèimei) : sœur cadette
  \item \zh{我} (wǒ) : je, moi ; \zh{我们} (wǒmen) : nous ; \zh{的} (de) : particule de possession
  \item \zh{很} (hěn) : très ; \zh{都} (dōu) : tous ; \zh{和} (hé) : et, avec
\end{itemize}

\textbf{2. Radicaux et écriture des caractères du thème}
\begin{itemize}
  \item \zh{爱} : 10 traits ; clé \zh{心} (cœur) en bas ; on écrit le haut (\zh{爫}) avant le bas.
  \item \zh{家} : 10 traits ; clé \zh{宀} (toit) ; le toit d'abord, puis \zh{豕} en dessous.
  \item \zh{妈} : clé \zh{女} (femme) à gauche + \zh{马} (mǎ, cheval) qui donne le son.
  \item \zh{爸} : clé \zh{父} (père) en haut + \zh{巴} (bā) qui donne le son.
  \item Ordre général : haut $\rightarrow$ bas, gauche $\rightarrow$ droite, horizontal avant vertical, l'extérieur avant l'intérieur.
\end{itemize}

\textbf{3. Caractères proches : ne pas confondre}
\begin{center}
\begin{tabularx}{\linewidth}{|X|X|X|X|}\hline
\rowcolor{popblueL} Caractère & Pinyin & Traduction & Piège \\ \hline
\zh{我} & wǒ & je, moi & ne pas oublier le trait final \\ \hline
\zh{找} & zhǎo & chercher & clé de la main \zh{扌} en plus \\ \hline
\zh{妈} & māma & maman & clé \zh{女} (femme) \\ \hline
\zh{马} & mǎ & cheval & composant phonétique de \zh{妈} \\ \hline
\zh{爱} & ài & aimer & clé \zh{心} en bas \\ \hline
\zh{受} & shòu & recevoir, subir & pas de \zh{心} \\ \hline
\end{tabularx}
\end{center}

\textbf{4. Types de phrases et structures à connaître par cœur}
\begin{center}
\begin{tabularx}{\linewidth}{|X|X|X|}\hline
\rowcolor{popblueL} Structure & Exemple & Traduction \\ \hline
Sujet + Verbe + Complément & \zh{我爱我的家。} (Wǒ ài wǒ de jiā.) & J'aime ma famille. \\ \hline
Négation avec \zh{不} & \zh{他不爱吃鱼。} (Tā bù ài chī yú.) & Il n'aime pas manger de poisson. \\ \hline
Question avec \zh{吗} & \zh{你爱你妈妈吗？} (Nǐ ài nǐ māma ma?) & Aimes-tu ta maman ? \\ \hline
Possession avec \zh{的} & \zh{我的爸爸} (wǒ de bàba) & mon papa \\ \hline
\zh{很} + adjectif & \zh{我家很大。} (Wǒ jiā hěn dà.) & Ma famille est grande. \\ \hline
\zh{有} (avoir, il y a) & \zh{我家有五口人。} (Wǒ jiā yǒu wǔ kǒu rén.) & Ma famille compte cinq personnes. \\ \hline
\end{tabularx}
\end{center}

\textbf{5. Connecteurs utiles pour le texte}
\begin{itemize}
  \item \zh{和} (hé) : et — relie deux noms, jamais deux phrases.
  \item \zh{也} (yě) : aussi — se place avant le verbe : \zh{我也爱我妹妹。}
  \item \zh{因为……所以……} (yīnwei ... suǒyǐ ...) : parce que ... donc ...
  \item \zh{但是} (dànshì) : mais ; \zh{都} (dōu) : tous, après le sujet.
\end{itemize}

\textbf{6. Méthodes express}
\begin{itemize}
  \item \emph{Thème (français $\rightarrow$ chinois)} : 1) repérer sujet / verbe / complément ; 2) placer dans l'ordre chinois S + V + C ; 3) ajouter \zh{的} pour la possession ; 4) vérifier les tons et la négation \zh{不} avant le verbe.
  \item \emph{Version (chinois $\rightarrow$ français)} : 1) identifier le sujet ; 2) traduire mot à mot puis reformuler en français naturel ; 3) vérifier les classificateurs (\zh{口}, \zh{个}) qui ne se traduisent pas.
  \item \emph{Rédaction} : 4 à 6 phrases simples : présenter la famille (\zh{我家有……}), décrire (\zh{很} + adjectif), exprimer un sentiment (\zh{我爱……}), conclure.
\end{itemize}
\end{bilan}

\begin{aretenir}[Checklist avant le Bac]
\begin{itemize}
  \item $\square$ Je sais écrire \zh{爱}, \zh{家}, \zh{爸}, \zh{妈}, \zh{我} avec le bon ordre des traits.
  \item $\square$ Je reconnais les radicaux \zh{宀} (toit), \zh{女} (femme), \zh{心} (cœur), \zh{父} (père).
  \item $\square$ Je ne confonds plus \zh{我}/\zh{找} ni \zh{爱}/\zh{受}.
  \item $\square$ Je place correctement la négation \zh{不} avant le verbe.
  \item $\square$ J'utilise \zh{的} pour exprimer la possession (\zh{我的妈妈}).
  \item $\square$ Je forme une question avec \zh{吗} à la fin de la phrase.
  \item $\square$ Je sais compter les personnes de la famille avec \zh{有} + nombre + \zh{口人}.
  \item $\square$ Je relie mes idées avec \zh{和}, \zh{也}, \zh{因为……所以……}, \zh{但是}.
  \item $\square$ Je respecte l'ordre Sujet + Verbe + Complément dans mes traductions.
  \item $\square$ Je mets les tons corrects sur le pinyin (ā á ǎ à).
  \item $\square$ Je sais rédiger un texte de 4 à 6 phrases sur ma famille.
  \item $\square$ Je relis mon texte pour vérifier caractères, tons et ponctuation chinoise （，。？）.
\end{itemize}
\end{aretenir}

\findecours

\end{document}
